% concatenated from main.tex
\documentclass[a4paper,11pt]{article}
\usepackage[utf8]{inputenc}

\usepackage{boxedminipage}
\usepackage{amsfonts}
\usepackage{amsmath} 
\usepackage{amssymb}
\usepackage{graphicx}
\usepackage{amsthm}
\usepackage{t1enc}
\usepackage{subfig}
\usepackage{authblk} % das package schaltet einen Affiliation operator frei
 \usepackage{mathabx}
 \usepackage{url}




\newtheorem{theorem}{Theorem}[section]
\newtheorem{lemma}[theorem]{Lemma}
\newtheorem{proposition}[theorem]{Proposition}
\newtheorem{corollary}[theorem]{Corollary}
\newtheorem{definition}[theorem]{Definition}
\newtheorem{example}[theorem]{Example}
\newtheorem{fact}[theorem]{Fact}
\newtheorem{remark}[theorem]{Remark}
\newtheorem{question}{Question}
\newtheorem{claim}{Claim}

\newcommand{\forceP}{\mathbb{P}}
\newcommand{\forceQ}{\mathbb{Q}}
\newcommand{\forceL}{\mathbb{L}}
\newcommand{\hybcomp}{\mathbin{\ltimes_{\mathrm{hyb}}}}

\newcommand{\last}{\mathsf{last}}


\newcommand{\NS}{\hbox{NS}_{\omega_1}}
\newcommand{\ZFC}{\mathsf{ZFC}}
\newcommand{\ZF}{\mathsf{ZF}}
\newcommand{\ZFP}{\mathsf{ZF}^-}
\newcommand{\AC}{\mathsf{AC}}
\newcommand{\PFA}{\mathsf{PFA}}
\newcommand{\BPFA}{\mathsf{BPFA}}
\newcommand{\MRP}{\mathsf{MRP}}
\newcommand{\MM}{\mathsf{MM}}
\newcommand{\MMpp}{\mathsf{MM^{++}}}
\newcommand{\CH}{\mathsf{CH}}
\newcommand{\MApp}{\mathsf{MA_{\omega_1} (\sigma-\text{closed})^{++}}}
\newcommand{\BMMpp}{\mathsf{BMM^{++}}}
\newcommand{\BMM}{\mathsf{BMM}}
\newcommand{\MA}{\mathsf{MA_{\omega_1}}}
\newcommand{\AD}{\mathsf{AD}}
\newcommand{\PD}{\mathsf{PD}}
\newcommand{\T}{\mathsf{T}}
\newcommand{\mouseM}{\mathcal{M}}
\newcommand{\mouseN}{\mathcal{N}}


\newcommand{\Reg}{\textrm{Reg}}
\newcommand{\Lim}{\mathop{\textrm{Lim}}}
\newcommand{\Suc}{\mathop{\textrm{Suc}}}
\newcommand{\dom}{\mathop{\textrm{dom}}}
\newcommand{\ran}{\mathop{\textrm{ran}}}
\newcommand{\spt}{\mathop{\textrm{spt}}}
\DeclareMathOperator{\ot}{ot}
\DeclareMathOperator{\cf}{cf}
\DeclareMathOperator{\supp}{supp}
\newcommand{\length}{\mathop{\textrm{length}}}
\newcommand\Th{\mathop{\textrm{Th}}}
\newcommand\Field{\mathop{\textrm{Field}}}



%\setlength{\parindent}{0pt}
%\setlength{\parskip}{2ex plus1ex minus 0.5ex}
%\renewcommand{\baselinestretch}{1.0}
%\setlength{\topmargin}{-0.62in}
%\setlength{\textheight}{9.10in}
%\setlength{\oddsidemargin}{-0.34in}
%\setlength{\evensidemargin}{-0.34in}
%\setlength{\textwidth}{6.95in}

\newcommand{\Col}{\textrm{Col}}
\newcommand{\Add}{\textrm{Add}}
\newcommand{\Ord}{\textrm{Ord}}
\newcommand{\Cof}{\textrm{Cof}}
\newcommand{\mba}{\mathbb{A}}
\newcommand{\mbb}{\mathbb{B}}
\newcommand{\mbc}{\mathbb{C}}
\newcommand{\mbbi}{\mathbb{B}_i}
\newcommand{\mbbj}{\mathbb{B}_j}
\newcommand{\mrg}{\mathring{G}}
\newcommand{\mrs}{\mathring{\sigma}}
\newcommand{\mbp}{\mathbb{P}}
\newcommand{\mbq}{\mathbb{Q}}
%\renewcommand{\qedsymbol}{QED}
\newcommand{\dbar}[1]{\bar{\bar{#1}}}
\newcommand{\mfm}{\mathfrak{M}}
\newcommand{\mcl}{\mathcal{L}}
\newcommand{\mcr}{\mathcal{R}}


\newcommand{\lra}{\longrightarrow}
\newcommand{\Lsharp}{L^{\#}}
\newcommand{\GCH}{\mathsf{GCH}}


\def\undertilde#1{\mathord{\vtop{\ialign{##\crcr
$\hfil\displaystyle{#1}\hfil$\crcr\noalign{\kern1.5pt\nointerlineskip}
$\hfil\tilde{}\hfil$\crcr\noalign{\kern1.5pt}}}}}

%opening
\title{On a local variant of the 12th Delfino problem --- the $\Pi$-side} 
\author{ Stefan Hoffelner$^{1}$  \footnote{The author's research was funded in whole by the Austrian Science Fund (FWF) Grant-DOI 10.55776/P37228. For the purpose of open access, the author has applied a CC BY public copyright license to any Author Accepted Manuscript version arising from this submission. }  }
\date{
    $^1$TU Wien \\
    \today
}

\begin{document}
\maketitle


\begin{abstract}
Assume that \(M_n\), the canonical inner model with \(n\) Woodin cardinals,
exists.  We force a model with continuum \(\aleph_2\) in which every
\(\boldsymbol{\Sigma}^1_{n+2}\) set of reals is Lebesgue measurable and has the
Baire property, the \(\Sigma^1_{n+2}\)- and
\(\Pi^1_{n+3}\)-uniformization properties hold, and the reals admit a
\(\Delta^1_{n+3}\)-definable well-order.  Thus regularity up to a fixed finite
projective level, together with a definable well-order of the reals at the
adjacent level, does not force the determinacy strength which would normally
explain that regularity, even when this package is strengthened by adjacent
\(\Sigma\)- and \(\Pi\)-uniformization.  In particular, this gives a negative
answer to a local form of Woodin's twelfth Delfino problem asked by Friedman--Schindler.
\end{abstract}


\section{Introduction}\label{sec:introduction}


A central theme in descriptive set theory is the interaction between
regularity properties, such as Lebesgue measurability and the Baire
property, and definable choice principles for projective sets of reals.
This paper concerns finite projective configurations in which regularity,
uniformization, and definable well-ordering occur together.  The canonical
source of such configurations is the hierarchy of mice with finitely many
Woodin cardinals.  Let \(M_k\) denote the canonical inner model with \(k\)
Woodin cardinals.  Steel's analysis shows that \(M_k\) has a
\(\Delta^1_{k+2}\)-definable well-order of its reals
\cite{Steel2,Steel3}.  On the other hand, the Woodin cardinals in these
models account, through the corresponding local determinacy and scale
analysis, for regularity and uniformization phenomena at lower projective
levels.

Thus \(M_{n+1}\) suggests a specific finite constellation in the
projective hierarchy.  Below the well-ordering level one has the
regularity and scale-theoretic consequences associated with the local
determinacy fragment; at the next level the fine structure of \(M_{n+1}\)
supplies a \(\Delta^1_{n+3}\)-definable well-order of the reals.  The
visible pattern relevant here is therefore regularity for boldface
\(\boldsymbol\Sigma^1_{n+2}\) sets, the corresponding finite
uniformization behavior, and a \(\Delta^1_{n+3}\)-definable well-order.

Before the present work, the systematic examples exhibiting this exact
constellation came from \(M_{n+1}\) and closely related canonical models;
in particular, they satisfied the corresponding
\(\boldsymbol\Delta^1_{n+1}\)-determinacy fragment.  Since the standard
source for this behavior of the projective sets of reals is the
\(M_{n+1}\)-picture, it is natural to ask whether this source is unique.
Does the finite projective behavior just described imply the corresponding
local determinacy hypothesis, or can the same behavior be obtained over a
weaker mouse?



The global antecedent of this recognition question is Woodin's twelfth
Delfino problem.  In one formulation, the problem asks whether $\mathsf{ZFC}$
together with the statements that every projective set is Lebesgue measurable
and has the Baire property, and that every projective relation admits a
projective uniformization, implies $\PD$.  Woodin showed that this theory
implies that $x^\dagger$ exists for every real $x$ and conjectured a
positive answer.  Steel gave a negative solution in 1997, and Schindler
later determined the precise consistency strength of the theory; see
\cite{schindler_delfino,Schindler2002CoreModel,CaicedoLoeweDelfino2019}.
The present paper studies
the corresponding finite recognition problem.

Friedman and Schindler isolated finite-level versions of this question in
\cite{friedman2003universally}.  Motivated by their formulation, we consider
the following exact-placement local form.

\begin{question}[Friedman--Schindler local problem, exact-placement form]
\label{question:fs-local}
Let $2\leq n<\omega$.  Suppose that every boldface
$\boldsymbol\Sigma^1_{n+2}$ set of reals is Lebesgue measurable and has the
Baire property, and suppose that the reals admit a
$\Delta^1_{n+3}$-definable well-order.  Must
$\boldsymbol\Delta^1_{n+1}$-determinacy hold?
\end{question}

Friedman and Schindler obtained a negative answer to a nearby local problem.
Working over inner models with finitely many strong cardinals, they produced
forcing extensions in which the relevant finite level of the projective
hierarchy is universally Baire and the reals admit a projective well-order.
Their construction gives the well-order at the desired level from a specific
coding real as a parameter, and parameter-free at a higher projective level.
Thus the exact lightface placement in Question~\ref{question:fs-local}, with
a $\Delta^1_{n+3}$ well-order and no additional real parameter, remained
open.

We resolve this exact-placement problem and add an adjacent uniformization
conclusion.  Throughout the introduction, uniformization statements written
without the word ``boldface'' are meant in the lightface sense: a pointclass
$\Gamma$ has uniformization if every lightface relation in $\Gamma$ has a
uniformizing function whose graph is again in $\Gamma$.  Boldface regularity
assertions are stated explicitly.

\begin{theorem}\label{thm:intro-main}
Let $1\leq n<\omega$, and assume that $M_n$, the canonical inner model with
$n$ Woodin cardinals, exists.  Then there is a forcing extension $W^\ast$ of
$M_n$ preserving all cardinals such that $2^{\aleph_0}=\aleph_2$ and:
\begin{enumerate}
   \item every boldface $\boldsymbol\Sigma^1_{n+2}$ set of reals is
   Lebesgue measurable;
   \item every boldface $\boldsymbol\Sigma^1_{n+2}$ set of reals has the
   Baire property;
   \item the $\Sigma^1_{n+2}$-uniformization property holds;
   \item the $\Pi^1_{n+3}$-uniformization property holds;
   \item the reals admit a $\Delta^1_{n+3}$-definable well-order.
\end{enumerate}
\end{theorem}

For $n\geq 2$, Theorem~\ref{thm:intro-main} gives a negative answer to
Question~\ref{question:fs-local}.  The proof records that the final extension
does not satisfy $\boldsymbol\Delta^1_{n+1}$-determinacy.  Hence the
regularity-and-well-ordering window suggested by $M_{n+1}$ can already be
realized by forcing over $M_n$.  Moreover, the construction gives the two
additional lightface uniformization conclusions
$\Sigma^1_{n+2}$-uniformization and $\Pi^1_{n+3}$-uniformization.

The first case, $n=1$, is not an instance of
Question~\ref{question:fs-local} as stated, but it illustrates the same
phenomenon at the lowest level covered by the theorem.  Starting from $M_1$,
we obtain a model in which all boldface $\boldsymbol\Sigma^1_3$ sets are
Lebesgue measurable and have the Baire property,
$\Sigma^1_3$-uniformization and $\Pi^1_4$-uniformization hold, and the reals
have a $\Delta^1_4$-definable well-order.  Thus adjacent
$\Pi$-uniformization can coexist with a projective well-order of the reals.

\paragraph{Relation with the companion paper.}
The present paper is the $\Pi$-side construction for the local Delfino
problem.  Its additional conclusion, beyond the exact-placement negative
answer, is the adjacent $\Pi^1_{n+3}$-uniformization theorem.  The forcing
mechanism introduced for this purpose is the derivative hierarchy of
allowable hybrid forcings, which controls the stability of proposed values
for sections of $\Pi^1_{n+3}$ relations through all later allowable
extensions.

The companion paper~\cite{HoffelnerLocalDelfinoI} treats a different
direction.  It obtains the same $\boldsymbol\Sigma^1_{n+2}$ regularity and
a $\Delta^1_{n+3}$ well-order, but combines them with Martin's Axiom and
with a global tail of $\Sigma$-uniformization, namely
$\Sigma^1_{n+2+m}$-uniformization for every $m\in\omega$.  That
construction is based on coding over codes.  The two papers share the same
exact-placement background, but their forcing mechanisms are separate: the
present construction is based on the local coding predicate and the
derivative hierarchy which yields the $\Pi^1_{n+3}$-uniformization
conclusion.


\paragraph{Proof overview.}
The construction has four main components.  First, over $M_n$ we add branches
through an $M_n$-definable independent sequence of Suslin trees and define a
local coding predicate $\Phi_n$ of complexity $\Sigma^1_{n+3}$.  The predicate
is arranged to be exact on the coding tags used for the well-order and for
the uniformization requirements.

Second, we use a hybrid forcing which combines a countable-support product of
$M_n$-Cohen reservoir coordinates with a finite-support c.c.c. iteration.
The reservoir coordinates provide fresh coding areas, while the iteration
meets the coding, uniformization, and regularity requirements.  This hybrid
form refines the forcing methods from \cite{HOFFELNER2023103292}.

Third, the adjacent $\Pi^1_{n+3}$-uniformization conclusion is obtained from
a derivative hierarchy of allowable hybrid forcings, refining the method of
\cite{Pi_Uniformization}.  At each derivative step, the construction tests
whether a proposed value for a section of a $\Pi^1_{n+3}$ relation is stable
under all further allowable extensions or can still be destroyed.  The
process is iterated to a stable class of allowable forcings, and the final
iteration uses only forcings from this stable class.

The same coding predicate gives the $\Delta^1_{n+3}$ well-order.  For each
pair of reals, the construction codes the relative order of their canonical
localized presentations.  Exactness of $\Phi_n$ on the well-order tags
ensures that exactly one of the two possible order tags is coded.  This gives
both a $\Sigma^1_{n+3}$ and a $\Pi^1_{n+3}$ definition of the order.

Finally, the regularity bookkeeping is kept separate from the coding
bookkeeping.  The final iteration places ordinary Cohen stages cofinally
often and, for every real parameter $a$, places relative random,
measure-amoeba, and category-amoeba stages over $L[T_{n+1},a]$ cofinally
often.  The covering statements obtained from the amoeba stages, together
with the small-generic absoluteness of the weakly homogeneous
Martin--Solovay tree $T_{n+1}$, yield Lebesgue measurability and the Baire
property for all boldface $\boldsymbol\Sigma^1_{n+2}$ sets by Hjorth's
argument.  The $\Sigma^1_{n+2}$-uniformization conclusion is obtained
separately from Steel's relativized capture analysis of $M_n(s)$ in the
small generic extensions used in the construction.

\paragraph{Organization.}
Section~\ref{sec:canonical-inner-models} recalls the inner-model-theoretic
background: the relevant fragments of Steel's comparison theory for $M_1$
and $M_n$, the recovery of the initial segments used in the construction,
and the Martin--Solovay trees $T_{n+1}$.  The next sections define the local
coding predicate, the hybrid forcings, and the derivative hierarchy of
allowable forcings.  The main construction is first carried out in full
detail over $M_1$.  The final section records the uniform lift to $M_n$ and
proves Theorem~\ref{thm:intro-main} in full generality.


\section{Canonical inner models with Woodin cardinals and the trees $T_n$}\label{sec:canonical-inner-models}

\subsection{The canonical inner model with one Woodin cardinal}\label{subsec:m1-background}

We recall the fragment of Steel's comparison theory which will be used later.  The ground model for the main construction is the canonical proper class mouse $M_1$, the minimal iterable proper class premouse with one Woodin cardinal.  We use the notation and comparison conventions of Steel's outline of inner model theory, and the projective definability analysis of Steel's work on projectively well-ordered inner models; see \cite{Steel3,Steel2}.  In particular, every proper initial segment of $M_1$ is $1$-small and $\omega$-sound, and the order of construction gives a canonical $\Delta^1_3$ well-order of the reals of $M_1$.

A premouse $\mathcal M$ is \emph{$1$-small above $\eta$} if whenever $E$ is an extender on the $\mathcal M$-sequence and $\eta<\operatorname{crit}(E)$, the initial segment $\mathcal J^{\mathcal M}_{\operatorname{crit}(E)}$ has no Woodin cardinal above $\eta$.  We say that $\mathcal M$ is \emph{$1$-small} if it is $1$-small above $0$.

We also recall the weak iterability notion used by Steel at the first odd level.  Let $\mathcal T$ be an $\omega$-maximal putative iteration tree on $\mathcal M$, let $b$ be a maximal branch through $\mathcal T$, and let $\alpha$ be a countable ordinal.  The branch $b$ is \emph{$\alpha$-good} if whenever $\mathcal N$ is either $\mathcal M_b^{\mathcal T}$ itself, or the $\alpha$-th linear iterate of an initial segment $\mathcal P\trianglelefteq \mathcal M_b^{\mathcal T}$ by one extender on the $\mathcal P$-sequence and its images, then either $\mathcal N$ is well-founded or $\alpha\in\operatorname{wfp}(\mathcal N)$.  A countable premouse $\mathcal M$ is \emph{$\Pi^1_2$-iterable} if player II wins the corresponding one-round weak iteration game: player I plays a countable putative $\omega$-maximal tree together with a countable ordinal $\alpha$, and player II either accepts a last well-founded model or plays a maximal $\alpha$-good branch.  Steel proves that this iterability condition is $\Pi^1_2$ in the codes.  For $1$-small mice it is the $n=1$ instance of the general $\Pi_n$-iterability analysis; see \cite[Lemma~1.7]{Steel2}.

We now fix the local notation for limit length trees.

\begin{definition}[The common part of a limit tree]\label{def:delta-common-part}
Let $\mathcal T$ be a $k$-maximal iteration tree of limit length on a premouse $\mathcal M$, where $k\leq\omega$.
\begin{enumerate}
   \item We set
   \[
      \delta(\mathcal T)=
      \sup\{\operatorname{lh}(E_\xi^{\mathcal T})\mid
             \xi+1<\operatorname{lh}(\mathcal T)\}.
   \]
   \item $\mathcal M(\mathcal T)$ denotes the common part of the models along $\mathcal T$ below $\delta(\mathcal T)$, i.e. the unique passive premouse $\mathcal P$ of height $\delta(\mathcal T)$ such that, for every extender $E_\xi^{\mathcal T}$ used in $\mathcal T$, $\mathcal P$ agrees with the corresponding model of the tree below $\operatorname{lh}(E_\xi^{\mathcal T})$.
\end{enumerate}
\end{definition}

We shall use the following form of Steel's branch uniqueness theorem, often called the zipper lemma.

\begin{theorem}[Steel's zipper lemma]\label{thm:zipper-local}
Let $\mathcal T$ be a $k$-maximal iteration tree of limit length on a premouse $\mathcal M$, where $k\leq\omega$, and let $b,c$ be distinct cofinal branches through $\mathcal T$.  Put $\delta=\delta(\mathcal T)$.  Suppose that $A\subseteq\delta$ and that $A,\delta\in\operatorname{wfp}(\mathcal M_b^{\mathcal T})\cap\operatorname{wfp}(\mathcal M_c^{\mathcal T})$.  Then one of the two branch models satisfies
\[
   \exists\kappa<\delta\,
   (\kappa\text{ is }A\text{-strong up to }\delta).
\]
\end{theorem}

\begin{definition}[$\mathcal Q$-structures]\label{def:q-structures-local}
Let $\mathcal T$ be a $k$-maximal iteration tree of limit length on a premouse $\mathcal M$, where $k\leq\omega$, and let $b$ be a cofinal well-founded branch through $\mathcal T$.  The $\mathcal Q$-structure $\mathcal Q(b,\mathcal T)$ is the least initial segment of $\mathcal M_b^{\mathcal T}$, if it exists, which either sees that $\delta(\mathcal T)$ is not Woodin over the common part $\mathcal M(\mathcal T)$, or projects strictly below $\delta(\mathcal T)$ at an allowed finite degree.  More explicitly, $\mathcal Q(b,\mathcal T)=\mathcal J_\gamma^{\mathcal M_b^{\mathcal T}}$ for the least $\gamma$ such that either
\[
   \mathcal J_{\gamma+1}^{\mathcal M_b^{\mathcal T}}
      \models ``\delta(\mathcal T)\text{ is not Woodin}'',
\]
or, for some $m<\omega$ allowed by the degree of the tree,
\[
   \rho_{m+1}(\mathcal J_\gamma^{\mathcal M_b^{\mathcal T}})
      <\delta(\mathcal T).
\]
If no such $\gamma$ exists, then $\mathcal Q(b,\mathcal T)$ is undefined.
\end{definition}

The point of the preceding definition is that, in the $1$-small context, a $\mathcal Q$-structure determines at most one good cofinal branch.  If two distinct branches had the same relevant well-founded $\mathcal Q$-structure, the zipper lemma would produce strength below $\delta(\mathcal T)$, contradicting the initial segment which witnesses that $\delta(\mathcal T)$ is not Woodin.

In the next few statements, an \emph{ordinary premouse} means a premouse in Steel's usual premouse language, with no real parameter, predicate parameter, or base set added to the structure.  This is only a terminological convention distinguishing these premice from mice over a real or over another base; no additional fine-structural notion is being introduced.

\begin{lemma}[Comparison with a $\Pi^1_2$-iterable mouse]\label{lem:local-comparison-m1}
Let $\mathcal M$ and $\mathcal N$ be countable ordinary premice.  Assume that both are $\omega$-sound and project to $\omega$, that $\mathcal M\triangleleft M_1$, and that $\mathcal N$ is $1$-small and $\Pi^1_2$-iterable.  Then the comparison of $\mathcal M$ with $\mathcal N$ is successful.  Consequently
\[
   \mathcal M\trianglelefteq\mathcal N
   \quad\text{or}\quad
   \mathcal N\trianglelefteq\mathcal M.
\]
\end{lemma}

\begin{proof}
Run the usual coiteration by least disagreement, producing trees $\mathcal T$ on $\mathcal M$ and $\mathcal U$ on $\mathcal N$.  The $\mathcal M$-side is governed by the strategy inherited from $M_1$.  The $\mathcal N$-side is governed by the winning strategy witnessing $\Pi^1_2$-iterability.

There is no new issue at successor stages.  Consider a countable limit stage and suppose first that both sides have reached the same comparison height, so that $\delta(\mathcal T)=\delta(\mathcal U)$.  Let $b$ be the branch selected on the $\mathcal M$-side.  Since $\mathcal M\triangleleft M_1$ and $M_1$ is $1$-small below its Woodin, the branch model $\mathcal M_b^{\mathcal T}$ has a $\mathcal Q$-structure witnessing the relevant failure of Woodinness at this common $\delta$.  Write this $\mathcal Q$-structure as $L_\alpha(\mathcal M(\mathcal T))$, using the standard coding of the common part.

Apply the $\Pi^1_2$-iterability of $\mathcal N$ to the tree $\mathcal U$ with ordinal parameter $\alpha$.  Player II supplies a maximal $\alpha$-good branch $c$.  Because the two sides of the comparison agree below the common $\delta$, the initial segment $L_\alpha(\mathcal M(\mathcal T))$ is also the corresponding $\mathcal Q$-structure for $c$ on the $\mathcal N$-side, provided the branch model is well-founded past that structure.  If there were a second sufficiently good cofinal branch through $\mathcal U$, then the zipper lemma, applied with a code for the common $\mathcal Q$-structure as parameter $A$, would produce strength below $\delta(\mathcal U)$.  This contradicts the witness to the failure of Woodinness coded by the $\mathcal Q$-structure.  Hence the good branch through $\mathcal U$ is unique, and it is fully well-founded.

If one side has stopped using extenders, the same argument applies with the last model on the stopped side replacing the common branch model.  Its relevant initial segment supplies the $\mathcal Q$-structure which witnesses that the remaining comparison height is not Woodin, and therefore determines the unique well-founded branch on the other side.

It remains to rule out a comparison of length $\omega_1$.  Suppose that such a putative comparison existed.  Force with the L\'evy collapse $\operatorname{Col}(\omega,\omega_1)$ over the ambient model.  In the collapse extension, the putative tree becomes countable, and $\Pi^1_2$-iterability is preserved.  Thus there is a branch supplied by the $\Pi^1_2$ strategy.  By the uniqueness just proved, this branch is ordinal definable from the ground-model data of the comparison.  Homogeneity of the collapse therefore puts the branch back in the ground model, contradicting the assumption that the comparison had no branch at stage $\omega_1$.

Thus the comparison terminates below $\omega_1$.  The terminal models are linearly ordered by initial segment.  Since both premice are $\omega$-sound and project to $\omega$, the usual no-drop argument for the shorter side pulls the conclusion back to the original premice.  Hence $\mathcal M\trianglelefteq\mathcal N$ or $\mathcal N\trianglelefteq\mathcal M$.
\end{proof}

\paragraph{Preservation convention for the outer models used below.}\label{par:steel-preservation-convention}
We shall use Lemma~\ref{lem:local-comparison-m1} only in the forcing extensions which occur in this paper.  The preservation fact needed there is the following: if $\mathcal M\triangleleft M_1$ is countable and belongs to one of these extensions, then $\mathcal M$ remains $\Pi^1_2$-iterable there.  This is the realizability preservation supplied by Steel's weak iteration-game analysis for initial segments of $M_1$, applied to the proper, $\omega_1$-preserving forcing extensions used below.  We do not claim that the set of all real codes for $\Pi^1_2$-iterable premice is absolute between arbitrary outer models with the same $\omega_1$.

\begin{lemma}[Definable cofinal system of $M_1$-initial segments]
\label{lem:definable-m1-initial-segments}
Let $M_1[G]$ be an $\omega_1$-preserving forcing extension of $M_1$ of the kind fixed in Paragraph~\ref{par:steel-preservation-convention}.  In $M_1[G]$ let
\[
\begin{split}
   \mathcal I=\bigl\{\mathcal N\mid {}&
      \mathcal N\text{ is a countable ordinary premouse, }\\
      &\mathcal N\text{ is }1\text{-small, }\mathcal N\text{ is }\omega\text{-sound, }\\
      &\rho_\omega(\mathcal N)=\omega,
        \text{ and }\mathcal N\text{ is }\Pi^1_2\text{-iterable}\bigr\}.
\end{split}
\]
Then $\mathcal I$ is $\Pi^1_2$-definable in the codes.  Moreover every member of $\mathcal I$ is of the form $\mathcal J_\eta^{M_1}$ for some $\eta<\omega_1$, and
\[
   \{\eta<\omega_1\mid \mathcal J_\eta^{M_1}\in\mathcal I\}
\]
is cofinal in $\omega_1$.
\end{lemma}

\begin{proof}
The clauses saying that a real codes a countable ordinary premouse, that the premouse is $1$-small, that it is $\omega$-sound, and that $\rho_\omega=\omega$ are arithmetic, after fixing the usual coding of countable premice by reals.  By Steel's analysis of the weak iteration games, $\Pi^1_2$-iterability is a $\Pi^1_2$ condition in the code.  Hence the displayed definition of $\mathcal I$ is $\Pi^1_2$.

We next show that no nonstandard premouse enters $\mathcal I$.  Work in $M_1[G]$ and let $\mathcal N\in\mathcal I$.  Since $\omega_1$ is preserved, $\mathcal N$ has countable height below the true $\omega_1^{M_1}$.  Choose $\eta<\omega_1$ such that $\operatorname{Ord}^{\mathcal N}<\eta$ and such that $\mathcal J_\eta^{M_1}$ is $\omega$-sound and projects to $\omega$.  The initial segment $\mathcal J_\eta^{M_1}$ is an ordinary premouse and is realizable inside $M_1$; hence it is $\Pi^1_2$-iterable in the present extension by Paragraph~\ref{par:steel-preservation-convention}.  Applying Lemma~\ref{lem:local-comparison-m1} to $\mathcal J_\eta^{M_1}$ and $\mathcal N$, the alternative $\mathcal J_\eta^{M_1}\triangleleft\mathcal N$ is impossible by the choice of $\eta$.  Therefore $\mathcal N\trianglelefteq\mathcal J_\eta^{M_1}$, so $\mathcal N$ is itself an initial segment of $M_1$.

Finally, the fine structure of $M_1$ gives cofinally many $\eta<\omega_1$ such that $\mathcal J_\eta^{M_1}$ is $\omega$-sound and projects to $\omega$.  These levels are ordinary premice, $1$-small, and realizable, hence $\Pi^1_2$-iterable in the present extension by the same preservation convention.  Therefore they belong to $\mathcal I$, proving cofinality.
\end{proof}

\begin{definition}[Recovering $M_1|\omega_1$]\label{def:recover-m1-omega1}
In any $\omega_1$-preserving forcing extension of $M_1$ of the kind fixed in Paragraph~\ref{par:steel-preservation-convention}, we write
\[
   \mathcal N=M_1|\omega_1
\]
for the assertion that
\[
   \mathcal N=\bigcup\{\mathcal M\mid \mathcal M\in\mathcal I\},
\]
where $\mathcal I$ is the class from Lemma~\ref{lem:definable-m1-initial-segments}.
Equivalently, $x\in\mathcal N$ iff $x$ belongs to some countable ordinary premouse which is $1$-small, $\omega$-sound, $\Pi^1_2$-iterable, and projects to $\omega$.
\end{definition}

\begin{lemma}[The recovered initial segment]\label{lem:recover-m1-omega1}
Let $M_1[G]$ be an $\omega_1$-preserving forcing extension of $M_1$ of the kind fixed in Paragraph~\ref{par:steel-preservation-convention}.  Then Definition~\ref{def:recover-m1-omega1} defines the true initial segment
\[
   M_1|\omega_1=\mathcal J^{M_1}_{\omega_1}.
\]
Moreover the definition is uniform in all such extensions.
\end{lemma}

\begin{proof}
By Lemma~\ref{lem:definable-m1-initial-segments}, every member of $\mathcal I$ is an initial segment $\mathcal J^{M_1}_\eta$ with $\eta<\omega_1$.  Hence the union in Definition~\ref{def:recover-m1-omega1} is contained in $M_1|\omega_1$.  Conversely, the same lemma gives cofinally many $\eta<\omega_1$ with $\mathcal J_\eta^{M_1}\in\mathcal I$.  Their union is $\mathcal J^{M_1}_{\omega_1}$.  The same formula defining $\mathcal I$ is used in every such extension.  We do not assert that the same real codes belong to $\mathcal I$ in different extensions, since new reals may code new countable putative iteration trees.  Rather, applying Lemma~\ref{lem:definable-m1-initial-segments} inside the given extension identifies the union of the premice satisfying that formula with the true $\mathcal J^{M_1}_{\omega_1}$.
\end{proof}

\paragraph{Remark.}\label{rem:m1-omega1-not-absolute}
The notation in Definition~\ref{def:recover-m1-omega1} is a convention for outer models of $M_1$ which preserve $\omega_1$.  The same first-order-looking assertion, if evaluated inside an arbitrary transitive model, need not imply that the object obtained is the true $M_1|\omega_1$.  Later, whenever this definition is used inside countable auxiliary models, the relevant countable premouse is also required externally to belong to the class $\mathcal I$.


We shall need a canonical diamond sequence which is available from the same fine structure.  The argument is Jensen's proof of diamond in $L$, with Steel's condensation theorem for initial segments of $M_1$ replacing condensation for $L$.

\begin{theorem}[Steel condensation, in the form used here]\label{thm:steel-condensation-m1}
Let $\mathcal M\trianglelefteq M_1$ be an $\omega$-sound initial segment and let
\[
   \pi:\bar N\rightarrow\mathcal M
\]
be the inverse of the transitive collapse of a sufficiently elementary substructure of $\mathcal M$.  Suppose that the critical point of $\pi$ is the relevant standard projectum of $\bar N$.  Then either
\begin{enumerate}
   \item $\bar N\trianglelefteq\mathcal M$, or
   \item $\bar N$ is an initial segment of a degree-zero ultrapower of an initial segment of $\mathcal M$ by an extender on the $M_1$-sequence whose length is that projectum.
\end{enumerate}
In the hulls used in the diamond argument below, the second alternative is impossible.  Hence the transitive collapse is an initial segment of $M_1$.
\end{theorem}

\begin{proof}
This is the standard condensation theorem for the Steel $M_1$-construction; see \cite[Theorem~5.1]{Steel3}.  We only spell out why the ultrapower alternative does not occur in the present application.  The hulls below are chosen so that their collapse has projectum equal to its internal $\omega_1$.  If the second alternative held, there would be an extender on the $M_1$-sequence indexed exactly at this internal $\omega_1$.  The lower part below that index is the collapse of the corresponding lower part of the hull and sees the index as the successor cut reached by the construction.  An extender indexed there would make the index inaccessible in the relevant extender model.  This contradicts the fact that the collapsed lower part computes it as its $\omega_1$.  Therefore only the initial-segment alternative remains.
\end{proof}

\begin{lemma}[A canonical $M_1$-diamond sequence]\label{lem:m1-diamond}
There is a sequence
\[
   \vec D=\langle D_\alpha\mid \alpha<\omega_1\rangle\in M_1
\]
with $D_\alpha\subseteq\alpha$ for all $\alpha<\omega_1$ such that
\[
   M_1\models ``\vec D\text{ is a }\diamondsuit_{\omega_1}\text{-sequence}''.
\]
Moreover $\vec D$ is uniformly definable over $M_1|\omega_1$ from the canonical order of construction of $M_1$.
\end{lemma}

\begin{proof}
Work in $M_1$.  Use the canonical well-order $<_{M_1}$ of the construction to define $\vec D$ recursively.  At a limit ordinal $\alpha<\omega_1$, suppose $\langle D_\beta\mid\beta<\alpha\rangle$ has been defined.  If there is a pair $(A,C)$ such that $A\subseteq\alpha$, $C\subseteq\alpha$ is club in $\alpha$, and
\[
   \forall\beta\in C\, (D_\beta\neq A\cap\beta),
\]
then let $(A_\alpha,C_\alpha)$ be the $<_{M_1}$-least such pair and set $D_\alpha=A_\alpha$.  If there is no such pair, set $D_\alpha=\emptyset$.  At successor stages we may again set $D_\alpha=\emptyset$.

The recursion is carried out over $M_1|\omega_1$.  Indeed, all objects considered at stage $\alpha$ are subsets of the countable ordinal $\alpha$ in $M_1$, and the order used to choose the least pair is the restriction of the canonical $M_1$ construction order.  Hence the resulting sequence is uniformly definable over $M_1|\omega_1$.

It remains to verify that $\vec D$ is a diamond sequence.  Suppose not.  Let $(A,C)$ be the $<_{M_1}$-least counterexample, so $A\subseteq\omega_1$, $C\subseteq\omega_1$ is club, and
\[
   \forall\alpha\in C\,(D_\alpha\neq A\cap\alpha).
\]
Choose an $\omega$-sound initial segment $\mathcal J^{M_1}_\theta$ containing $A$, $C$, and the sequence $\vec D$, and then take a countable elementary substructure
\[
   X\prec \mathcal J^{M_1}_\theta
\]
with $A,C,\vec D\in X$ and with $\alpha=X\cap\omega_1\in C$.  Let
\[
   \pi:\bar X\rightarrow X
\]
be the inverse of the transitive collapse.  By Theorem~\ref{thm:steel-condensation-m1}, $\bar X$ is an initial segment of $M_1$.  Consequently the recursive construction of $\vec D$ inside $\bar X$ is exactly the initial part
\[
   \langle D_\beta\mid\beta<\alpha\rangle.
\]
Moreover the collapse sends $A$ to $A\cap\alpha$ and $C$ to $C\cap\alpha$, and by elementarity $\bar X$ sees $(A\cap\alpha,C\cap\alpha)$ as the $<_{M_1}$-least witness that the previous sequence is not diamond on $\alpha$.  Therefore the recursion at stage $\alpha$ gives
\[
   D_\alpha=A\cap\alpha.
\]
Since $\alpha\in C$, this contradicts the choice of $(A,C)$.  Thus no counterexample exists, and $\vec D$ is a $\diamondsuit_{\omega_1}$-sequence in $M_1$.
\end{proof}

We fix once and for all the $<_{M_1}$-least sequence $\vec D$ satisfying Lemma~\ref{lem:m1-diamond}.  In later sections this sequence will be used to define a canonical $\omega_1$-sequence of independent Suslin trees over $M_1$.  The construction of those trees is postponed until the point where the coding machinery needs them.

\subsection{The canonical inner model with $n$-many Woodin cardinals}\label{subsec:mn-background}

We now record the higher-level analogue of Subsection~\ref{subsec:m1-background}.  Throughout this subsection $1\leq n<\omega$ is fixed, and $M_n$ denotes the canonical minimal proper class premouse with $n$ Woodin cardinals, in the sense of Steel's construction of tame mice with full background extenders.  Thus $M_0=L$, and for $n>0$ the model $M_n$ is obtained from the Steel background construction by stopping at the first failure of $n$-smallness, or as the limit of the $n$-small construction if no such failure occurs.  We shall use only the following standard consequences of Steel's analysis.

First, if there are $n$ Woodin cardinals, then $M_n$ exists and satisfies that there are $n$ Woodin cardinals.  Secondly, every proper initial segment of $M_n$ is $n$-small and $\omega$-sound.  Thirdly, the order of construction of $M_n$ gives a canonical construction well-order of the reals of $M_n$, and Steel's projective analysis shows that this well-order is $\Delta^1_{n+2}$ over the reals of $M_n$; see \cite{Steel2,Steel3}.  The case $n=1$ is exactly the situation isolated in the previous subsection.

\begin{definition}[$n$-smallness]\label{def:n-smallness-general}
Let $\mathcal M$ be a premouse and let $\eta<\operatorname{Ord}^{\mathcal M}$.  We say that $\mathcal M$ is \emph{$n$-small above $\eta$} if whenever $E$ is an extender on the $\mathcal M$-sequence and
\[
   \eta<\operatorname{crit}(E),
\]
then the initial segment $\mathcal J^{\mathcal M}_{\operatorname{crit}(E)}$ does not have $n$ Woodin cardinals above $\eta$.  We say that $\mathcal M$ is \emph{$n$-small} if it is $n$-small above $0$.
\end{definition}

We shall also use Steel's finite-level iterability condition.  The exact definition depends on the parity of $n$.  For even $n$, one uses the $n$-round weak iteration game in which the branches played by player II are required to be sufficiently definable over the trees played by player I.  For odd $n$, one adds the corresponding $\alpha$-goodness requirement at the last round.  In both cases the definition is arranged so that realizable mice satisfy it, and so that it is projectively simple in the codes.

\begin{definition}[Steel $\Pi_n$-iterability]\label{def:pin-iterability}
Let $\mathcal M$ be a countable premouse and let $\eta<\operatorname{Ord}^{\mathcal M}$ be a cutpoint.  We say that $\mathcal M$ is \emph{$\Pi_n$-iterable above $\eta$} if player II has a winning strategy in Steel's game $\mathcal G(\mathcal M,\eta,n)$ for $\Pi_n$-iterability above $\eta$.  We say simply that $\mathcal M$ is \emph{$\Pi_n$-iterable} if it is $\Pi_n$-iterable above $0$.
\end{definition}

For later reference we isolate the three consequences of the definition which are used in this paper.

\begin{fact}[Steel]\label{fact:steel-pin-iterability}
For each fixed $1\leq n<\omega$ the following hold.
\begin{enumerate}
   \item The relation
   \[
      ``x\text{ codes a countable premouse which is }\Pi_n\text{-iterable}''
   \]
   is $\boldsymbol\Pi^1_{n+1}$ in the real code $x$.
   \item If $\mathcal M$ is a countable initial segment of $M_n$, then $\mathcal M$ is $\Pi_n$-iterable in the Steel sense, in all outer models considered in this paper in which the relevant realizability argument is preserved.
   \item If $\mathcal M$ is countable, $n$-small and realizable, and $\mathcal N$ is countable, $n$-small and $\Pi_n$-iterable, then the Steel comparison of $\mathcal M$ and $\mathcal N$ is successful, provided the two mice have the same lower part.  In the applications below the lower part is empty, so this is the $\eta=0$ case.
\end{enumerate}
\end{fact}

The next lemma is the direct analogue of Lemma~\ref{lem:local-comparison-m1}.  We state it separately because it is the form in which it will be used later.

\begin{lemma}[Comparison with a $\Pi_n$-iterable mouse]\label{lem:local-comparison-mn}
Let $\mathcal M$ and $\mathcal N$ be countable ordinary premice.  Assume that both are $\omega$-sound and project to $\omega$, that $\mathcal M\triangleleft M_n$, and that $\mathcal N$ is $n$-small and $\Pi_n$-iterable.  Then the comparison of $\mathcal M$ with $\mathcal N$ is successful.  Consequently
\[
   \mathcal M\trianglelefteq\mathcal N
   \quad\text{or}\quad
   \mathcal N\trianglelefteq\mathcal M .
\]
\end{lemma}

\begin{proof}
This is the $\eta=0$ instance of Steel's comparison theorem for $n$-small $\Pi_n$-iterable mice.  Since $\mathcal M\triangleleft M_n$, the premouse $\mathcal M$ is realizable in the background construction of $M_n$.  Since $\mathcal N$ is $n$-small and $\Pi_n$-iterable, the $\mathcal N$-side has exactly the amount of iterability required by Steel's comparison lemma.  Because both premice are ordinary premice in the parameter-free sense fixed above, the lower parts agree trivially.

The proof is the usual least-disagreement coiteration.  The $\mathcal M$-side uses the realization strategy inherited from the construction of $M_n$, while the $\mathcal N$-side uses the $\Pi_n$-iterability strategy.  At limit stages, the good branch is characterized by the corresponding $\mathcal Q$-structure: for $n=1$ this is the branch uniqueness argument described in Subsection~\ref{subsec:m1-background}, and for $n>1$ the assertion is proved by induction on $n$, because the $\mathcal Q$-structure appearing at a limit stage is $(n-1)$-small and $\Pi_{n-1}$-iterable above the new cutpoint.  If two distinct good branches were available, Steel's zipper argument would produce the forbidden strength below the comparison height, contradicting $n$-smallness at the relevant level.  Thus the comparison has well-founded branches and reaches terminal models which are linearly ordered by initial segment.

Finally, since both premice are $\omega$-sound and project to $\omega$, the standard no-drop pullback argument gives the initial-segment relation already between the original premice.  Hence $\mathcal M\trianglelefteq\mathcal N$ or $\mathcal N\trianglelefteq\mathcal M$.
\end{proof}

\paragraph{Preservation convention for the outer models used below.}\label{par:steel-mn-preservation-convention}
As in the $M_1$ case, we shall only use the preceding comparison in forcing extensions which occur in this paper.  The preservation fact needed is the following one: if $\mathcal M\triangleleft M_n$ is countable and belongs to such an extension, then $\mathcal M$ remains $\Pi_n$-iterable there.  This is the higher-level version of the realizability preservation used in Paragraph~\ref{par:steel-preservation-convention}.  We do not claim that the class of all real codes for $\Pi_n$-iterable premice is absolute between arbitrary outer models with the same $\omega_1$.

For the coding arguments we only need to recover the lower part $M_n|\omega_1$.  Thus we isolate the lower-part initial segments of $M_n$, rather than the larger class of all $n$-small $\Pi_n$-iterable mice.  This distinction is harmless for $n=1$, but it is important at higher odd levels, where one has to avoid the familiar nonstandard $\Pi_n$-iterable mice.  The lower-part restriction below excludes these objects from the definition used in the coding apparatus.

\begin{definition}[Lower-part $M_n$-approximations]\label{def:mn-lower-part-approximations}
Let $M_n[G]$ be an $\omega_1$-preserving forcing extension of $M_n$ of the kind fixed in Paragraph~\ref{par:steel-mn-preservation-convention}.  In $M_n[G]$ let $\mathcal I_n$ be the class of all $\mathcal N$ such that
\begin{enumerate}
   \item $\mathcal N$ is a countable passive ordinary premouse;
   \item $\mathcal N$ is a lower-part premouse, i.e. $\mathcal N$ has no Woodin cardinal;
   \item $\mathcal N$ is $n$-small, $\omega$-sound, and $\rho_\omega(\mathcal N)=\omega$;
   \item $\mathcal N$ is $\Pi_n$-iterable.
\end{enumerate}
\end{definition}

\begin{lemma}[Definable cofinal system of $M_n$-initial segments]\label{lem:definable-mn-initial-segments}
Let $M_n[G]$ be an $\omega_1$-preserving forcing extension of $M_n$ of the kind fixed above.  Then $\mathcal I_n$ is $\boldsymbol\Pi^1_{n+1}$-definable in the codes.  Moreover every member of $\mathcal I_n$ is of the form $\mathcal J^{M_n}_\eta$ for some $\eta<\omega_1$, and
\[
   \{\eta<\omega_1\mid \mathcal J^{M_n}_\eta\in\mathcal I_n\}
\]
is cofinal in $\omega_1$.
\end{lemma}

\begin{proof}
The clauses saying that a real codes a countable passive ordinary premouse, that the premouse is lower-part, $n$-small, $\omega$-sound, and satisfies $\rho_\omega=\omega$, are arithmetic in the usual code for countable premice.  By Fact~\ref{fact:steel-pin-iterability}, the $\Pi_n$-iterability clause is $\boldsymbol\Pi^1_{n+1}$.  Therefore $\mathcal I_n$ is $\boldsymbol\Pi^1_{n+1}$ in the codes.

We next prove that no nonstandard lower-part premouse enters $\mathcal I_n$.  Work in $M_n[G]$ and let $\mathcal N\in\mathcal I_n$.  Since $\omega_1$ is preserved, the height of $\mathcal N$ is below the true $\omega_1^{M_n}$.  Choose $\eta<\omega_1$ such that $\operatorname{Ord}^{\mathcal N}<\eta$ and such that $\mathcal J^{M_n}_\eta$ is passive, lower-part, $\omega$-sound, and projects to $\omega$.  The fine structure of $M_n$ gives cofinally many such $\eta$.  The initial segment $\mathcal J^{M_n}_\eta$ is realizable inside $M_n$, and hence is $\Pi_n$-iterable in the present extension by the preservation convention.

Apply Lemma~\ref{lem:local-comparison-mn} to $\mathcal J^{M_n}_\eta$ and $\mathcal N$.  The alternative $\mathcal J^{M_n}_\eta\trianglelefteq\mathcal N$ is impossible by the choice of $\eta$, because $\operatorname{Ord}^{\mathcal N}<\eta$.  Therefore
\[
   \mathcal N\trianglelefteq \mathcal J^{M_n}_\eta,
\]
and hence $\mathcal N$ is itself an initial segment of $M_n$.

Conversely, let $\eta<\omega_1$ be such that $\mathcal J^{M_n}_\eta$ is passive, lower-part, $\omega$-sound, and projects to $\omega$.  Then $\mathcal J^{M_n}_\eta$ is $n$-small and realizable in the Steel construction of $M_n$, and by the preservation convention it is $\Pi_n$-iterable in $M_n[G]$.  Thus $\mathcal J^{M_n}_\eta\in\mathcal I_n$.  Since such $\eta$ are cofinal in $\omega_1$, the cofinality assertion follows.
\end{proof}

\begin{definition}[Recovering $M_n|\omega_1$]\label{def:recover-mn-omega1}
In any $\omega_1$-preserving forcing extension of $M_n$ of the kind fixed above, we write
\[
   \mathcal N=M_n|\omega_1
\]
if
\[
   \mathcal N=\bigcup\mathcal I_n,
\]
where $\mathcal I_n$ is the class from Definition~\ref{def:mn-lower-part-approximations}, computed in that extension.
\end{definition}

\begin{lemma}[Correctness of the recovery]\label{lem:recover-mn-omega1}
Let $M_n[G]$ be an $\omega_1$-preserving forcing extension of $M_n$ of the kind fixed above.  Then Definition~\ref{def:recover-mn-omega1} defines the true initial segment
\[
   M_n|\omega_1=\mathcal J^{M_n}_{\omega_1}.
\]
Moreover the definition is uniform in all such extensions.
\end{lemma}

\begin{proof}
By Lemma~\ref{lem:definable-mn-initial-segments}, every member of $\mathcal I_n$ is an initial segment $\mathcal J^{M_n}_\eta$ with $\eta<\omega_1$.  Hence the union in Definition~\ref{def:recover-mn-omega1} is contained in $M_n|\omega_1$.  Conversely, the same lemma gives cofinally many $\eta<\omega_1$ with $\mathcal J^{M_n}_\eta\in\mathcal I_n$.  The union of these initial segments is $\mathcal J^{M_n}_{\omega_1}$.

The same formula defining $\mathcal I_n$ is used in every relevant extension.  We do not assert that the same real codes belong to $\mathcal I_n$ in different extensions, since new reals may code new countable putative iteration trees.  Rather, applying Lemma~\ref{lem:definable-mn-initial-segments} inside the given extension identifies the union of the premice satisfying the formula with the true $\mathcal J^{M_n}_{\omega_1}$.
\end{proof}

\paragraph{Remark.}\label{rem:mn-omega1-not-absolute}
The notation $M_n|\omega_1$ in Definition~\ref{def:recover-mn-omega1} is a convention for the outer models of $M_n$ used in this paper.  It should not be read as saying that an arbitrary transitive model which internally satisfies the same first-order-looking definition has computed the true $M_n|\omega_1$.  Later, whenever countable auxiliary models are used, the relevant countable premouse is also required externally to belong to $\mathcal I_n$.

The condensation and diamond arguments from the previous subsection lift without change once $M_1$ is replaced by $M_n$.

\begin{theorem}[Steel condensation for $M_n$, in the form used here]\label{thm:steel-condensation-mn}
Let $\mathcal M\trianglelefteq M_n$ be an $\omega$-sound initial segment and let
\[
   \pi:\bar N\rightarrow\mathcal M
\]
be the inverse of the transitive collapse of a sufficiently elementary substructure of $\mathcal M$.  Suppose that the critical point of $\pi$ is the relevant standard projectum of $\bar N$.  Then either
\begin{enumerate}
   \item $\bar N\trianglelefteq\mathcal M$, or
   \item $\bar N$ is an initial segment of a degree-zero ultrapower of an initial segment of $\mathcal M$ by an extender on the $M_n$-sequence whose length is that projectum.
\end{enumerate}
In the hulls used in the diamond argument below, the second alternative is impossible.  Hence the transitive collapse is an initial segment of $M_n$.
\end{theorem}

\begin{proof}
This is the condensation theorem for initial segments of the Steel $M_n$-construction.  The proof is the same fine-structural argument as in the $M_1$ case, with $n$-smallness replacing $1$-smallness.  The only point needed below is the exclusion of the ultrapower alternative.  The hulls in the diamond argument are chosen so that the collapsed structure has projectum equal to its internal $\omega_1$.  If the ultrapower alternative occurred, an extender on the $M_n$-sequence would be indexed at this internal $\omega_1$.  The lower part below that index computes the index as its first uncountable cardinal, while the presence of such an extender would make it inaccessible in the relevant extender model.  This contradiction leaves only the initial-segment alternative.
\end{proof}

\begin{fact}[Steel's projective well-order of $M_n$]\label{fact:mn-projective-wellorder}
The construction order $<_{M_n}$ restricted to the reals of $M_n$ is a $\Delta^1_{n+2}$ well-order.  Equivalently, for reals $x,y\in M_n$, the assertion that $x$ is constructed before $y$ in $M_n$ is given by a $\Sigma^1_{n+2}$ formula and also by a $\Pi^1_{n+2}$ formula.
\end{fact}

We shall use this fact only as a source of canonical choices inside \(M_n\).
In particular, using \(<_{M_n}\) we fix once and for all the canonical
almost disjoint family \[D=\langle d_\xi\mid \xi<\omega_1\rangle\in M_n\]
used for almost disjoint coding, and the canonical \(\diamondsuit_{\omega_1}\)-
sequence used in the bookkeeping.  No later argument uses the well-order
to choose elements from projective sections in the final model.



\begin{lemma}[A canonical $M_n$-diamond sequence]\label{lem:mn-diamond}
There is a sequence
\[
   \vec D^{\,n}=\langle D^n_\alpha\mid \alpha<\omega_1\rangle\in M_n
\]
with $D^n_\alpha\subseteq\alpha$ for all $\alpha<\omega_1$ such that
\[
   M_n\models ``\vec D^{\,n}\text{ is a }\diamondsuit_{\omega_1}\text{-sequence}''.
\]
Moreover $\vec D^{\,n}$ is uniformly definable over $M_n|\omega_1$ from the canonical order of construction of $M_n$.
\end{lemma}

\begin{proof}
Work in $M_n$.  Use the canonical well-order $<_{M_n}$ of the construction to define $\vec D^{\,n}$ recursively.  At a limit ordinal $\alpha<\omega_1$, suppose $\langle D^n_\beta\mid\beta<\alpha\rangle$ has already been defined.  If there is a pair $(A,C)$ such that $A\subseteq\alpha$, $C\subseteq\alpha$ is club in $\alpha$, and
\[
   \forall\beta\in C\,(D^n_\beta\neq A\cap\beta),
\]
then let $(A_\alpha,C_\alpha)$ be the $<_{M_n}$-least such pair and set $D^n_\alpha=A_\alpha$.  If there is no such pair, set $D^n_\alpha=\emptyset$.  At successor stages we again set $D^n_\alpha=\emptyset$.

The recursion is carried out over $M_n|\omega_1$, because at stage $\alpha$ all relevant objects are subsets of the countable ordinal $\alpha$ in $M_n$, and the order used to choose the least pair is the restriction of the canonical construction order of $M_n$.

Suppose toward a contradiction that $\vec D^{\,n}$ is not a diamond sequence in $M_n$.  Let $(A,C)$ be the $<_{M_n}$-least counterexample, so $A\subseteq\omega_1$, $C\subseteq\omega_1$ is club, and
\[
   \forall\alpha\in C\,(D^n_\alpha\neq A\cap\alpha).
\]
Choose an $\omega$-sound initial segment $\mathcal J^{M_n}_\theta$ containing $A$, $C$, and the sequence $\vec D^{\,n}$, and take a countable elementary substructure
\[
   X\prec \mathcal J^{M_n}_\theta
\]
with $A,C,\vec D^{\,n}\in X$ and with $\alpha=X\cap\omega_1\in C$.  Let
\[
   \pi:\bar X\rightarrow X
\]
be the inverse of the transitive collapse.  By Theorem~\ref{thm:steel-condensation-mn}, the collapse $\bar X$ is an initial segment of $M_n$.  Therefore the recursive construction of $\vec D^{\,n}$ inside $\bar X$ is exactly the initial part
\[
   \langle D^n_\beta\mid\beta<\alpha\rangle.
\]
The collapse sends $A$ to $A\cap\alpha$ and $C$ to $C\cap\alpha$, and by elementarity $\bar X$ sees $(A\cap\alpha,C\cap\alpha)$ as the $<_{M_n}$-least witness that the previous sequence is not diamond on $\alpha$.  Hence the recursion at stage $\alpha$ gives
\[
   D^n_\alpha=A\cap\alpha.
\]
This contradicts $\alpha\in C$.  Therefore no counterexample exists, and $\vec D^{\,n}$ is a $\diamondsuit_{\omega_1}$-sequence in $M_n$.
\end{proof}

We fix once and for all the $<_{M_n}$-least sequence $\vec D^{\,n}$ satisfying Lemma~\ref{lem:mn-diamond}.  In the higher-level version of the construction, this sequence plays exactly the role played by the $M_1$-diamond sequence in the detailed $M_1$ case: it gives the canonical lower-part bookkeeping from which the corresponding sequence of independent Suslin trees and the localized coding apparatus are built.



\subsection{The trees $T_n$, weak homogeneity, and small generic absoluteness}\label{subsec:tn-weak-homogeneity}

We next fix the tree representations of the projective pointclasses which will be used in the regularity argument.  The main application in this paper is the case $n=2$: the Martin--Solovay tree $T_2\in M_1$ represents the universal $\boldsymbol\Sigma^1_3$ set.  Since the higher-level notation is no harder, we record the construction uniformly.  Thus, for $n\geq 2$, the tree $T_n$ will be a canonical tree belonging to $M_{n-1}$ whose projection is the universal $\boldsymbol\Sigma^1_{n+1}$ set of reals.  In the final model we shall use only the following consequence of its construction: membership in $p[T_n]$ is absolute to the small generic extensions in which the relevant real appears.

We recall the form of weak homogeneity used below.  The notation follows Steel's exposition of the Martin--Solovay theorem in terms of towers of measures~\cite{steel2009derived}.  If $Z$ is a set and $\kappa$ is a cardinal, let $\operatorname{meas}_\kappa(Z)$ be the set of $\kappa$-additive measures on $Z^{<\omega}$.  If $\mu\in\operatorname{meas}_\kappa(Z)$, its dimension is the unique $m<\omega$ such that $\mu$ concentrates on $Z^m$.  If $\mu$ has dimension $m$, $\nu$ has dimension $k$, and $m\leq k$, we say that $\nu$ projects to $\mu$ if for every $A\subseteq Z^m$,
\[
   A\in\mu
   \quad\Longleftrightarrow\quad
   \{u\in Z^k\mid u\upharpoonright m\in A\}\in\nu .
\]
A tower $\langle\mu_i\mid i<\omega\rangle$ is countably complete if, whenever $A_i\in\mu_i$ for every $i$, there is an $f\in Z^\omega$ such that $f\upharpoonright\operatorname{dim}(\mu_i)\in A_i$ for every $i$.

\begin{definition}[Weakly homogeneous tree]\label{def:weakly-homogeneous-tree}
Let $T$ be a tree on $\omega\times Z$.  For $x\in\omega^\omega$ write
\[
   T_x=\{u\in Z^{<\omega}\mid \forall k<|u|\ ((x\upharpoonright k,u\upharpoonright k)\in T)\}.
\]
We say that $T$ is \emph{$\kappa$-weakly homogeneous} if there is a countable set
\[
   \mathcal M_T\subseteq\operatorname{meas}_\kappa(Z)
\]
closed under projections such that, for every real $x$,
\[
   x\in p[T]
\]
if and only if there is a countably complete tower $\langle\mu_i\mid i<\omega\rangle$ from $\mathcal M_T$ such that each $\mu_i$ concentrates on the corresponding finite section of $T_x$.  Equivalently, $T$ is weakly homogeneous in the sense of a weak homogeneity system whose range is countable.

When defined from such a witnessing system, we call $p[T]$ a $\kappa$-weakly homogeneously Suslin set.
\end{definition}

\begin{definition}[Absolute complements and universal Baireness]\label{def:absolute-complements}
Let $T$ be a tree on $\omega\times Z$ and $S$ a tree on $\omega\times Z'$.  We say that $T$ and $S$ are \emph{$\kappa$-absolute complements} if, for every forcing extension by a partial order of size $<\kappa$,
\[
   p[T]=\omega^\omega\setminus p[S].
\]
A set of reals is \emph{$\kappa$-universally Baire} if it is the projection of a tree which has a $\kappa$-absolute complement.
\end{definition}

The connection between the preceding two notions is the Martin--Solovay tree construction.  If $\bar\mu$ is a weak homogeneity system for $T$ and $\Theta$ is sufficiently large, the Martin--Solovay tree
\[
   \operatorname{ms}(\bar\mu,\Theta)
\]
searches for a coherent descending sequence of ordinal ranks through the ultrapowers determined by the measures in $\bar\mu$.  A branch through this Martin--Solovay tree is precisely a continuous certificate that the corresponding tower is ill-founded.  The following is the form of the Martin--Solovay theorem used in the rest of the paper~\cite[Theorem~2.20 and Corollary~2.21]{steel2009derived}.

\begin{theorem}[Martin--Solovay, Steel's formulation]\label{thm:ms-weak-homogeneous-ub}
Let $T$ be $\kappa$-weakly homogeneous via a weak homogeneity system $\bar\mu$, and let $\Theta>|T|^+$.  Then $T$ and $\operatorname{ms}(\bar\mu,\Theta)$ are $\kappa$-absolute complements.  In particular, $p[T]$ is $\kappa$-universally Baire.  Moreover, in every forcing extension by a partial order of size $<\kappa$, the measures in $\bar\mu$ lift to a weak homogeneity system witnessing the same assertion for the same ground-model tree $T$.
\end{theorem}

\begin{proof}
This is the Martin--Solovay theorem for weakly homogeneous trees.  The key point is that a weak homogeneity system gives, for each real $x$, a countable tree of possible measure towers.  If one of these towers is countably complete, then $x\in p[T]$.  If no such tower is countably complete, the associated ill-foundedness is witnessed uniformly by a descending sequence of ordinal ranks; these ranks form a branch through $\operatorname{ms}(\bar\mu,\Theta)_x$, provided $\Theta$ is chosen above the size of the relevant tree.

Small forcing preserves the measure towers in the required way: the measures extend canonically to the forcing extension, and functions in the extension are represented modulo the lifted measures by ground-model functions.  Consequently the same Martin--Solovay tree remains the complementing tree in every $<\kappa$-generic extension.  This is exactly Steel's proof of the weakly homogeneous case of the Martin--Solovay theorem.
\end{proof}

We shall use Theorem~\ref{thm:ms-weak-homogeneous-ub} only in a localized form.  Let $M$ be one of the mice $M_{n-1}$ and let $\delta$ be its least Woodin cardinal.  If $T\in M$ is $\kappa$-weakly homogeneous in $M$ for every $\kappa<\delta$, then for every forcing $\mathbb P\in M$ of $M$-cardinality $<\delta$, and every $M$-generic $G\subseteq\mathbb P$, the same tree $T$ and its Martin--Solovay complement compute the same projective set in $M[G]$.  Equivalently, for reals $x\in M[G]$, membership in $p[T]$ can be checked by the ground-model tree $T$ and is not changed by passing to further forcing extensions of size below the relevant completeness bound.

\begin{corollary}[Small generic absoluteness for a fixed weakly homogeneous tree]\label{cor:small-generic-absoluteness-fixed-tree}
Let $M$ be a transitive model containing a $\kappa$-weakly homogeneous tree $T$ and a witnessing system $\bar\mu$.  Let $S=\operatorname{ms}(\bar\mu,\Theta)$ for sufficiently large $\Theta$.  If $G$ is generic over $M$ for a forcing of $M$-cardinality $<\kappa$, then in $M[G]$,
\[
   p[T]=\omega^\omega\setminus p[S].
\]
In particular, if a projective formula $\varphi(x)$ is represented over $M$ by the complementing pair $(T,S)$, then for every real $x\in M[G]$,
\[
   M[G]\models \varphi(x)
   \quad\Longleftrightarrow\quad
   x\in p[T]^{M[G]} .
\]
\end{corollary}

\begin{proof}
The first assertion is Theorem~\ref{thm:ms-weak-homogeneous-ub}.  For the final assertion, the projective formula is represented in $M$ by the pair $(T,S)$, and the pair remains an absolute complementing pair in $M[G]$.  Thus exactly one of the trees has a branch over the real $x$ in $M[G]$, and this is the same truth value assigned by the projective definition in the small generic extension.
\end{proof}

\paragraph{Remark.}\label{rem:localized-smallness-of-absoluteness}
Later forcing notions need not be regarded globally as small over the relevant mouse.  What is used is the local consequence: every real and every name appearing in the coding construction is read in a bounded regular subforcing, and the relevant subforcing has size below the completeness bound of the homogeneity system.  The assertion above is therefore applied in the intermediate model generated by that local support.

\subsection{Defining the canonical trees $T_n$ inside $M_{n-1}$}\label{subsec:canonical-trees-tn}

We now choose the particular trees to which the preceding subsection will be applied.  Fix $n\geq 2$.  Work in $M_{n-1}$, and let
\[
   \delta^n_0<\delta^n_1<\cdots<\delta^n_{n-2}
\]
be the Woodin cardinals of $M_{n-1}$.  We write simply $\delta_0$ for the least of them when $n$ is fixed.

Let $U_{n+1}\subseteq\omega^\omega$ be the standard universal $\boldsymbol\Sigma^1_{n+1}$ set, with the first real coordinate coding the index and the remaining coordinates coding the parameter and the argument.  We fix this universal set once and for all using the usual Moschovakis parametrization conventions for projective pointclasses~\cite{Moschovakis}.  Thus every boldface $\boldsymbol\Sigma^1_{n+1}$ set is a section of $U_{n+1}$.

The following theorem is the precise object needed later.  It is a standard consequence of the Martin--Steel analysis of projective scales, the Martin--Solovay construction, and Steel's tree-production argument for mice with finitely many Woodin cardinals~\cite{martin2008tree,Steel2,steel2009derived}.

\begin{theorem}[Canonical weakly homogeneous tree for $\boldsymbol\Sigma^1_{n+1}$]\label{thm:canonical-tn-existence}
In $M_{n-1}$ there is a tree
\[
   T_n\subseteq (\omega\times\lambda_n)^{<\omega}
\]
for some ordinal $\lambda_n$ of $M_{n-1}$ such that
\[
   p[T_n]=U_{n+1}
\]
inside $M_{n-1}$.  Moreover, for every $\kappa<\delta_0$, the tree $T_n$ is $\kappa$-weakly homogeneous in $M_{n-1}$.
\end{theorem}

\begin{proof}
We recall the construction, since the parity shift is a common source of confusion.  One begins with the projective pointclass immediately below $\boldsymbol\Sigma^1_{n+1}$.  By the periodicity theorems and the Martin--Steel scale analysis, the relevant universal set at that lower level has a scale whose associated tree is homogeneous in $M_{n-1}$.  This is the higher analogue of the familiar Martin--Solovay representation of the complete $\boldsymbol\Sigma^1_3$ set by the tree $T_2$ over $M_1$.

If the lower-level universal set occurs with the correct polarity, the tree of the scale already gives the required homogeneous representation.  If the polarity is the complementary one, one applies the Martin--Solovay construction to the homogeneity system.  The Martin--Solovay tree represents the complement by searching for coherent descending ordinal ranks in the corresponding ultrapowers.  This is the step which accounts for the odd/even alternation in the projective hierarchy.

Finally, the passage from the lower-level matrix to the universal $\boldsymbol\Sigma^1_{n+1}$ set is an existential real projection.  Homogeneous representations are stable under this operation in the weak sense: the additional real witness is absorbed into the weak choice of a countably complete tower.  Equivalently, the projection of a homogeneously Suslin representation is weakly homogeneously Suslin.  Thus one obtains a tree $T_n$ with $p[T_n]=U_{n+1}$ and with $\kappa$-complete weak homogeneity systems for every $\kappa<\delta_0$.

All objects used in this construction are chosen inside $M_{n-1}$.  The completeness of the systems below the least Woodin follows from the extender strength available in $M_{n-1}$ and the standard Steel tree-production argument.
\end{proof}

For definiteness, we make the choice canonical.

\begin{definition}[The canonical tree $T_n$]\label{def:canonical-tn}
For $n\geq 2$, $T_n$ denotes the $<_{M_{n-1}}$-least tree $T$ for which $M_{n-1}$ verifies the conclusion of Theorem~\ref{thm:canonical-tn-existence}.  Along with $T_n$ we fix the $<_{M_{n-1}}$-least coherent choice of weak homogeneity systems
\[
   \bar\mu^n_\kappa
   \qquad (\kappa<\delta^n_0)
\]
which witness that $T_n$ is $\kappa$-weakly homogeneous.  We also fix, for sufficiently large $\Theta_n$, the Martin--Solovay complement
\[
   S_n=\operatorname{ms}(\bar\mu^n_\kappa,\Theta_n)
\]
whenever a completeness bound $\kappa$ has been specified.  The value of $S_n$ may depend on the chosen bound $\kappa$, but this will never matter: in each application we choose $\kappa$ above the size of the relevant forcing.
\end{definition}

\begin{lemma}[Absoluteness of the canonical $T_n$ representation]\label{lem:tn-small-generic-absoluteness}
Let $n\geq 2$, let $\mathbb P\in M_{n-1}$ have $M_{n-1}$-cardinality $<\kappa<\delta^n_0$, and let $G\subseteq\mathbb P$ be $M_{n-1}$-generic.  Then in $M_{n-1}[G]$ the ground-model tree $T_n$ and the corresponding Martin--Solovay complement $S_n$ are complements.  Consequently, for every real $x\in M_{n-1}[G]$,
\[
   M_{n-1}[G]\models x\in U_{n+1}
   \quad\Longleftrightarrow\quad
   x\in p[T_n]^{M_{n-1}[G]} .
\]
The same equivalence remains true in all further forcing extensions of $M_{n-1}[G]$ by forcing notions of size below the remaining completeness bound.
\end{lemma}

\begin{proof}
By Definition~\ref{def:canonical-tn}, $T_n$ is $\kappa$-weakly homogeneous in $M_{n-1}$ via $\bar\mu^n_\kappa$.  Theorem~\ref{thm:ms-weak-homogeneous-ub} gives a Martin--Solovay complement $S_n$ such that $T_n$ and $S_n$ are $\kappa$-absolute complements.  Since $|\mathbb P|^{M_{n-1}}<\kappa$, the complementing relation holds in $M_{n-1}[G]$.

The tree $T_n$ represents the universal $\boldsymbol\Sigma^1_{n+1}$ set in the ground model, and the complementing pair remains absolute in the extension.  Therefore a real $x$ in the extension satisfies the projective universal formula exactly when the $T_n$-section above $x$ is ill-founded.  The final sentence is the same argument applied once more to the lifted homogeneity system.
\end{proof}

\begin{corollary}[The case used in the main construction]\label{cor:t2-sigma13-absoluteness}
In $M_1$ there is a canonical tree $T_2$ such that
\[
   p[T_2]=U_3,
\]
where $U_3$ is the universal $\boldsymbol\Sigma^1_3$ set.  If $\delta$ is the Woodin cardinal of $M_1$, then for every $\kappa<\delta$, $T_2$ is $\kappa$-weakly homogeneous.  Hence the $T_2$ representation of $\boldsymbol\Sigma^1_3$ truth is absolute to all forcing extensions of $M_1$ obtained by forcing of size $<\kappa$, and locally to all later intermediate extensions whose relevant support has size $<\kappa$.
\end{corollary}

\begin{proof}
This is Lemma~\ref{lem:tn-small-generic-absoluteness} with $n=2$.  The tree $T_2$ is the Martin--Solovay tree in the form used by Hjorth and Solovay~\cite{martin1969basis,Hjorth}: it represents the universal $\boldsymbol\Sigma^1_3$ set and carries weak homogeneity systems below the Woodin cardinal of $M_1$.
\end{proof}

\paragraph{Remark (how $T_n$ will be used).}\label{rem:tn-use-later}
The tree $T_n$ is not an additional coding device.  Its role is to make the
lower projective truth predicates robust in small generic extensions.  In the
proof of Lebesgue measurability and the Baire property, the only instance
needed is $T_2$: after a real parameter $a$ appears, the bookkeeping adds
Random and Cohen stages over bases containing $L[T_2,a]$, and amoeba stages over
the same bases cover the null and meager Borel sets coded there.  The weak
homogeneity and Martin--Solovay complementing pair allow the
$\boldsymbol\Sigma^1_3$ definition to be read correctly in those generic
extensions.

\section{The $M_1$-ground model and the $\Sigma^1_4$ coding apparatus}
\label{sec:M1-large-continuum-coding}

We now describe the ground model used for the detailed $M_1$-case and the coding predicate.  The target value of the continuum in the final extension is $\aleph_2$, and the later bookkeeping iteration will have length $\omega_2$.  Since the preliminary branch forcing is c.c.c., this is the same $\omega_2$ as in $M_1$.

\subsection{The branch extension $W$}

Using the canonical $\diamondsuit$-sequence of $M_1|\omega_1$, fix once and for all the canonical $M_1$-definable sequence
\[
   \vec S=\langle S_\xi\mid \xi<\omega_1\rangle
\]
of independent Suslin trees.  Thus for every finite set $e\subseteq\omega_1$, the product
\[
   \prod_{\xi\in e}S_\xi
\]
is again a Suslin tree.  We identify each tree with a subset of $\omega_1$ by the fixed $M_1$-definable coding of countable normal trees.  This convention is used throughout the rest of the paper.

Let
\[
   Br=\prod_{\xi<\omega_1}^{\mathrm{fs}} S_\xi
\]
be the finite-support branch product, computed in $M_1$.   By independence of $\vec S$, every finite subproduct is
Suslin, and the usual $\Delta$-system argument shows that $Br$ is c.c.c.  If
$H\subseteq Br$ is $M_1$-generic, write
\[
   b_\xi=\bigcup\{p(\xi):p\in H,\xi\in\operatorname{dom}(p)\}
\]
for the cofinal branch through $S_\xi$ added by the $\xi$-th coordinate.  We set
\[
   W=M_1[H].
\]
Thus $W$ is just the extension obtained by adding a branch through every tree in the fixed sequence $\vec S$.  Since $Br$ has size $\omega_1$ and is c.c.c., $W$ preserves cardinals and satisfies $2^{\aleph_0}=\aleph_1$.  The later length-$\omega_2$ iteration will add $\aleph_2$ many reals and will be responsible for the final value $2^{\aleph_0}=\aleph_2$.

The point of having all branches present in $W$ is that a coding forcing may use selected branches from $H$ as parameters.  The point of using the finite-support product is that every real in a local subextension depends on only countably many branch coordinates.  This support feature will later allow us to compare the information decoded from a real with the branch coordinates which were actually used to produce that real.

\subsection{Cohen coding areas}

We fix the canonical $M_1$-definable almost disjoint family of reals
\[
   D=\langle d_\xi\mid \xi<\omega_1\rangle.
\]
It will be used at the last step of the coding, when a subset of $\omega_1$ is almost-disjointly coded into a real.

We also fix the $<_{M_1}$-least bijection
\[
   \rho: ({}^{<\omega_1}2)^{M_1}\longrightarrow \omega_1
\]
or, equivalently, the corresponding $M_1$-definable bijection between countable subsets of $\omega_1$ and $\omega_1$.  Let $\mathbb C_{M_1}$ be the $M_1$-version of the $\sigma$-closed $\omega_1$-Cohen forcing: conditions are countable binary sequences from $M_1$, ordered by end-extension.  If $g\subseteq\omega_1$ is $\mathbb C_{M_1}$-generic, then every proper initial segment $g\upharpoonright\alpha$ belongs to $M_1$, and we define the coding area determined by $g$ by
\[
   C_g=\{\rho(g\upharpoonright\alpha)\mid \alpha<\omega_1\}\subseteq\omega_1.
\]
We shall say that $C\subseteq\omega_1$ is an $M_1$-Cohen coding area if $C=C_g$ for some such $g$.  The standard fusion argument for $\mathbb C_{M_1}$ gives the two facts we need: each initial segment of $C_g$ is computed in the relevant $M_1$-initial segment, and coding areas coming from mutually generic $M_1$-Cohen subsets of $\omega_1$ are almost disjoint in that their intersection is countable.

\begin{definition}[Selected branch coordinates]\label{def:selected-branch-coordinates}
For \(C\subseteq\omega_1\) and \(u\in2^\omega\), define
\[
   \operatorname{Sel}(C,u)=
   \{\tau_u(\gamma,n)\mid \gamma\in C,\ n<\omega\},
\]
where
\[
   \tau_u(\gamma,n)=
   \begin{cases}
      \omega\cdot\gamma+2n, &\text{if } n\notin u,\\
      \omega\cdot\gamma+2n+1, &\text{if } n\in u.
   \end{cases}
\]
Thus \(\operatorname{Sel}(C,u)\) is the set of branch coordinates selected by
the area \(C\) when it writes the real \(u\).
\end{definition}

\begin{lemma}[Countably many coding areas]\label{lem:countably-many-coding-areas}
Let \(I\) be countable.  Let \((C_i:i\in I)\) be \(M_1\)-Cohen coding areas
coming from mutually generic reservoir coordinates, and let \(u_i\in2^\omega\)
for \(i\in I\).  Let \(B\in[\omega_1]^\omega\), and put
\[
   U=B\cup\bigcup_{i\in I}\operatorname{Sel}(C_i,u_i).
\]
Then the following hold.
\begin{enumerate}
\item If \(C\subseteq\omega_1\) has size \(\omega_1\) and
\[
   \forall i\in I\ (C\cap C_i\text{ is bounded in }\omega_1),
\]
then there are \(\omega_1\)-many \(\gamma\in C\) such that
\[
   \{\omega\cdot\gamma+k\mid k<\omega\}\cap U=\emptyset.
\]

\item If \(i_0\in I\), \(C\cap C_{i_0}\) is unbounded in \(\omega_1\), and
\(w\in2^\omega\) satisfies \(w\ne u_{i_0}\), then, for any
\(n<\omega\) with \(w(n)\ne u_{i_0}(n)\), there are \(\omega_1\)-many
\(\gamma\in C\cap C_{i_0}\) such that
\[
   \tau_w(\gamma,n)\notin U .
\]
\end{enumerate}
\end{lemma}

\begin{proof}
For \(i\ne j\) in \(I\), the set \(C_i\cap C_j\) is bounded in \(\omega_1\).
Since \(I\) is countable, there is \(\theta_0<\omega_1\) such that
\[
   C_i\cap C_j\subseteq\theta_0
   \qquad(i\ne j\text{ in }I).
\]
Let
\[
   B^\ast=\{\gamma<\omega_1\mid
      \exists k<\omega\ (\omega\cdot\gamma+k\in B)\}.
\]
Then \(B^\ast\) is countable.

For (1), choose \(\theta_1<\omega_1\) such that
\[
   C\cap C_i\subseteq\theta_1
   \qquad(i\in I).
\]
Every \(\gamma\in C\setminus(\theta_1\cup B^\ast)\) has the required property:
it does not lie in any \(C_i\), and its \(\omega\)-block does not meet \(B\).

For (2), let \(n<\omega\) be such that \(w(n)\ne u_{i_0}(n)\).  If
\[
   \gamma\in (C\cap C_{i_0})\setminus(\theta_0\cup B^\ast),
\]
then \(\gamma\notin C_i\) for all \(i\ne i_0\), and the \(\omega\)-block of
\(\gamma\) does not meet \(B\).  At the area \(C_{i_0}\), the coordinate
selected by \(u_{i_0}\) at \(n\) is the other member of the pair
\[
   \{\omega\cdot\gamma+2n,\ \omega\cdot\gamma+2n+1\}.
\]
Hence \(\tau_w(\gamma,n)\notin \operatorname{Sel}(C_{i_0},u_{i_0})\), and it is
not in any other \(\operatorname{Sel}(C_i,u_i)\).  Thus
\(\tau_w(\gamma,n)\notin U\).  There are \(\omega_1\)-many such \(\gamma\).
\end{proof}

\subsection{Writing a real into the Suslin sequence}

Let $w\in2^\omega$ be a real in a generic extension of $W$.  Suppose first that $g$ is an $M_1$-Cohen subset of $\omega_1$ and put $C=C_g$.  For each $\gamma\in C$ we use the $\omega$-block
\[
   \langle S_{\omega\cdot\gamma+k}\mid k<\omega\rangle
\]
of the fixed sequence $\vec S$.  The bit $w(n)$ is represented by choosing the branch through exactly one of the two trees
\[
   S_{\omega\cdot\gamma+2n},\qquad S_{\omega\cdot\gamma+2n+1}.
\]
Equivalently, in the notation fixed above, the selected coordinate is
\(\tau_w(\gamma,n)\).  The branch used at the coordinate \((\gamma,n)\) is
\(b_{\tau_w(\gamma,n)}\), already present in \(W\).  Thus any inner model which can correctly compute the trees $\vec S$ and which contains, for all $\gamma\in C$ and $n<\omega$, the selected branches
\[
   \langle b_{\tau_w(\gamma,n)}\mid \gamma\in C,
      n<\omega\rangle
\]
can read the characteristic function of $w$ from the pattern
\begin{align*}
   n\in w
   &\Longleftrightarrow S_{\omega\cdot\gamma+2n+1}
      \text{ has an }\omega_1\text{-branch},\\
   n\notin w
   &\Longleftrightarrow S_{\omega\cdot\gamma+2n}
      \text{ has an }\omega_1\text{-branch},
\end{align*}
for every $\gamma\in C$.

To turn this branch pattern into a projective predicate, we do not code the branch data directly as a bare subset of $\omega_1$.  Instead we first package the relevant information into a transitive $\aleph_1$-sized model and then code that model.  Let
\[
   \mathcal B_{g,w}=\langle b_{\tau_w(\gamma,n)}
      \mid \gamma\in C_g,\ n<\omega\rangle
\]
be the indexed sequence of selected branches, coded as a set of triples
$(\gamma,n,b_{\tau_w(\gamma,n)})$.  Thus the domain of this sequence already determines $C_g$, but we shall keep $C_g$ as an explicit object for readability.

Choose a transitive model
\[
   \mathcal N=\mathcal N_{g,w}
\]
of size $\aleph_1$ such that
\begin{enumerate}
   \item $\mathcal N\models\ZFC^-+``\aleph_1\text{ exists}''$ and $\omega_1^{\mathcal N}=\omega_1$;
   \item $m_\infty=\mathcal J^{M_1}_{\omega_1}$ belongs to $\mathcal N$;
   \item $C_g$ and $\mathcal B_{g,w}$ belong to $\mathcal N$;
   \item using $m_\infty$, the model $\mathcal N$ computes the fixed sequence $\vec S$, the map $\rho$, and the almost disjoint family $D$, and it sees the branch pattern determined by $\mathcal B_{g,w}$ as coding the real $w$.
\end{enumerate}
Here $w$ is not used as a distinguished piece of the code.  It is recovered from the selected branch pattern, while $m_\infty$ supplies the definition of the Suslin sequence from which the pattern is read.

Fix a canonical code
\[
   X_{\mathcal N}=X_{g,w}\subseteq\omega_1
\]
for the structure $(\mathcal N,\in,m_\infty,C_g,\mathcal B_{g,w})$.  Concretely, choose the $<_{M_1}$-least enumeration of $\mathcal N$ of length $\omega_1$ whose first distinguished entries are $m_\infty$, $C_g$, and $\mathcal B_{g,w}$, and let $X_{\mathcal N}$ code the resulting well-founded extensional relation together with these distinguished constants.  Any model which decodes $X_{\mathcal N}$ therefore recovers the transitive collapse $\mathcal N$ and the three named objects.

We now apply David's reshaping trick to $X_{\mathcal N}$.  Choose an ordinal $\lambda$ of size $\aleph_1$ such that
\[
   L_\lambda[X_{\mathcal N}]\models \ZFC^-+``\aleph_1\text{ exists}''.
\]
Inside $L_\lambda[X_{\mathcal N}]$, take a continuous increasing sequence
\[
   \langle M_\xi\mid \xi<\omega_1\rangle
\]
of countable elementary submodels, and put
\[
   c_\xi=M_\xi\cap\omega_1,
   \qquad
   C=\{c_\xi\mid\xi<\omega_1\}.
\]
The reshaped set $Y=Y_{g,w}\subseteq\omega_1$ codes $X_{\mathcal N}$ on the odd ordinals and codes the club $C$ on the even ordinals.  More explicitly, the odd part of $Y$ is a fixed copy of $X_{\mathcal N}$, while the even part codes the increasing enumeration of $C$ by the usual interval pattern: before $c_0$ it codes a well-order of type $c_0$, between $c_\xi$ and $c_{\xi+1}$ it codes a well-order of type $c_{\xi+1}$, and it is empty on the remaining intervals.

Consequently, if $M$ is a countable transitive model of $\ZFC^-+``\aleph_1\text{ exists}''$ and $Y\cap\omega_1^M\in M$, then $M$ recognizes that $\omega_1^M\in C$.  Inside its version of $L[Y\cap\omega_1^M]$, it recovers $X_{\mathcal N}\cap\omega_1^M$ and hence a local transitive model $\bar{\mathcal N}$ of size $\aleph_1^M$.  This local model contains the corresponding initial segment $m$ of $M_1$, the local coding area, and the local selected branch sequence.  Since $m$ defines the local version of $\vec S$, the model $\bar{\mathcal N}$ can read the same branch pattern and verify the local statement that the pattern codes $w$.

Finally, over the model containing $Y$, force with the Jensen--Solovay almost disjoint coding forcing
\[
   \mathbb A_D(Y)
\]
relative to the fixed $M_1$-definable family $D$.  If $r_Y$ is the generic real, then
\[
   \alpha\in Y
   \quad\Longleftrightarrow\quad
   r_Y\cap d_\alpha\text{ is finite}
\]
for every $\alpha<\omega_1$, and this equivalence remains absolute to all outer models with the same interpretation of the initial segment of $D$.

For a real $w$ we denote by
\[
   \operatorname{Code}(w)
\]
the two-step forcing which first adds an $M_1$-Cohen subset $g$ of $\omega_1$ and then almost-disjointly codes the reshaped set $Y_{g,w}$ into a real:
\[
   \operatorname{Code}(w)=
   \mathbb C_{M_1}*\dot{\mathbb A}_D(\dot Y_{\dot g,w}).
\]
When $w$ is a name over an initial segment of the later iteration, the same formula is read as the corresponding name for this two-step forcing.

Later, once hybrid presentations have been introduced, we use the same notation
\(\operatorname{Code}(\dot u)\) for the hybrid coding operation associated with
this two-step forcing.  Thus the \(\mathbb C_{M_1}\)-factor is inserted as a
fresh reservoir product coordinate, while the almost-disjoint forcing
\(\dot{\mathbb A}_D(\dot Y_{\dot g,\dot u})\) is inserted as the corresponding
real-adding coordinate.  In other words, \(\operatorname{Code}(\dot u)\) denotes
an ordinary two-step forcing in the local description above, but denotes this
product/iteration insertion when it occurs inside a hybrid presentation.

\subsection{The predicates $\operatorname{LocalCode}$, $\Psi$, and $\Phi$}

We isolate the projective content of the preceding construction.  The notation
\[
   M\models\operatorname{LocalCode}(Y,w,m)
\]
will be used for the following first-order assertion over a transitive model $M$ of $\ZFC^-+``\aleph_1\text{ exists}''$, with $m\in M$ as a displayed parameter.  The set $Y\subseteq\omega_1^M$ is a reshaped David code: its odd part decodes a set $X_{\mathcal N}\subseteq\omega_1^M$, its even part supplies the reshaping club, and $X_{\mathcal N}$ codes a well-founded extensional structure whose transitive collapse is a model
\[
   \bar{\mathcal N}
\]
of size $\aleph_1^M$ such that
\begin{itemize}
   \item $\bar{\mathcal N}\models\ZFC^-+``\aleph_1\text{ exists}''$;
   \item $\bar{\mathcal N}$ contains the distinguished initial segment $m\in\mathcal I$ with $\omega_1^m=\omega_1^M$;
   \item using $m$, the model $\bar{\mathcal N}$ computes the local versions of $\vec S$, $\rho$, and the almost disjoint family $D$;
   \item $\bar{\mathcal N}$ contains an $M_1$-Cohen coding area $C$ computed from $m$, i.e. a set of the form
   \[
      C=\{\rho^m(g\upharpoonright\alpha)\mid \alpha<\omega_1^m\}
   \]
   for the $M_1$-Cohen subset $g\subseteq\omega_1^m$ decoded in $\bar{\mathcal N}$, and it contains an indexed sequence of selected branches
   \[
      \mathcal B=\langle c_{\gamma,n}\mid \gamma\in C,
         n<\omega\rangle;
   \]
   \item the branch sequence verifies, in $\bar{\mathcal N}$, the pattern coding $w$:
   \begin{align*}
      n\in w
      &\Longleftrightarrow
      \bar{\mathcal N}\models
      ``S^m_{\omega\cdot\gamma+2n+1}
        \text{ has an }\omega_1^m\text{-branch}'',\\
      n\notin w
      &\Longleftrightarrow
      \bar{\mathcal N}\models
      ``S^m_{\omega\cdot\gamma+2n}
        \text{ has an }\omega_1^m\text{-branch}''
   \end{align*}
   for every $\gamma\in C$ and every $n<\omega$.
\end{itemize}
Here $S^m_\xi$ denotes the $\xi$-th tree of the sequence computed from $m$.  Thus $Y$ does not merely list branches; it codes a local universe in which the relevant $M_1$-initial segment defines the Suslin sequence, and the selected branches then determine the real $w$.

Let $D^m=\langle d^m_\alpha\mid\alpha<\omega_1^m\rangle$ denote the canonical almost disjoint family as computed in $m$.  We define $\Psi(r,w)$ by the following formula:
\begin{align*}
\Psi(r,w)\quad\Longleftrightarrow\quad
\forall M\,\forall m\Big(&
   M \text{ is a countable transitive model of }
      \ZFC^-+``\aleph_1\text{ exists}'' \\
&\wedge\ r,w,m\in M \\
&\wedge\ m\in\mathcal I \\
&\wedge\ \omega_1^m=\omega_1^M \\
&\Longrightarrow
   \exists Y\in M\Big[
      Y\subseteq\omega_1^M \\
&\qquad\wedge\,
      \forall\alpha<\omega_1^M\,
      (\alpha\in Y\Longleftrightarrow r\cap d^m_\alpha
         \text{ is finite}) \\
&\qquad\wedge\,
      M\models\operatorname{LocalCode}(Y,w,m)
   \Big]\Big).
\end{align*}
Thus $\Psi(r,w)$ says that every relevant countable transitive model which has the correct $M_1$-initial segment at its $\omega_1$ decodes from $r$, using the local almost disjoint family, a reshaped code for a local transitive model whose branch sequence codes $w$.

The assertion $\Psi(r,w)$ is a $\boldsymbol\Pi^1_3$ statement.  The only non-arithmetic ingredient is the recognition of the correct $M_1$-initial segments, which is $\boldsymbol\Pi^1_2$ via the class $\mathcal I$; the remaining decoding and branch-pattern checks are absolute first-order checks inside the countable transitive model.  Therefore
\[
   \Phi(w)\quad\Longleftrightarrow\quad \exists r\,\Psi(r,w)
\]
is a $\boldsymbol\Sigma^1_4$ predicate.  We read $\Phi(w)$ as ``$w$ is coded into the fixed Suslin sequence $\vec S$''.

\begin{lemma}[The inner model decoded by a $\Phi$-witness]\label{lem:inner-model-decoded-by-Phi}
Assume that $r$ witnesses $\Phi(w)$, that is, assume that $\Psi(r,w)$ holds.  Put
\[
   m_\infty=\bigcup\mathcal I=\mathcal J^{M_1}_{\omega_1}.
\]
Then the inner model $L[r,m_\infty]$ contains the real $w$ and contains the objects decoded from $r$: a reshaped set $Y_r\subseteq\omega_1$, a model code $X_{\mathcal N,r}\subseteq\omega_1$, the transitive model $\mathcal N_r$ decoded from it, an $M_1$-Cohen coding area $C_r$, and the selected branch sequence through the trees of the fixed sequence $\vec S$.  Moreover $L[r,m_\infty]$ sees the branch pattern coding $w$; more explicitly, in $L[r,m_\infty]$ the decoded objects satisfy
\begin{align*}
   \forall \gamma\in C_r\,\forall n<\omega\quad
   n\in w
   &\Longleftrightarrow
   S_{\omega\cdot\gamma+2n+1}\text{ has an }\omega_1\text{-branch},\\
   \forall \gamma\in C_r\,\forall n<\omega\quad
   n\notin w
   &\Longleftrightarrow
   S_{\omega\cdot\gamma+2n}\text{ has an }\omega_1\text{-branch}.
\end{align*}
Thus a witness $r$ for $\Phi(w)$ does not merely certify the projective statement externally; together with $\mathcal J^{M_1}_{\omega_1}$ it generates an actual inner model in which the $w$-pattern is present.
\end{lemma}

\begin{proof}
All decoding procedures used below are the fixed canonical ones.  In particular, once $r$ and an initial segment $m\in\mathcal I$ are given, the almost-disjoint decoding below $\omega_1^m$ is unique.

First $L[r,m_\infty]$ can form the almost-disjoint decoding of $r$ with respect to the canonical family computed from $m_\infty$:
\[
   Y_r=\{\alpha<\omega_1\mid r\cap d^{m_\infty}_\alpha
      \text{ is finite}\}.
\]
This set belongs to $L[r,m_\infty]$ and is uniquely determined by $r$ and $m_\infty$.

We claim that the David decoding of $Y_r$ in $L[r,m_\infty]$ succeeds and produces the asserted model code and branch data.  Suppose not.  Then some finite part of the canonical decoding fails: either the reshaping test fails, or the odd part does not decode a well-founded extensional structure of the required kind, or the resulting collapse does not contain the required distinguished objects, or one of the branch-pattern clauses fails.  Choose a sufficiently large regular $\Theta$ and take, in the ambient universe, a countable elementary submodel
\[
   X\prec H_\Theta
\]
containing $r$, $w$, $m_\infty$, and a witness to this failure.  Let
\[
   \pi:X\longrightarrow \bar M
\]
be the transitive collapse, and put $m=\pi(m_\infty)$.  By the recovery and condensation facts for the cofinal system $\mathcal I$, the model $m$ is a member of $\mathcal I$ and
\[
   \omega_1^m=\omega_1^{\bar M}.
\]
Moreover $\bar M$ is a countable transitive model of
$\ZFC^-+``\aleph_1\text{ exists}''$, and $r,w,m\in\bar M$.

Inside $\bar M$, the almost-disjoint decoding of $r$ with respect to $D^m$ is exactly
\[
   Y_r\cap\omega_1^{\bar M}.
\]
Indeed, for every $\alpha<\omega_1^{\bar M}$, the real $d^m_\alpha$ is the corresponding initial member of the canonical almost disjoint family computed from $m_\infty$, and hence
\[
   \alpha\in Y_r\cap\omega_1^{\bar M}
   \quad\Longleftrightarrow\quad
   r\cap d^m_\alpha\text{ is finite}.
\]
The reshaping and odd/even decoding are first-order procedures over the relevant transitive model, so the failure chosen above reflects to $\bar M$.  Thus
\[
   \bar M\not\models
      \operatorname{LocalCode}(Y_r\cap\omega_1^{\bar M},w,m).
\]
But $\Psi(r,w)$ applies to the countable transitive model $\bar M$ and the initial segment $m$.  It therefore yields some $Y\in\bar M$ such that
\[
   \forall\alpha<\omega_1^{\bar M}\,
   \bigl(\alpha\in Y\Longleftrightarrow r\cap d^m_\alpha
      \text{ is finite}\bigr)
\]
and
\[
   \bar M\models\operatorname{LocalCode}(Y,w,m).
\]
By uniqueness of the almost-disjoint decoding, $Y=Y_r\cap\omega_1^{\bar M}$, a contradiction.  Hence the global decoding in $L[r,m_\infty]$ succeeds.  Let $X_{\mathcal N,r}$ be the model code decoded from the odd part of $Y_r$, let $\mathcal N_r$ be the transitive collapse of the structure coded by $X_{\mathcal N,r}$, and let $C_r$ and the selected branch sequence be the distinguished objects decoded inside $\mathcal N_r$.

It remains to see that the real read from this branch pattern is the original real $w$.  Since the global decoding has succeeded, $L[r,m_\infty]$ defines a real $w^*$ by reading the decoded branches: for $n<\omega$, put $n\in w^*$ iff, equivalently for every $\gamma\in C_r$, the decoded branch sequence witnesses that
\[
   S_{\omega\cdot\gamma+2n+1}
   \text{ has an }\omega_1\text{-branch}.
\]
The local-code clauses ensure that this definition is independent of the choice of $\gamma\in C_r$.

Suppose toward a contradiction that $w^*\neq w$, and choose $n<\omega$ such that $w^*(n)\neq w(n)$.  Again take a sufficiently large countable elementary submodel $X\prec H_\Theta$ in the ambient universe containing $r,w,m_\infty,n$ and enough of the decoded global data to witness the value of $w^*(n)$, and collapse it to $\bar M$, with $m=\pi(m_\infty)$.  As above, $\bar M$ is a relevant countable transitive model, $m\in\mathcal I$, $\omega_1^m=\omega_1^{\bar M}$, and the almost-disjoint decoding of $r$ in $\bar M$ is $Y_r\cap\omega_1^{\bar M}$.  By elementarity and the absoluteness of the decoding procedure, $\bar M$ reads the $n$-th bit of the local branch pattern as $w^*(n)$.  On the other hand, $\Psi(r,w)$ says that the same uniquely decoded set satisfies
\[
   \bar M\models
      \operatorname{LocalCode}(Y_r\cap\omega_1^{\bar M},w,m),
\]
so $\bar M$ reads the $n$-th bit of that same branch pattern as $w(n)$.  This contradicts $w^*(n)\neq w(n)$.

Therefore $w^*=w$.  Since $w^*$ is definable in $L[r,m_\infty]$ from the decoded objects, we have $w\in L[r,m_\infty]$.  The decoded objects also belong to $L[r,m_\infty]$ by construction, and the displayed branch-pattern equivalences hold there.  This proves the lemma.
\end{proof}



\subsection{Localized presentations of names}\label{subsec:localized-presentations}

We now fix the presentation convention which will be used throughout the
construction.  Recall that the first forcing over \(M_1\) is the finite-support
branch product
\[
   Br=\prod_{\xi<\omega_1}^{\mathrm{fs}}S_\xi,
\]
which adds a cofinal branch through every tree in the fixed independent
sequence \(\vec S\).  If \(H\subseteq Br\) is \(M_1\)-generic, we write
\[
   W=M_1[H].
\]
All later forcing is described over this branch extension.  A local hybrid
presentation over \(W\) will be denoted by \(\mathcal R\), and its initial
presentation up to stage \(\beta\) by \(\mathcal R_\beta\).  Once the
presentation is fixed, \(\mathbb P_\beta\) denotes the post-branch forcing over
\(W\) associated with \(\mathcal R_\beta\).  If \(\mathcal R=\mathcal R_\delta\)
is a completed presentation, then \(\mathbb P=\mathbb P_\delta\) denotes its
terminal post-branch forcing.

The corresponding full forcing over \(M_1\), including the initial branch block,
is written
\[
   Br\hybcomp\mathcal R_\beta .
\]
Here \(Br\hybcomp\mathcal R_\beta\) means: first force with \(Br\), form
\(W=M_1[H]\), and then force over \(W\) with the post-branch forcing
\(\mathbb P_\beta\) associated with \(\mathcal R_\beta\).  The symbol
\(\hybcomp\) is only notation for this composition with a post-branch hybrid
presentation; it is not an ordinary iteration symbol.  Equivalently, if
\(H\subseteq Br\) is generic, then the quotient of \(Br\hybcomp\mathcal R_\beta\)
over \(H\) is canonically isomorphic to \(\mathbb P_\beta\).

We also fix the convention for the forcings which deal with Lebesgue measure and
the Baire property; we call them the \emph{regularity forcings}.  If \(W'\) is a
transitive inner model of the current ambient universe, with the relevant
\(\omega_1\) and enough set theory to compute Borel codes, measure, category,
and the corresponding forcing notions, then
\[
   \mathbb B_{\mathrm{rand}}^{W'}
\]
denotes the Random algebra as computed in \(W'\).  Its conditions are the
positive Borel sets coded in \(W'\), modulo null equivalence as computed in
\(W'\).  We also write
\[
   \mathbb A_{\mathcal N}^{W'}
   \qquad\text{and}\qquad
   \mathbb A_{\mathcal M}^{W'}
\]
for the standard amoeba forcing for the null ideal, respectively for the meager
ideal, as computed in \(W'\).  Thus \(\mathbb A_{\mathcal N}^{W'}\) adds a
Borel null set covering every Borel null set coded in \(W'\), and
\(\mathbb A_{\mathcal M}^{W'}\) adds a Borel meager set covering every Borel
meager set coded in \(W'\).  We use fixed Borel-code presentations of these
amoeba forcings in which they are \(\sigma\)-linked.  In particular, these
forcings preserve Suslin trees from the ground model.

The only bases used for regularity stages are finite-real-parameter bases of
the form
\[
   L[T_2,a_0,\ldots,a_{k-1}].
\]
Here \(T_2\in M_1\) is the fixed Martin--Solovay tree from
Subsection~\ref{subsec:canonical-trees-tn}, and the reals
\(a_0,\ldots,a_{k-1}\) are read from the current generic extension.  When a
regularity forcing computed in such a base is used inside a larger ambient
extension, it is not recomputed in the larger universe.  The conditions and the
order are the ones computed in the displayed base.

The post-branch forcing is a hybrid of a product and an iteration.  We now give
its official syntactic form.  A \emph{resolved hybrid bookkeeping} of length
\(\beta\) is a sequence
\[
   F=\langle F(\eta)\mid \eta<\beta\rangle
\]
whose \(\eta\)-th value is decoded over the already constructed initial forcing
as one of the following tags:
\[
   \mathbf 1,
   \qquad
   \mathrm{Coh},
   \qquad
   \operatorname{Code}(\dot u_\eta),
\]
or one of the relative regularity tags
\[
   (\mathrm{Rand};\dot a^\eta_0,\ldots,\dot a^\eta_{k_\eta-1}),
   \qquad
   (\mathrm{Am}_{\mathcal N};\dot a^\eta_0,\ldots,\dot a^\eta_{k_\eta-1}),
   \qquad
   (\mathrm{Am}_{\mathcal M};\dot a^\eta_0,\ldots,\dot a^\eta_{k_\eta-1}).
\]
Here \(k_\eta<\omega\), \(\dot u_\eta\) is a name for a real over the current
initial forcing, and each \(\dot a^\eta_i\) is a name for a real over the
current initial forcing.  The word ``resolved'' is important: this bookkeeping
tells the hybrid presentation what forcing to use next.  The later allowability
recursions may start from higher-level bookkeeping data, such as a well-order
pair or a uniformization tuple, but before a hybrid stage is formed that data is
resolved into one of the tags listed above.

Given such an \(F\), we recursively build a post-branch hybrid presentation
\[
   \langle \mathcal R_\eta\mid \eta\leq\beta\rangle
\]
together with its associated post-branch forcings
\(\langle\mathbb P_\eta\mid \eta\leq\beta\rangle\).  We put
\(\mathbb P_0=\mathbf 1\).  At a successor stage
\(\eta+1\), after \(\mathbb P_\eta\) has been constructed, we first add a fresh
reservoir coordinate \(\mathbb C_\eta\), a copy of the \(M_1\)-computed
\(\omega_1\)-Cohen forcing.  Thus the preliminary forcing at this stage is
\[
   \widehat{\mathbb P}_{\eta+1}=\mathbb P_\eta\times\mathbb C_\eta,
\]
where \(\dot g_\eta\) denotes the \(\mathbb C_\eta\)-generic.  The second-step
forcing \(\dot{\mathbb Q}_\eta\) over \(\widehat{\mathbb P}_{\eta+1}\) is
determined by the decoded tag \(F(\eta)\):
\begin{enumerate}
   \item If the tag is \(\mathbf 1\), then \(\dot{\mathbb Q}_\eta\) is the
   trivial forcing.

   \item If the tag is \(\mathrm{Coh}\), then \(\dot{\mathbb Q}_\eta\) is
   ordinary Cohen forcing.

   \item If the tag is \(\operatorname{Code}(\dot u_\eta)\), then
   \[
      \dot{\mathbb Q}_\eta
      =\dot{\mathbb A}_D(\dot Y_{\dot g_\eta,\dot u_\eta}).
   \]
   This is the Jensen--Solovay almost-disjoint coding forcing which codes the
   reshaped set determined by the fresh reservoir generic and the real named by
   \(\dot u_\eta\).

   \item If the tag is
   \((\mathrm{Rand};\dot a^\eta_0,\ldots,\dot a^\eta_{k_\eta-1})\), put
   \[
      \dot W_\eta=L[T_2,\dot a^\eta_0,\ldots,
         \dot a^\eta_{k_\eta-1}],
   \]
   with all real names lifted to \(\widehat{\mathbb P}_{\eta+1}\).  The
   second-step forcing is
   \[
      \dot{\mathbb Q}_\eta
      =\dot{\mathbb B}_{\mathrm{rand}}^{\dot W_\eta},
   \]
   the Random algebra computed in the displayed base.

   \item If the tag is
   \((\mathrm{Am}_{\mathcal N};\dot a^\eta_0,\ldots,
      \dot a^\eta_{k_\eta-1})\), respectively
   \((\mathrm{Am}_{\mathcal M};\dot a^\eta_0,\ldots,
      \dot a^\eta_{k_\eta-1})\), put
   \[
      \dot W_\eta=L[T_2,\dot a^\eta_0,\ldots,
         \dot a^\eta_{k_\eta-1}],
   \]
   again with all real names lifted to \(\widehat{\mathbb P}_{\eta+1}\).  The
   second-step forcing is the measure-amoeba, respectively category-amoeba,
   forcing computed in \(\dot W_\eta\).
\end{enumerate}
Finally we
set
\[
   \mathbb P_{\eta+1}=\widehat{\mathbb P}_{\eta+1}*\dot{\mathbb Q}_\eta.
\]
At limit stages we take the mixed-support limit: countable support on the
reservoir product coordinates and finite support on the c.c.c. second-step
coordinates.  Thus every successor stage carries a fresh reservoir coordinate,
but \(\Phi\) reads such a coordinate only when the second step is an
almost-disjoint coding forcing.  We call the resulting object a \emph{local
hybrid presentation over \(W\)}.

We shall compare names by means of canonical localized presentations.  Since
\(W=M_1[H]\), we use the fixed \(M_1\)-coding of \(W\)-objects by
\(Br\)-names.  Whenever an object of \(W\) is used as a parameter, it is
represented by its \(<_{M_1}\)-least \(Br\)-name.  Thus the collection of codes
for the names and supports below is ordered by the canonical
\(M_1\)-construction order.

\begin{definition}[Admissible supports and canonical lifts]
\label{def:admissible-supports-canonical-lifts}
Let \(\mathcal R_\beta\) be a local hybrid presentation over \(W\), with resolved
bookkeeping \(F\), and let \(A\subseteq\beta\).  We define, by induction on
\(\eta\leq\beta\), the following objects simultaneously:
\begin{enumerate}
   \item the assertion that \(A\cap\eta\) is admissible for
   \(\mathcal R_\eta\);
   \item the restricted hybrid presentation \(\mathcal R_{A\cap\eta}\) and its
   forcing \(\mathbb P_{A\cap\eta}\);
   \item the canonical complete embedding
   \[
      i_A^\eta:\mathbb P_{A\cap\eta}\longrightarrow\mathbb P_\eta;
   \]
   \item the canonical lift of names along this embedding.
\end{enumerate}
At \(\eta=0\), the restricted presentation is trivial and the embedding is the
identity.

At a limit \(\lambda\leq\beta\), \(A\cap\lambda\) is admissible if
\(A\cap\eta\) is admissible for every \(\eta<\lambda\) and the restricted
presentations are coherent.  In that case \(\mathcal R_{A\cap\lambda}\) is the
mixed-support limit of the earlier restricted presentations, and
\(i_A^\lambda\) is the coordinatewise direct limit of the embeddings
\(i_A^\eta\), \(\eta<\lambda\).

Suppose now that \(\eta<\beta\) and that \(A\cap\eta\) has already been declared
admissible.  If \(\eta\notin A\), then the restricted presentation skips the
\(\eta\)-th stage:
\[
   \mathbb P_{A\cap(\eta+1)}=\mathbb P_{A\cap\eta}.
\]
The embedding \(i_A^{\eta+1}\) extends \(i_A^\eta\) by placing the trivial
condition at the fresh reservoir coordinate \(\mathbb C_\eta\) and at the
second-step coordinate \(\dot{\mathbb Q}_\eta\).

If \(\eta\in A\), then the \(\eta\)-th stage must be readable over the earlier
restricted forcing.  We first keep the fresh reservoir coordinate and form
\[
   \widehat{\mathbb P}_{A,\eta+1}
      =\mathbb P_{A\cap\eta}\times\mathbb C_\eta,
\]
with the embedding induced by \(i_A^\eta\times\mathrm{id}_{\mathbb C_\eta}\).
The second-step forcing is then defined according to the resolved tag
\(F(\eta)\):
\begin{enumerate}
   \item For the tags \(\mathbf 1\) and \(\mathrm{Coh}\), the restricted stage
   uses the same trivial, respectively ordinary Cohen, second-step forcing.

   \item If \(F(\eta)=\operatorname{Code}(\dot u_\eta)\), then there must be a
   \(\mathbb P_{A\cap\eta}\)-name \(\dot u^A_\eta\) such that
   \[
      \mathbb P_\eta\Vdash
      \dot u_\eta=(\dot u^A_\eta)^{\uparrow\eta}.
   \]
   If no such name exists, then \(A\cap(\eta+1)\) is not admissible.  If such a
   name exists, we use the \(<_{M_1}\)-least one and put
   \[
      \dot{\mathbb Q}^A_\eta
      =\dot{\mathbb A}_D(\dot Y_{\dot g_\eta,\dot u^A_\eta})
   \]
   over \(\widehat{\mathbb P}_{A,\eta+1}\).

   \item If \(F(\eta)\) is a relative regularity tag
   \[
      (\mathrm{Rand};\dot a^\eta_0,\ldots,\dot a^\eta_{k_\eta-1}),
      \quad
      (\mathrm{Am}_{\mathcal N};\dot a^\eta_0,\ldots,
         \dot a^\eta_{k_\eta-1}),
      \quad\text{or}\quad
      (\mathrm{Am}_{\mathcal M};\dot a^\eta_0,\ldots,
         \dot a^\eta_{k_\eta-1}),
   \]
   then for every \(i<k_\eta\) there must be a
   \(\mathbb P_{A\cap\eta}\)-name \(\dot a^{\eta,A}_i\) such that
   \[
      \mathbb P_\eta\Vdash
      \dot a^\eta_i=(\dot a^{\eta,A}_i)^{\uparrow\eta}.
   \]
   If one of these names does not exist, then \(A\cap(\eta+1)\) is not
   admissible.  If they exist, use the \(<_{M_1}\)-least such tuple and let
   \[
      \dot W^A_\eta
      =L[T_2,
         \dot a^{\eta,A}_0,\ldots,
         \dot a^{\eta,A}_{k_\eta-1}].
   \]
   The restricted stage uses the same kind of regularity forcing computed in
   \(\dot W^A_\eta\): Random, measure-amoeba, or category-amoeba according to
   the tag.
\end{enumerate}
If the relevant requirement is satisfied, then
\[
   \mathbb P_{A\cap(\eta+1)}
   =\widehat{\mathbb P}_{A,\eta+1}*\dot{\mathbb Q}^A_\eta,
\]
and the embedding \(i_A^{\eta+1}\) is the natural two-step extension of
\(i_A^\eta\).

Once \(i_A^\eta\) has been defined, the canonical lift of a
\(\mathbb P_{A\cap\eta}\)-name \(\sigma\) to a \(\mathbb P_\eta\)-name is the
recursive name translation
\[
   \sigma^{\uparrow\eta}
   =
   \{(\tau^{\uparrow\eta},i_A^\eta(p))\mid (\tau,p)\in\sigma\}.
\]
When \(\eta=\beta\) and \(A\cap\beta=A\) is admissible, we simply write
\(\mathbb P_A\), \(i_A^\beta\), and \(\sigma^{\uparrow\beta}\).
\end{definition}

Thus, if \(G_\beta\subseteq\mathbb P_\beta\) is generic and
\(G_A=(i_A^\beta)^{-1}[G_\beta]\), then \(G_A\) is
\(\mathbb P_A\)-generic and
\[
   (\sigma^{\uparrow\beta})^{G_\beta}=\sigma^{G_A}
\]
for every \(\mathbb P_A\)-name \(\sigma\).  The embedding
\(i_A^\beta\) identifies \(\mathbb P_A\) with a complete subforcing of
\(\mathbb P_\beta\).  In what follows we suppress the embedding and regard
\(\mathbb P_A\) as this complete subforcing.

If \(\vec\tau\) is a finite tuple of \(\mathbb P_\beta\)-names for reals or
natural numbers, then an admissible support \(A\subseteq\beta\) is
\emph{admissible for \(\vec\tau\)} if every member of \(\vec\tau\) is the
canonical lift of a \(\mathbb P_A\)-name.  There may be many admissible supports
for \(\vec\tau\), and there need not be a unique inclusion-minimal one.  We
therefore set
\[
   \operatorname{supp}_{\mathrm{loc}}(\vec\tau,\mathbb P_\beta)
   :=
   \text{the }<_{M_1}\text{-least admissible support for }\vec\tau.
\]
The notation
\[
   \mathbb P_{\operatorname{supp}_{\mathrm{loc}}(\vec\tau,\mathbb P_\beta)}
\]
always means the complete subforcing determined by this least admissible
support.  The definition of \(\operatorname{supp}_{\mathrm{loc}}\) has no size
clause.  Countable localization is proved later for the presentations which
come from the allowable recursion.

\begin{definition}[Localized presentation]\label{def:localized-presentation}
Let \(\dot x\) be a \(\mathbb P_\beta\)-name for a real.  A
\emph{localized presentation} of \(\dot x\) is a pair \((A,\sigma)\) such that
\(A\subseteq\beta\) is an admissible support, \(\sigma\) is a
\(\mathbb P_A\)-name for a real, and
\[
   1_{\mathbb P_\beta}\Vdash
   \dot x=\sigma^{\uparrow\beta}.
\]
We order localized presentations by the canonical \(<_{M_1}\)-order of their
codes.  The \emph{canonical localized presentation} of \(\dot x\) is the
\(<_{M_1}\)-least localized presentation of \(\dot x\).
\end{definition}

If \(G_\beta\subseteq\mathbb P_\beta\) is generic and \((A_i,\sigma_i)\) is the
canonical localized presentation of \(\dot z_i\), then
\[
   z_i=\dot z_i^{G_\beta}=\sigma_i^{G_{A_i}},
\]
where \(G_{A_i}\) is the induced \(\mathbb P_{A_i}\)-generic filter.  For two
names \(\dot z_0,\dot z_1\) we write
\[
   \dot z_0\triangleleft_{\mathrm{loc}}\dot z_1
\]
if the canonical localized presentation of \(\dot z_0\) is \(<_{M_1}\)-below
the canonical localized presentation of \(\dot z_1\).  This is the comparison
used by the well-order coding stages below.


\begin{definition}[Regularity-stage supports and branch footprints]
\label{def:admissible-random-base}
Let \(\mathcal R_\beta\) be a local hybrid presentation over \(W=M_1[H]\), and
let \(\eta<\beta\) be a relative regularity stage.  Suppose that the resolved
tag at \(\eta\) is one of
\[
   (\mathrm{Rand};\dot a^\eta_0,\ldots,\dot a^\eta_{k_\eta-1}),
   \qquad
   (\mathrm{Am}_{\mathcal N};\dot a^\eta_0,\ldots,
      \dot a^\eta_{k_\eta-1}),
   \qquad
   (\mathrm{Am}_{\mathcal M};\dot a^\eta_0,\ldots,
      \dot a^\eta_{k_\eta-1}).
\]
The base used at this stage is, by definition,
\[
   \dot W_\eta=L[T_2,\dot a^\eta_0,\ldots,
      \dot a^\eta_{k_\eta-1}],
\]
with the evident meaning \(L[T_2]\) when \(k_\eta=0\).  The regularity forcing
at the stage is the Random algebra, measure-amoeba forcing, or category-amoeba
forcing computed in this displayed base, according to the tag.

A \emph{support witness} for the regularity stage is a tuple
\[
   (A_{\mathrm{rb}},
    \langle \dot a^{\mathrm{rb}}_i\mid i<k_\eta\rangle,
    B^0_{\mathrm{rb}},
    \dot B_{\mathrm{rb}})
\]
such that:
\begin{enumerate}
   \item \(A_{\mathrm{rb}}\subseteq\eta\) is a countable admissible support, and
   each \(\dot a^{\mathrm{rb}}_i\) is a \(\mathbb P_{A_{\mathrm{rb}}}\)-name for
   a real;

   \item for every \(i<k_\eta\),
   \[
      \mathbb P_\eta\Vdash
      \dot a^\eta_i=(\dot a^{\mathrm{rb}}_i)^{\uparrow\eta};
   \]

   \item every object of \(W=M_1[H]\) occurring in the codes of the names
   \(\dot a^{\mathrm{rb}}_i\) is represented by its \(<_{M_1}\)-least
   \(Br\)-name, and the union of the supports of these \(Br\)-names is
   contained in the countable set \(B^0_{\mathrm{rb}}\subseteq\omega_1\).
\end{enumerate}
Let
\[
   \dot W_{\mathrm{rb}}
   =L[T_2,
      \dot a^{\mathrm{rb}}_0,\ldots,
      \dot a^{\mathrm{rb}}_{k_\eta-1}]
\]
be the base read on the support \(A_{\mathrm{rb}}\).  Thus the canonical lift of
\(\dot W_{\mathrm{rb}}\) to stage \(\eta\) is forced to be the stage base
\(\dot W_\eta\).

Using the notation from Definition~\ref{def:selected-branch-coordinates}, let
\(E(A_{\mathrm{rb}})\) be the set of explicit coding coordinates in
\(A_{\mathrm{rb}}\).  If \(\nu\in E(A_{\mathrm{rb}})\), let \(\dot C_\nu\) be the
\(\mathbb P_{A_{\mathrm{rb}}}\)-name for the coding area of the reservoir
generic at \(\nu\), and let \(\dot u_\nu\) be the
\(\mathbb P_{A_{\mathrm{rb}}}\)-name coded at \(\nu\).  The branch footprint of
the witness is the \(\mathbb P_{A_{\mathrm{rb}}}\)-name
\[
   \dot B_{\mathrm{rb}}
   =
   \check B^0_{\mathrm{rb}}
   \cup
   \bigcup_{\nu\in E(A_{\mathrm{rb}})}
      \operatorname{Sel}(\dot C_\nu,\dot u_\nu).
\]
Thus \(\dot B_{\mathrm{rb}}\) is generated by countably many coding areas and a
countable set of branch coordinates.  It may have size \(\omega_1\).

Among all support witnesses for the stage, choose the \(<_{M_1}\)-least one.  We
call the associated quadruple
\[
   (A^{\mathrm{reg}}_\eta,
    \dot W^{\mathrm{reg}}_\eta,
    B^{0,\mathrm{reg}}_\eta,
    \dot B^{\mathrm{reg}}_\eta)
\]
the \emph{canonical support witness} for the regularity stage, where
\[
   \dot W^{\mathrm{reg}}_\eta
   =L[T_2,
      \dot a^{\mathrm{reg}}_0,\ldots,
      \dot a^{\mathrm{reg}}_{k_\eta-1}]
\]
is the base read on the least support witness.  The regularity forcing is always
computed in the relevant \(L[T_2,\vec a]\)-base.  It is not recomputed after
passing to a larger ambient extension.
\end{definition}

\begin{remark}[Regularity bases and omission]
Definition~\ref{def:admissible-random-base} records the support and branch-footprint data which keep the regularity stages usable
in omitting models.  If a stage used Random or amoeba forcing computed over the
whole branch extension \(W=M_1[H]\), then the stage would not in general be
interpretable over
\[
   M_1[H\upharpoonright(\omega_1\setminus\{\xi\})].
\]
In the definition above the base is instead generated from finitely many real
parameters on a countable admissible support:
\[
   L[T_2,a_0,\ldots,a_{k-1}].
\]
The tree \(T_2\) is fixed ground-model data and contributes no branch-product
coordinate.  The branch-product coordinates relevant to the base are precisely
those used by the \(W\)-parameters occurring in the names for the reals
\(a_i\), together with the coordinates generated by explicit coding stages in
the admissible support on which those names are read.  This is recorded by the
name \(\dot B^{\mathrm{reg}}_\eta\).

The branch footprint of such a base is not required to be countable; it is
generated by countably many coding areas and countably many ground branch
coordinates.  Its complement has size \(\omega_1\): above one bound, each block
meets at most one generating coding area, and from each pair
\(\{\omega\cdot\gamma+2n,\omega\cdot\gamma+2n+1\}\) at most one coordinate is
selected.  This is the support form used in the branch omission and
no-unwanted-code arguments.

Thus, for a real parameter \(a\) with a countable local presentation,
\(L[T_2,a]\) is a base of the allowed form, and similarly for any finite tuple
of real parameters read on a countable admissible support.  The bookkeeping
inserts relative Random, relative measure-amoeba, and relative category-amoeba
stages over such bases cofinally often.
\end{remark}


\subsection{Allowable and \(\alpha\)-allowable local coding iterations}
\label{subsec:allowable-local-coding-iterations}

We now formulate the forcing recursion using the localized-presentation
convention.  Fix, once and for all, a universal lightface sequence
\[
   \langle A_m\mid m<\omega\rangle
\]
of \(\Pi^1_4\) subsets of \((2^\omega)^2\) in the \(M_1\)-case.  Real parameters are handled later by the standard pairing argument, i.e. by applying the lightface list to relations in the paired variable \((a,x)\).  In
the higher \(M_n\)-case this sequence is replaced by the corresponding
universal lightface sequence of \(\Pi^1_{n+3}\) sets.  A bookkeeping function
is an element of \(M_1\).  Its values are \(M_1\)-descriptions which are decoded
relative to the current associated post-branch forcing \(\mathbb P_\beta\).
We suppress the decoding map and write, for example,
\[
   F(\beta)=(\mathrm{Rand};\dot a_0,\ldots,\dot a_{k-1}),
   \qquad
   F(\beta)=\dot w,
   \qquad
   F(\beta)=(\dot z_0,\dot z_1),
   \qquad
   F(\beta)=(\dot x,\dot z,\dot m),
\]
when the decoded value has the corresponding form.  The regularity tags may
also be
\[
   (\mathrm{Am}_{\mathcal N};\dot a_0,\ldots,\dot a_{k-1})
   \quad\text{or}\quad
   (\mathrm{Am}_{\mathcal M};\dot a_0,\ldots,\dot a_{k-1}).
\]
Here the \(\dot a_i\)'s are \(\mathbb P_\beta\)-names for reals, and the base
associated with the tag is, by definition,
\[
   L[T_2,\dot a_0,\ldots,\dot a_{k-1}].
\]
Such a regularity tag is used only when it has a canonical admissibility witness
in the sense of Definition~\ref{def:admissible-random-base}; equivalently, the
finitely many real parameters generating the base are read on a countable
admissible support.  In the last case, \(\dot m\) is a name for a natural number
coding the relevant set \(A_m\), \(\dot x\) is the section parameter, and
\(\dot z\) is the bookkeeping guess for a possible value in the \(x\)-section.
The name \(\dot z\) is disposable: at successor allowability stages it is used
only when it is not selected as the protected tentative value.

The successor stages use the same hybrid product/iteration form as above.  At
every successor stage one first adds a fresh, independent copy of the
\(M_1\)-Cohen coding-area forcing.  The bookkeeping value then determines the
second-step forcing over the product with this fresh reservoir coordinate.  If
the stage explicitly codes a real \(\dot u\), this second step is the
Jensen--Solovay almost-disjoint coding determined by \(\dot u\).  Over the
current intermediate model the coding quotient is canonically isomorphic to the
two-step forcing
\[
   \mathbb C_{M_1}*\dot{\mathbb A}_D(\dot Y_{\dot g,\dot u}),
\]
but in the global hybrid presentation the first factor is inserted as an
actual product factor:
\[
   (\text{current forcing}\times \mathbb C_\beta)
   *\dot{\mathbb A}_D(\dot Y_{\dot g_\beta,\dot u}).
\]
Here \(\mathbb C_\beta\) is the fresh reservoir coordinate for the stage and
\(\dot g_\beta\) is its canonical generic.  At relative Random, relative amoeba,
ordinary Cohen, and trivial stages the same reservoir coordinate is still
present, but no almost-disjoint code is written from it.  The coding areas
obtained from distinct reservoir coordinates are mutually almost disjoint by
the standard fusion argument for \(\mathbb C_{M_1}\).


\begin{definition}[\(0\)-allowable local hybrid presentations]
\label{def:zero-local-allowable}
Let \(F\in M_1\) be a bookkeeping function of length
\(\delta\leq\omega_2\).  A sequence
\[
   \mathcal R=\langle \mathcal R_\beta\mid \beta\leq\delta\rangle
\]
over \(W\) is a \emph{\(0\)-allowable local hybrid presentation relative to
\(F\)} if it is obtained by the following recursion.

First, \(\mathcal R_0\) is the trivial post-branch presentation.  Thus the associated quotient forcing over \(W\) is
\[
   \mathbb P_0=\mathbf 1,
\]
whereas the corresponding full forcing over \(M_1\) is simply \(Br\).
At limit stages, \(\mathcal R_\eta\) is the mixed-support limit of the preceding
hybrid construction, and \(\mathbb P_\eta\) is the associated post-branch
forcing.  Suppose that \(\mathcal R_\beta\) has been constructed, and write
\(\mathbb P_\beta\) for its associated post-branch forcing.
Choose a fresh reservoir coordinate \(\mathbb C_\beta\), let
\(\dot g_\beta\) be its canonical generic, and put
\[
   \widehat{\mathbb P}_{\beta+1}
      =
   \mathbb P_\beta\times\mathbb C_\beta.
\]
The second-step forcing \(\dot{\mathbb Q}_\beta\) over
\(\widehat{\mathbb P}_{\beta+1}\) is determined by the decoded value of
\(F(\beta)\).

\begin{enumerate}
   \item If \(F(\beta)\) is one of the relative regularity tags
   \[
      (\mathrm{Rand};\dot a^\beta_0,\ldots,
         \dot a^\beta_{k_\beta-1}),
      \qquad
      (\mathrm{Am}_{\mathcal N};\dot a^\beta_0,\ldots,
         \dot a^\beta_{k_\beta-1}),
      \qquad
      (\mathrm{Am}_{\mathcal M};\dot a^\beta_0,\ldots,
         \dot a^\beta_{k_\beta-1}),
   \]
   and this stage has a canonical admissibility witness in the sense of
   Definition~\ref{def:admissible-random-base}, let
   \[
      (A^{\mathrm{reg}}_\beta,
       \dot W^{\mathrm{reg}}_\beta,
       B^{0,\mathrm{reg}}_\beta,
       \dot B^{\mathrm{reg}}_\beta)
   \]
   be that witness.  Thus, for the \(<_{M_1}\)-least tuple of
   \(\mathbb P_{A^{\mathrm{reg}}_\beta}\)-names for reals
   \(\langle\dot a^{\beta,\mathrm{reg}}_i\mid i<k_\beta\rangle\) reading the
   displayed real names, we have
   \[
      \dot W^{\mathrm{reg}}_\beta
      =L[T_2,
          \dot a^{\beta,\mathrm{reg}}_0,
          \ldots,
          \dot a^{\beta,\mathrm{reg}}_{k_\beta-1}].
   \]
   Let \(\dot W^{\mathrm{reg},+}_\beta\) be the canonical
   \(\widehat{\mathbb P}_{\beta+1}\)-name obtained by lifting this base, equivalently
   the name for
   \[
      L[T_2,
        (\dot a^{\beta,\mathrm{reg}}_0)^{\uparrow},\ldots,
        (\dot a^{\beta,\mathrm{reg}}_{k_\beta-1})^{\uparrow}]
   \]
   over the stage \(\widehat{\mathbb P}_{\beta+1}\).  The second-step forcing
   \(\dot{\mathbb Q}_\beta\) is the member of the following list corresponding
   to the tag:
   \[
      \dot{\mathbb B}_{\mathrm{rand}}^{\dot W^{\mathrm{reg},+}_\beta},
      \qquad
      \dot{\mathbb A}_{\mathcal N}^{\dot W^{\mathrm{reg},+}_\beta},
      \qquad
      \dot{\mathbb A}_{\mathcal M}^{\dot W^{\mathrm{reg},+}_\beta}.
   \]
   Thus the relative Random or amoeba forcing is computed in the lifted base
   \(L[T_2,(\dot a^{\beta,\mathrm{reg}}_i)^{\uparrow}:i<k_\beta]\).  It is not
   recomputed in the larger ambient extension.  The local support
   \(A^{\mathrm{reg}}_\beta\) and the footprint name
   \(\dot B^{\mathrm{reg}}_\beta\) are recorded as data of this stage.  If the
   displayed regularity tag has no admissibility witness, put
   \(\dot{\mathbb Q}_\beta=\mathbf 1\).

   \item If \(F(\beta)\) is an ordinary Cohen tag, then
   \[
      \dot{\mathbb Q}_\beta=\dot{\mathbb C}_\omega.
   \]

   \item If \(F(\beta)=\dot w\), where \(\dot w\) is a
   \(\mathbb P_\beta\)-name for a real, let
   \[
      \dot u_\beta=\dot w
   \]
   and put
   \[
      \dot{\mathbb Q}_\beta
      =
      \dot{\mathbb A}_D(\dot Y_{\dot g_\beta,\dot u_\beta}).
   \]

   \item If
   \[
      F(\beta)=(\dot x,\dot z,\dot m),
   \]
   where \(\dot x,\dot z\) are \(\mathbb P_\beta\)-names for reals and
   \(\dot m\) is a \(\mathbb P_\beta\)-name for a natural number, let
   \[
      \dot u_\beta=\langle \dot x,\dot z,\dot m\rangle
   \]
   and put
   \[
      \dot{\mathbb Q}_\beta
      =
      \dot{\mathbb A}_D(\dot Y_{\dot g_\beta,\dot u_\beta}).
   \]

   \item If
   \[
      F(\beta)=(\dot z_0,\dot z_1),
   \]
   where \(\dot z_0,\dot z_1\) are \(\mathbb P_\beta\)-names for reals, let
   \((A_i,\sigma_i)\) be the canonical localized presentation of \(\dot z_i\),
   for \(i<2\).  If
   \[
      (A_0,\sigma_0)<_{M_1}(A_1,\sigma_1),
   \]
   let
   \[
      \dot u_\beta=\langle \dot z_0,\dot z_1\rangle
   \]
   and put
   \[
      \dot{\mathbb Q}_\beta
      =
      \dot{\mathbb A}_D(\dot Y_{\dot g_\beta,\dot u_\beta}).
   \]
   If instead
   \[
      (A_1,\sigma_1)<_{M_1}(A_0,\sigma_0),
   \]
   let
   \[
      \dot u_\beta=\langle \dot z_1,\dot z_0\rangle
   \]
   and put
   \[
      \dot{\mathbb Q}_\beta
      =
      \dot{\mathbb A}_D(\dot Y_{\dot g_\beta,\dot u_\beta}).
   \]
   If the two canonical localized presentations are equal, put
   \[
      \dot{\mathbb Q}_\beta=\mathbf 1.
   \]
\end{enumerate}
If none of the preceding cases applies, put
\[
   \dot{\mathbb Q}_\beta=\mathbf 1.
\]
Finally define
\[
   \mathbb P_{\beta+1}
      =
   \widehat{\mathbb P}_{\beta+1}*\dot{\mathbb Q}_\beta.
\]
In cases (3), (4), and the nontrivial subcases of (5), we say that stage
\(\beta\) \emph{explicitly codes} the name \(\dot u_\beta\).  A forcing
admitting such a presentation is called \emph{locally allowable}, or
\emph{\(0\)-allowable}.
\end{definition}





\begin{lemma}
\label{lem:basic-zero-allowable}
Let \(\mathcal R_\delta\) be a \(0\)-allowable local hybrid presentation over \(W\).
Then the following hold.
\begin{enumerate}
   \item \textbf{Normal form.}  Over the branch extension \(W\), the
   post-branch quotient forcing described by \(\mathcal R_\delta\) can be
   rearranged into an equivalent presentation of the form
   \[
      \mathbb C_{\mathrm{res}}*\dot{\mathbb S}.
   \]
   Here \(\mathbb C_{\mathrm{res}}\) is the countable-support product of all
   fresh \(M_1\)-Cohen reservoir coordinates introduced at successor stages of
   the hybrid presentation, including stages whose second-step forcing is
   Random, amoeba, ordinary Cohen, or trivial.  The forcing \(\dot{\mathbb S}\) is a
   finite-support iteration of c.c.c. real-adding forcings: ordinary Cohen
   forcing, fixed relative Random algebras, fixed relative amoeba forcings over
   their associated \(L[T_2,\vec a]\)-bases, Jensen--Solovay almost-disjoint coding forcings
   computed from the relevant reservoir generics and localized names, and
   trivial factors.  Equivalently, the full forcing \(Br\hybcomp\mathcal R_\delta\) over \(M_1\)
   is equivalent to
   \[
      Br*(\mathbb C_{\mathrm{res}}*\dot{\mathbb S}).
   \]

   \item The post-branch quotient \(\mathbb P_\delta\) is proper and preserves \(\omega_1\); equivalently, the full forcing \(Br\hybcomp\mathcal R_\delta\) preserves \(\omega_1\) over \(M_1\).
   Moreover, if \(\delta\leq\omega_2\), then the corresponding forcing preserves \(\omega_2\).

   \item Coding areas are almost disjoint: if \(g\) and \(h\) are distinct
   \(M_1\)-Cohen subsets of \(\omega_1\), then the corresponding coding areas
   \(C_g\) and \(C_h\) have bounded, hence countable, intersection.  In
   particular this holds for the reservoir generics added at distinct successor
   stages of the hybrid presentation.

   \item \textbf{Non-interference of coding blocks.}  A coding block attached to
   a reservoir coordinate is read only from that reservoir area and from the
   corresponding almost-disjoint coding real.  Later stages use fresh reservoir
   coordinates.  Hence later forcing cannot alter the decoding of an earlier
   block, and a decoded area which agrees with the reservoir area of an earlier
   explicit coding stage yields the real intentionally written at that stage.
   Reservoir coordinates whose second-step forcing is Random, amoeba, ordinary
   Cohen, or trivial write no almost-disjoint code and therefore create no coding
   block.

   \item If a stage explicitly codes a name \(\dot u_\beta\) and
   \(u_\beta=(\dot u_\beta)^{G_\beta}\), then in the next intermediate
   extension there is a real \(r_\beta\) such that \(\Psi(r_\beta,u_\beta)\)
   holds.  In particular \(\Phi(u_\beta)\) holds at that stage.  This instance
   of \(\Phi(u_\beta)\) persists to all later locally allowable extensions
   which preserve \(\omega_1\) and use fresh reservoir coordinates.
\end{enumerate}
\end{lemma}

\begin{proof}
For (1), observe that the first component of every successor stage is a fresh
copy of the \(M_1\)-computed \(\omega_1\)-Cohen reservoir forcing.  These
coordinates are independent of one another and are ordered with countable
support.  They may therefore be moved to a single product part
\(\mathbb C_{\mathrm{res}}\).  The second-step forcing at a stage is then read
over the product extension.  If the stage is a relative regularity stage, the
forcing remains the fixed forcing computed in the lifted admissible base; it is
not recomputed in the larger universe.  If the stage is an explicit coding
stage, the almost-disjoint forcing is computed from the lifted localized name
and from the corresponding reservoir generic.  Ordinary Cohen and trivial
stages are unaffected.  Canonical lifting of localized presentations preserves
the well-order comparison clauses.  This gives the stated normal form, and the
forcing over \(M_1\) is obtained by placing the initial branch product
\(Br\) in front.

For (2), use the normal form from (1).  The branch product \(Br\) is c.c.c. and
has size \(\aleph_1\).  The reservoir part \(\mathbb C_{\mathrm{res}}\) is the
\(M_1\)-computed countable-support product of copies of \(\mathbb C_{M_1}\).  We
use its closure before the branch product is performed: over \(M_1\) this
reservoir block is \(\sigma\)-closed, hence it preserves \(\omega_1\) and
preserves the finite-product Suslinity of the fixed sequence \(\vec S\).  Thus,
in the reservoir extension, \(Br\) is still c.c.c.  After the branch product, the
real-adding part \(\dot{\mathbb S}\) is a finite-support iteration of c.c.c.
forcings: ordinary Cohen forcing, fixed relative Random algebras, fixed relative
amoeba forcings over admissible bases, Jensen--Solovay almost-disjoint coding
forcings, and trivial factors.  Hence \(\dot{\mathbb S}\) is c.c.c. and therefore
proper.  This commuted presentation proves preservation for the full forcing
over \(M_1\), and therefore for the post-branch quotient over \(W\).  Notice that
we do not use the post-branch reservoir quotient as a literally \(\sigma\)-closed
forcing over \(W\).

If \(\delta\leq\omega_2\), the same normal form gives preservation of
\(\omega_2\).  Since \(M_1\) satisfies \(\GCH\) and \(Br\) has size \(\aleph_1\), the branch
extension \(W\) satisfies \(\CH\) and \( (\omega_2)^\omega=\omega_2\).
The usual \(\Delta\)-system argument therefore shows that the
countable-support product of at most \(\omega_2\) many \(\aleph_1\)-sized
reservoir coordinates has the \(\aleph_2\)-chain condition.
In the reservoir extension, every real-adding iterand in \(\dot{\mathbb S}\) has
size at most \(\aleph_1\): this is clear for ordinary Cohen forcing,
Jensen--Solovay almost-disjoint coding, and trivial forcing, and for relative
relative regularity forcing it follows because admissible bases have only
\(\aleph_1\)-many relevant Borel codes.  A second \(\Delta\)-system argument
therefore shows that the finite-support real-adding iteration has the
\(\aleph_2\)-chain condition.  Thus \(\omega_2\) is preserved.

Clause (3) follows from the definition of the coding areas.  If \(g\ne h\), let
\(\xi<\omega_1\) be the first coordinate at which they differ.  Then no proper
initial segment of \(g\) of length above \(\xi\) can equal a proper initial
segment of \(h\) of length above \(\xi\).  Since \(\rho\) is one-to-one,
\(C_g\cap C_h\) is contained in the codes of the initial segments of length at
most \(\xi\), and is therefore bounded in \(\omega_1\).

For (4), the decoding predicate is local to a single reservoir area and to the
almost-disjoint real added for the reshaped set associated with that area.
Distinct reservoir areas are almost disjoint by (3), and all later stages use
fresh reservoir coordinates.  Thus later coding blocks may add new codes, but
they do not change the old almost-disjoint pattern and cannot make a disjoint
block decode the same old object.  If a purported decoding uses the same
reservoir area as an earlier explicit coding stage, then the David decoding on
that area is the one produced at that stage, so it yields the intentionally
coded real.  If the reservoir area belongs to a non-coding stage, there is no
corresponding almost-disjoint coding real and hence no coding block to decode.

For (5), after \(g_\beta\) is added, the set \(Y_{g_\beta,u_\beta}\) is exactly
the David reshaping of the model code associated with \(u_\beta\), and the real
\(r_\beta\) added by \(\mathbb A_D(Y_{g_\beta,u_\beta})\) almost-disjointly
decodes this \(Y\) over every relevant countable model.  Hence
\(\Psi(r_\beta,u_\beta)\) holds by the definition of \(\Psi\).  The persistence
follows because later locally allowable stages preserve \(\omega_1\), use fresh
reservoir coordinates, and do not change the already added real \(r_\beta\) or
the initial segment of the almost-disjoint family used in the decoding.
\end{proof}


We next fix the branch-footprint notation used in the omission arguments.  This
notation is not part of the definition of admissibility.

Let \(A\subseteq\delta\) be a countable admissible support, let
\(\vec\sigma\) be a finite tuple of \(\mathbb P_A\)-names for reals, and let
\(G_\delta\subseteq\mathbb P_\delta\) be generic.  The names in
\(\vec\sigma\), the bookkeeping entries used on \(A\), and the names appearing
in the stages of \(\mathcal R_A\) may use parameters from \(W=M_1[H]\).  Each such
parameter is represented by its \(<_{M_1}\)-least \(Br\)-name.  Let
\[
   B_0(A,\vec\sigma)\subseteq\omega_1
\]
be the union of the supports of these \(Br\)-names.  If \(\eta\in A\) is a
relative Random, measure-amoeba, or category-amoeba stage, and
\[
   (A^{\mathrm{reg}}_\eta,
    \dot W^{\mathrm{reg}}_\eta,
    B^{0,\mathrm{reg}}_\eta,
    \dot B^{\mathrm{reg}}_\eta)
\]
is the canonical admissibility witness for its base, then we also put
\(B^{0,\mathrm{reg}}_\eta\) into \(B_0(A,\vec\sigma)\).  Thus
\(B_0(A,\vec\sigma)\) is countable.

Let \(R(A)\) be the set of relative Random and relative amoeba stages in
\(A\), and let \(E(A)\) be the set of explicit coding stages in \(A\).  If
\(\eta\in R(A)\), set
\[
   B^{\mathrm{reg},G}_\eta
   =
   (\dot B^{\mathrm{reg}}_\eta)^{G_{A^{\mathrm{reg}}_\eta}}.
\]
If \(\eta\in E(A)\), let \(C^G_\eta\) be the coding area determined by the
reservoir generic at \(\eta\), and let \(u^G_\eta\) be the real coded at
\(\eta\).  Define
\[
\begin{split}
   U(A,\vec\sigma,G_\delta)=
   B_0(A,\vec\sigma)
   &\cup
   \bigcup_{\eta\in R(A)}B^{\mathrm{reg},G}_\eta \\
   &\cup
   \bigcup_{\eta\in E(A)}
      \operatorname{Sel}(C^G_\eta,u^G_\eta).
\end{split}
\]
Thus \(U(A,\vec\sigma,G_\delta)\), informally the branch footprint of
\(A\) and \(\vec\sigma\) in the generic extension, records the branch
coordinates used by the restricted presentation and the displayed names.  It is
generated by countably many coding areas and a countable set of branch-product
coordinates.  It may have size \(\omega_1\).  At a coding
stage only the selected coordinates
\(\operatorname{Sel}(C^G_\eta,u^G_\eta)\) are included; the other member of
\(\{\omega\cdot\gamma+2n,\omega\cdot\gamma+2n+1\}\) is not included
unless it is selected by another coding stage or belongs to a regularity-base
footprint.


\begin{lemma}[Countable localization and branch footprints]
\label{lem:countable-localization-real-names}
Let \(\mathcal R_\delta\) be a \(0\)-allowable local hybrid presentation over
\(W\), and let \(\mathbb P_\delta\) be its associated post-branch forcing.  If
\(\vec\tau\) is a finite tuple of \(\mathbb P_\delta\)-names for reals or
natural numbers, then there are a countable admissible support
\(A\subseteq\delta\) and a tuple \(\vec\sigma\) of \(\mathbb P_A\)-names such
that
\[
   \mathbb P_\delta\Vdash
   \vec\tau=\vec\sigma^{\uparrow\delta}.
\]
Consequently, if \(G_\delta\subseteq\mathbb P_\delta\) is generic and
\(x\in W[G_\delta]\cap2^\omega\), then
\[
   x\in W[G_A]
\]
for some countable admissible \(A\subseteq\delta\), where
\(G_A=G_\delta\cap\mathbb P_A\).

Moreover, for every such \(A\), every finite tuple \(\vec\sigma\), and every
\(G_\delta\), the branch footprint has the form
\[
   U(A,\vec\sigma,G_\delta)
   =B\cup\bigcup_{j\in J}\operatorname{Sel}(C_j,u_j),
\]
where \(J\) is countable, \(B\in[\omega_1]^\omega\), each \(C_j\) is a reservoir
coding area occurring in the local presentation, and each \(u_j\in2^\omega\).
Thus the conclusions of Lemma~\ref{lem:countably-many-coding-areas} apply to
\(U(A,\vec\sigma,G_\delta)\).
\end{lemma}

\begin{proof}
We prove the localization statement by induction on \(\delta\).  We use the
normal form from Lemma~\ref{lem:basic-zero-allowable}, but in the commuted full
presentation over \(M_1\).  Thus the full forcing is equivalent to
\[
   \mathbb C_{\mathrm{res}}*(Br*\dot{\mathbb S}),
\]
where \(\mathbb C_{\mathrm{res}}\) is the \(M_1\)-computed countable-support
product of the reservoir coordinates, and, in the \(\mathbb C_{\mathrm{res}}\)-extension,
\(Br*\dot{\mathbb S}\) is a finite-support c.c.c. iteration of the branch
product together with the real-adding second-step coordinates.  The reservoir
block is used here before the branch product.  It is \(\sigma\)-closed over
\(M_1\), and hence adds no reals, and no countable sequences of
\(M_1\)-ordinals.  We do not use the post-branch reservoir quotient over
\(W=M_1[H]\) as a closed forcing.

Replace \(\vec\tau\) by equivalent names in this commuted presentation.  In the
\(\mathbb C_{\mathrm{res}}\)-extension, take nice names for these reals and
natural numbers with respect to the c.c.c. forcing \(Br*\dot{\mathbb S}\).  Thus,
for each bit of each real name, and for each natural-number name, there is a
countable maximal antichain deciding the relevant value.  Each condition in
such an antichain has finite branch support and finite support in the
real-adding second-step coordinates.  Since \(\mathbb C_{\mathrm{res}}\) adds no
countable sequences of ground ordinals, the union of the second-step supports
appearing in all these antichains is represented by a countable set of hybrid
coordinates.  Let this set be \(A_0\subseteq\delta\).  The branch coordinates
which occur in the \(W\)-parameters of these names are recorded separately in
\(B_0\); they are not coordinates of the post-branch forcing.

We now close \(A_0\) under the dependencies needed to make the restricted
hybrid presentation meaningful.  Given \(A_n\), define \(A_{n+1}\) as follows.
If \(\eta\in A_n\) is a relative Random, measure-amoeba, or category-amoeba
stage, let
\[
   (A^{\mathrm{reg}}_\eta,
    \dot W^{\mathrm{reg}}_\eta,
    B^{0,\mathrm{reg}}_\eta,
    \dot B^{\mathrm{reg}}_\eta)
\]
be the canonical support witness for the regularity base at \(\eta\).  Thus
\[
   \dot W^{\mathrm{reg}}_\eta
      =L[T_2,\dot a^\eta_0,
             \ldots,\dot a^\eta_{k_\eta-1}]
\]
for finitely many names for reals read on
\(A^{\mathrm{reg}}_\eta\).  Put \(A^{\mathrm{reg}}_\eta\subseteq A_{n+1}\).
If \(\eta\in A_n\) is an explicit coding stage, apply the induction hypothesis
below \(\eta\) to the name which is coded at \(\eta\), and add the resulting
countable admissible support to \(A_{n+1}\).  Ordinary Cohen and trivial stages
require no additional support.  Finally let
\[
   A=\bigcup_{n<\omega}A_n.
\]
Then \(A\) is countable.

By construction, every nontrivial stage kept in \(A\) has all earlier data
needed to interpret its tag: regularity stages have their finite real-parameter
base read on the included support, and coding stages have the coded real read
on the included support.  Hence \(A\) is admissible in the sense of
Definition~\ref{def:admissible-supports-canonical-lifts}, and the canonical
embedding identifies \(\mathbb P_A\) with a complete subforcing of
\(\mathbb P_\delta\).  The nice names chosen above mention only stages in
\(A\), after replacing every kept stage by its restricted version.  Therefore
they define a tuple \(\vec\sigma\) of \(\mathbb P_A\)-names such that
\[
   \mathbb P_\delta\Vdash
   \vec\tau=\vec\sigma^{\uparrow\delta}.
\]
This proves the localization statement.  The assertion for a real
\(x\in W[G_\delta]\) follows by applying the first part to a name for \(x\).

It remains to record the footprint form.  By definition,
\(B_0(A,\vec\sigma)\) is countable.  Each explicit coding stage in \(A\)
contributes one set \(\operatorname{Sel}(C,u)\).  If \(\eta\in A\) is a
relative Random, measure-amoeba, or category-amoeba stage, then the canonical
support witness for its base contributes
\[
   B^{0,\mathrm{reg}}_\eta
   \cup
   \bigcup_{\nu\in E(A^{\mathrm{reg}}_\eta)}
      \operatorname{Sel}(C_\nu,u_\nu).
\]
Here \(A^{\mathrm{reg}}_\eta\) is the countable support on which the finitely
many real parameters generating the displayed \(L[T_2,\vec a]\)-base are read,
and \(B^{0,\mathrm{reg}}_\eta\) is the countable branch-product support of the
\(W\)-parameters occurring in those names.  Since \(A\) is countable, all
selected coding-area contributions can be enumerated as
\[
   (C_j,u_j)_{j\in J}
   \qquad\text{with }J\text{ countable},
\]
together with one countable set \(B\subseteq\omega_1\).  Hence
\[
   U(A,\vec\sigma,G_\delta)
   =B\cup\bigcup_{j\in J}\operatorname{Sel}(C_j,u_j).
\]
Lemma~\ref{lem:countably-many-coding-areas} applies to every set of this form.
\end{proof}

\begin{lemma}
\label{lem:omit-unused-branch-coordinate}
Let \(\mathcal R_\delta\) be a \(0\)-allowable local hybrid presentation over
\(W\), let \(G_\delta\subseteq\mathbb P_\delta\) be generic, and let
\(A\subseteq\delta\) be a countable admissible support.  Let \(\vec\sigma\) be a
finite tuple of \(\mathbb P_A\)-names for reals.  Suppose that
\(\xi<\omega_1\) satisfies
\[
   \xi\notin U(A,\vec\sigma,G_\delta).
\]
Let \(H\subseteq Br\) be the branch-product generic used to form \(W\), and set
\[
   W_\xi=M_1[H\upharpoonright(\omega_1\setminus\{\xi\})].
\]
Then \(\mathcal R_A\) is interpretable over \(W_\xi\).  If
\(G_A=G_\delta\cap\mathbb P_A\), then \(G_A\) is generic over \(W_\xi\), the
interpretations of the names in \(\vec\sigma\) belong to \(W_\xi[G_A]\), and
\[
   W_\xi[G_A]\models ``S_\xi\text{ is Suslin}.''
\]
\end{lemma}

\begin{proof}
We argue by induction over the restricted presentation \(\mathcal R_A\).  Since
\(\xi\notin B_0(A,\vec\sigma)\), every ground-model parameter from
\(W=M_1[H]\) used by \(\vec\sigma\), by the bookkeeping on \(A\), or by a
stage of \(\mathcal R_A\), is represented by a \(<_{M_1}\)-least \(Br\)-name whose
support omits \(\xi\).  Hence all such parameters are in \(W_\xi\).
Moreover the product on the branch coordinates different from \(\xi\) preserves
\(S_\xi\), by independence of \(\vec S\).

Suppose first that \(\eta\in A\) is a relative Random, measure-amoeba, or
category-amoeba stage.  Let
\[
   (A^{\mathrm{reg}}_\eta,
    \dot W^{\mathrm{reg}}_\eta,
    B^{0,\mathrm{reg}}_\eta,
    \dot B^{\mathrm{reg}}_\eta)
\]
be the admissibility witness for its base.  Since \(A\) is admissible,
\[
   A^{\mathrm{reg}}_\eta\subseteq A\cap\eta.
\]
Also
\[
   B^{0,\mathrm{reg}}_\eta\subseteq B_0(A,\vec\sigma)
   \quad\text{and}\quad
   (\dot B^{\mathrm{reg}}_\eta)^{G_{A^{\mathrm{reg}}_\eta}}
      \subseteq U(A,\vec\sigma,G_\delta).
\]
Thus \(\xi\) is outside the branch footprint of the base.  By
Definition~\ref{def:admissible-random-base}, this witness is obtained from
finitely many \(\mathbb P_{A^{\mathrm{reg}}_\eta}\)-names for reals
\[
   \dot a^\eta_0,\ldots,\dot a^\eta_{k_\eta-1}
\]
and
\[
   \dot W^{\mathrm{reg}}_\eta
   =L[T_2,\dot a^\eta_0,\ldots,\dot a^\eta_{k_\eta-1}].
\]
The set \(B^{0,\mathrm{reg}}_\eta\) records the branch-product supports of the
\(W\)-parameters occurring in the codes of these real names, and
\((\dot B^{\mathrm{reg}}_\eta)^{G_{A^{\mathrm{reg}}_\eta}}\) records the branch
coordinates generated by the explicit coding stages in
\(A^{\mathrm{reg}}_\eta\).  Since both are contained in
\(U(A,\vec\sigma,G_\delta)\) and \(\xi\notin U(A,\vec\sigma,G_\delta)\), the
finite tuple
\[
   \langle \dot a^\eta_i\mid i<k_\eta\rangle
\]
is interpreted in \(W_\xi[G_{A^{\mathrm{reg}}_\eta}]\) with the same value as
in \(W[G_{A^{\mathrm{reg}}_\eta}]\).  Hence the base interpreted in the
\(\xi\)-omitting model is the same model
\[
   L[T_2,a^\eta_0,\ldots,a^\eta_{k_\eta-1}]
\]
as in the full extension; here \(T_2\in M_1\) is fixed ground-model data.  The
iterand at \(\eta\) is therefore the same relative Random, measure-amoeba, or
category-amoeba forcing computed in this base; it is not recomputed in the
ambient model.  The relative Random algebra is Knaster, and the two amoeba
forcings used here are \(\sigma\)-linked.  Hence the iterand preserves
\(S_\xi\).

If \(\eta\in A\) is an ordinary Cohen or trivial stage, the conclusion is
immediate.  The reservoir coordinate at \(\eta\) is the fixed
\(M_1\)-Cohen forcing and is independent of the branch coordinate \(\xi\); in
the normal form of Lemma~\ref{lem:basic-zero-allowable} it is placed before the
branch product and preserves \(S_\xi\).

It remains to consider an explicit coding stage \(\eta\in A\).  Let
\(\dot u_\eta\) be the name coded at \(\eta\), and let \(C_\eta^G\) be the
coding area of the reservoir generic.  By admissibility, \(\dot u_\eta\) is
read on an earlier support contained in \(A\).  The branch coordinates used by
the David code at this stage are precisely
\[
   \operatorname{Sel}(C_\eta^G,u_\eta^G).
\]
Since this set is contained in \(U(A,\vec\sigma,G_\delta)\) and
\(\xi\notin U(A,\vec\sigma,G_\delta)\), the branch \(b_\xi\) is not used in
forming \(Y_{g_\eta,u_\eta}\).  Hence the same reshaped set and the same
almost-disjoint coding forcing are computed over the \(\xi\)-omitting model.
The Jensen--Solovay almost-disjoint forcing is Knaster, and so preserves
\(S_\xi\).

At limit stages the restricted presentation uses mixed support.  The support in
\(A\) is countable on reservoir coordinates and finite on the c.c.c. real-adding
coordinates at each condition.  The preceding successor analysis therefore
iterates along \(A\).  Thus \(\mathcal R_A\) is defined over \(W_\xi\),
\(G_A\) is generic over \(W_\xi\), the tuple \(\vec\sigma^{G_A}\) belongs to
\(W_\xi[G_A]\), and no stage of the restricted forcing adds a branch through
\(S_\xi\).  Hence \(S_\xi\) remains Suslin in \(W_\xi[G_A]\).
\end{proof}

\begin{lemma}[No unwanted \(\Phi\)-codes]\label{lem:no-unwanted-codes}
Let \(\mathcal R_\delta\) be a \(0\)-allowable local hybrid presentation over \(W\),
and let \(G_\delta\subseteq\mathbb P_\delta\) be generic.  Suppose
that \(w\in W[G_\delta]\cap 2^\omega\) and
\[
   W[G_\delta]\models \Phi(w).
\]
Then there is an explicit coding stage \(\beta<\delta\) with coded name
\(\dot u_\beta\) such that
\[
   w=(\dot u_\beta)^{G_\beta}.
\]
Equivalently, Random stages, amoeba stages, ordinary Cohen stages, trivial
stages, and coding stages for other reals do not create new instances of
\(\Phi\).
\end{lemma}

\begin{proof}
Fix \(r\in W[G_\delta]\) such that \(\Psi(r,w)\).  Put
\[
   m_\infty=\mathcal J^{M_1}_{\omega_1}.
\]
By Lemma~\ref{lem:inner-model-decoded-by-Phi}, the model
\(L[r,m_\infty]\) decodes a coding area \(C_r\subseteq\omega_1\) and a selected
branch pattern.  Thus \(|C_r|=\omega_1\), and for every \(\gamma\in C_r\) and
\(n<\omega\), if
\[
   \xi=\tau_w(\gamma,n),
\]
then \(L[r,m_\infty]\) contains an \(\omega_1\)-branch through \(S_\xi\).

First suppose that \(r\in W\).  Since \(W=M_1[H]\) is obtained by the
finite-support branch product, there is \(B\in[\omega_1]^\omega\) such that
\(r\in M_1[H\upharpoonright B]\).  Since \(w\in L[r,m_\infty]\), also
\(w\in M_1[H\upharpoonright B]\).  Choose \(\gamma\in C_r\) such that
\[
   \{\omega\cdot\gamma+k\mid k<\omega\}\cap B=\emptyset.
\]
Choose \(n<\omega\), and set \(\xi=\tau_w(\gamma,n)\).  Then
\(\xi\notin B\).  By independence of \(\vec S\), the tree \(S_\xi\) is Suslin in
\(M_1[H\upharpoonright B]\).  But
\[
   L[r,m_\infty]\subseteq M_1[H\upharpoonright B]
\]
contains an \(\omega_1\)-branch through \(S_\xi\), a contradiction.  Hence
\(r\notin W\).

Choose a countable admissible support \(A\subseteq\delta\) and a
\(\mathbb P_A\)-name \(\sigma\) such that
\[
   \sigma^{G_A}=r,
   \qquad G_A=G_\delta\cap\mathbb P_A.
\]
Enlarge \(A\), if necessary, so that it contains the local supports needed for
this name.  Let \(\vec\sigma=(\sigma)\).  By the definition of the footprint,
there are a countable set \(J\), a countable set \(B\subseteq\omega_1\), coding
areas \((C_j:j\in J)\), and reals \((u_j:j\in J)\), such that
\[
   U(A,\vec\sigma,G_\delta)
   =
   B\cup\bigcup_{j\in J}\operatorname{Sel}(C_j,u_j).
\]
Each pair \((C_j,u_j)\) comes from an explicit coding stage of
\(\mathcal R_\delta\).  The set \(B\) contains the branch-product supports of the
\(W\)-parameters in the names under consideration and the countable parts of the
regularity-base witnesses.  The coding areas \((C_j:j\in J)\) are pairwise
almost disjoint.

If there is \(j\in J\) such that \(u_j=w\), then \(w\) is the real explicitly
coded at the corresponding stage, and we are done.  We therefore assume
\[
   \forall j\in J\quad u_j\ne w.
\]
Let
\[
   J_1=\{j\in J\mid C_r\cap C_j\text{ is unbounded in }\omega_1\}.
\]

If \(J_1\ne\emptyset\), choose \(j_0\in J_1\).  Then \(u_{j_0}\ne w\).  Choose
\(n<\omega\) with
\[
   u_{j_0}(n)\ne w(n).
\]
By Lemma~\ref{lem:countably-many-coding-areas}, applied to the countable family
\((C_j:j\in J)\), the set \(B\), the set \(C_r\), and the index \(j_0\), there
is \(\gamma\in C_r\cap C_{j_0}\) such that
\[
   \tau_w(\gamma,n)\notin U(A,\vec\sigma,G_\delta).
\]
Set \(\xi=\tau_w(\gamma,n)\).

If \(J_1=\emptyset\), then every \(C_r\cap C_j\), \(j\in J\), is bounded.  By
Lemma~\ref{lem:countably-many-coding-areas} there is \(\gamma\in C_r\) such that
\[
   \{\omega\cdot\gamma+k\mid k<\omega\}
   \cap U(A,\vec\sigma,G_\delta)=\emptyset.
\]
Choose \(n<\omega\), and set \(\xi=\tau_w(\gamma,n)\).

In both cases
\[
   \xi\notin U(A,\vec\sigma,G_\delta),
\]
and \(L[r,m_\infty]\) contains an \(\omega_1\)-branch through \(S_\xi\).  By
Lemma~\ref{lem:omit-unused-branch-coordinate}, the restricted extension
\(W_\xi[G_A]\) contains \(r\), and \(S_\xi\) remains Suslin in
\(W_\xi[G_A]\).  Since \(m_\infty\in M_1\subseteq W_\xi[G_A]\), we have
\[
   L[r,m_\infty]\subseteq W_\xi[G_A].
\]
This contradicts the Suslinity of \(S_\xi\) in \(W_\xi[G_A]\).  The contradiction
shows that the earlier case must occur: \(w\) is explicitly coded at some stage
\(\beta<\delta\).
\end{proof}


We now define the derivative hierarchy of allowable forcings.  The point of
the hierarchy is to separate genuine candidates for a uniformizing value from
guesses which can be destroyed by further allowable forcing.  A successor step
asks the following local question.  Given the bookkeeping data \((x,A_m)\) in
the current extension, is there already, in the local subextension determined
by the canonical support of these data, a real \(y\) such that \((x,y)\) is in
the \(x\)-section of \(A_m\) and such that lower-rank allowable tails cannot
force it out of that section?  If the answer is yes, \(y\) becomes a tentative
value of the least possible rank, and the construction codes some other real
on the same section.  If the answer is no, the construction treats the
bookkeeping value as a guess and codes that guess.  The hierarchy is defined
by transfinite recursion.  Limit levels are pure intersections; successor
levels are obtained by the local test just described.

For \(\rho\)-allowability we work with pairs
\[
   (\mathcal R,\dot I)\in W,
\]
where \(\mathcal R\) is a mixed-support local hybrid presentation over \(W\) and
\(\dot I\) is a name, over the terminal post-branch forcing associated with
\(\mathcal R\), for a set of quadruples \((x,y,m,\xi)\).  After interpretation by a generic filter, \(I\) is the set
of tentative uniformizing values, and the ordinal \(\xi\) records the rank at
which the value was selected.  The base level is
\[
   0\text{-allowable}=\text{locally allowable},
\]
with \(\dot I\) the canonical name for the empty set.

\begin{definition}[Limit \(\lambda\)-allowability]
\label{def:limit-lambda-allowable-local}
Assume that \(\lambda\) is a nonzero limit ordinal and that
\(\zeta\)-allowability has already been defined for every \(\zeta<\lambda\).
Let \(\mathcal A^W_\zeta\) denote the class of pairs
\((\mathcal R,\dot I)\in W\) which are \(\zeta\)-allowable over \(W\).  We define
\[
   \mathcal A^W_\lambda
   =
   \bigcap_{\zeta<\lambda}\mathcal A^W_\zeta.
\]
Thus \((\mathcal R,\dot I)\) is \emph{\(\lambda\)-allowable} iff it is
\(\zeta\)-allowable for every \(\zeta<\lambda\).  
\end{definition}

\begin{definition}[Relative \(\rho\)-allowable extensions]
\label{def:relative-rho-allowable-local}
Assume that \(\rho\)-allowability has been defined.  If
\((\mathcal R,\dot I)\) and \((\mathcal S,\dot J)\) are pairs belonging to
\(W\), we write
\[
   (\mathcal S,\dot J)\triangleright_\rho(\mathcal R,\dot I)
\]
if \(\mathcal S\) is obtained by appending to \(\mathcal R\) a mixed-support local
hybrid tail which satisfies the same \(\rho\)-allowability clauses, starting
with the auxiliary name \(\dot I\) instead of the empty name.  Let \(P\) and
\(P'\) be the terminal post-branch forcings associated with \(\mathcal R\) and
\(\mathcal S\), respectively.  If \(G\subseteq P\) is generic, then \(P'/G\) is
called a \emph{\(\rho\)-allowable tail} over the interpreted pair
\((\mathcal R,\dot I)^G\). 
\end{definition}

\begin{definition}[Successor \(\rho\)-allowability]
\label{def:rho-allowable-local}
Assume that \(\rho=\alpha+1\) and that \(\zeta\)-allowability and relative
\(\zeta\)-allowability have been defined for every \(\zeta<\rho\).  Let
\(F\in M_1\) be a bookkeeping function of length \(\delta\leq\omega_2\).  We
say that \((\mathcal R_\delta,\dot I_\delta)\) is \emph{\(\rho\)-allowable
relative to \(F\)} if it is obtained by the following recursion in \(W\).

At each stage \(\beta\leq\delta\) we construct a hybrid presentation
\(\mathcal R_\beta\) and a \(\mathbb P_\beta\)-name \(\dot I_\beta\), where
\(\mathbb P_\beta\) denotes the post-branch forcing associated with
\(\mathcal R_\beta\).
We start with $W$ and with \(\dot I_0\) the
canonical name for the empty set.  At a limit stage \(\eta\leq\delta\),
\(\mathcal R_\eta\) is the mixed-support limit of the preceding hybrid
presentations, and \(\dot I_\eta\) is the canonical \(\mathbb P_\eta\)-name
forced to be the union of the earlier interpreted auxiliary names.

Suppose that \(\beta<\delta\) and that \((\mathcal R_\beta,\dot I_\beta)\) has
already been constructed.  The value \(F(\beta)\) is decoded relative to
\(\mathbb P_\beta\).  If the decoded value is a relative Random tag, a relative amoeba tag, a Cohen
tag, a direct coding tag \(\dot w\), or a well-order tag
\((\dot z_0,\dot z_1)\), then the next stage is the hybrid operation
prescribed in Definition~\ref{def:zero-local-allowable}, and
\(\dot I_{\beta+1}\) is the canonical lift of \(\dot I_\beta\).

It remains to describe the uniformization tag.  Suppose that
\[
   F(\beta)=(\dot x,\dot z,\dot m),
\]
where \(\dot x,\dot z\) are \(\mathbb P_\beta\)-names for reals and \(\dot m\)
is a \(\mathbb P_\beta\)-name for a natural number.  Let
\[
   A=\operatorname{supp}_{\mathrm{loc}}((\dot x,\dot m),\mathbb P_\beta).
\]
A candidate value in the test below is a \(\mathbb P_A\)-name \(\dot y_0\) for
a real, lifted canonically to a \(\mathbb P_\beta\)-name.  We distinguish two
alternatives.

\begin{enumerate}
\item[(a)] There are an ordinal \(\zeta<\rho\) and a \(\mathbb P_A\)-name
\(\dot y_0\) for a real such that
\[
   \mathbb P_\beta\Vdash
   (\dot x,\dot y_0^{\uparrow\beta})\in A_{\dot m},
\]
and no relative \(\zeta\)-allowable extension in \(W\) can force this pair
out of the relevant section.  More explicitly, for every relative
\(\zeta\)-allowable extension
\[
   (\mathcal S,\dot J)\triangleright_\zeta(\mathcal R_\beta,\dot I_\beta)
\]
in \(W\), write \(P_{\mathcal S}\) for the terminal post-branch forcing
associated with \(\mathcal S\).  If the displayed names are lifted to
\(P_{\mathcal S}\), then
\[
   \mathbb P_\beta\Vdash
   \text{``there is no condition in }
   P_{\mathcal S}/\dot G_\beta
   \text{ which forces }
   (\dot x^{\uparrow\mathcal S},\dot y_0^{\uparrow\mathcal S})
      \notin A_{\dot m^{\uparrow\mathcal S}}\text{.''}
\]
If such witnesses exist, choose first the least possible \(\zeta\), and then
choose the \(<_{M_1}\)-least localized presentation of a witnessing name
\(\dot y_0\).

In this case \(\dot y_0\) is declared to be a tentative value of rank
\(\zeta\).  The next forcing does not code this tentative value.  Instead it
codes another real on the same section: if the bookkeeping guess \(\dot z\) is
forced to be different from \(\dot y_0^{\uparrow\beta}\), use \(\dot z\);
otherwise use the \(<_{M_1}\)-least localized name \(\dot z'\) for a real
forced to be different from \(\dot y_0^{\uparrow\beta}\).  Denote the resulting
name by \(\dot u\).  The next coding operation is
\[
   \operatorname{Code}(\langle\dot x,\dot u,\dot m\rangle),
\]
understood in the hybrid sense of Definition~\ref{def:zero-local-allowable}.
The auxiliary name is updated by letting \(\dot I_{\beta+1}\) be the canonical
\(\mathbb P_{\beta+1}\)-name forced to be
\[
   \dot I_\beta^{\dot G_\beta}
   \cup
   \{(\dot x^{\dot G_\beta},
       (\dot y_0^{\uparrow\beta})^{\dot G_\beta},
       \dot m^{\dot G_\beta},\check\zeta)\}.
\]

\item[(b)] If clause~(a) fails, then the construction guesses.  The bookkeeping
guess \(\dot z\) is used, unless the bookkeeping code does not provide a usable
real name, in which case we use the \(<_{M_1}\)-least default real name.  The
next coding operation is
\[
   \operatorname{Code}(\langle\dot x,\dot z,\dot m\rangle),
\]
again understood as the corresponding hybrid product/iteration step, and no
new tentative value is added:
\[
   \dot I_{\beta+1}=\dot I_\beta
\]
up to the canonical lift to the longer hybrid presentation.
\end{enumerate}

If the decoded bookkeeping value is none of the recognized tags, the stage is
trivial and \(\dot I_{\beta+1}\) is the canonical lift of \(\dot I_\beta\).  If
there is a bookkeeping function \(F\in M_1\) such that \((\mathcal R,\dot I)\) is
\(\rho\)-allowable relative to \(F\), then we simply say that
\((\mathcal R,\dot I)\) is \(\rho\)-allowable.
\end{definition}

\begin{lemma}[Shrinking of allowability]
\label{lem:shrinking-allowability-local}
Work in \(W\).  Let \(\alpha<\beta\).  Let \(\mathcal R\) be a hybrid presentation,
and let \(\mathbb P\) be its terminal post-branch forcing.  Suppose that
\(\mathcal R\) is \(\beta\)-allowable, say there is a
\(\mathbb P\)-name \(\dot I^\beta\) such that \((\mathcal R,\dot I^\beta)\) is
\(\beta\)-allowable.  Then there is a \(\mathbb P\)-name \(\dot I^\alpha\) such
that \((\mathcal R,\dot I^\alpha)\) is \(\alpha\)-allowable.  In particular, the
classes of underlying allowable hybrid presentations are decreasing with the
rank of allowability.
\end{lemma}

\begin{proof}
We prove the statement for the underlying hybrid presentation.  The auxiliary
name is allowed to change, and in general it should change, because the set of
tentative values records the ranks at which values were selected.

If \(\beta\) is a limit ordinal, the assertion is immediate from
Definition~\ref{def:limit-lambda-allowable-local}.  If \(\alpha\) is a nonzero
limit, it is enough to prove the assertion for all \(\xi<\alpha\).  Thus it
remains to consider the case in which \(\beta\) is a successor ordinal and
\(\alpha\) is either \(0\) or a successor ordinal.

Write \(\beta=\theta+1\), and let \(F\in M_1\) witness that
\((\mathcal R_\delta,\dot I^\beta_\delta)\) is \(\beta\)-allowable.  We construct
a new \(M_1\)-bookkeeping function \(F^\alpha\) such that the
\(\alpha\)-allowable recursion guided by \(F^\alpha\) produces the same
underlying hybrid presentation
\[
   \langle\mathcal R_\eta\mid \eta\leq\delta\rangle.
\]
At trivial, relative Random, relative amoeba, Cohen, direct coding, and well-order coding stages
we keep the same bookkeeping entry.  These clauses are independent of the rank
of allowability.

At a uniformization stage, suppose first that the original construction uses
clause~(b).  Then no witness for clause~(a) exists at any rank below \(\beta\),
hence none exists at any rank below \(\alpha\).  The \(\alpha\)-construction
therefore also uses clause~(b), coding the same real.

Suppose next that the original construction uses clause~(a).  Let
\(\xi<\beta\) be the least rank witnessing clause~(a), let \(\dot y_0\) be the
chosen local name, and let \(\dot u\) be the name actually coded by the stage.
If \(\xi<\alpha\), then the same witness is available in the
\(\alpha\)-allowable construction, and the same name \(\dot u\) is coded.  If
\(\alpha\leq\xi<\beta\), then by minimality of \(\xi\) there is no witness to
clause~(a) of rank below \(\alpha\).  In this case we alter the bookkeeping
entry so that it decodes to the uniformization tuple
\[
   (\dot x,\dot u,\dot m),
\]
where \(\dot u\) is the name actually coded by the original construction.  If
\(\alpha=0\), this is a \(0\)-allowable coding stage.  If \(\alpha\) is a
successor, clause~(a) fails below \(\alpha\), so clause~(b) prescribes the same
coding operation.  Thus the same underlying forcing stage is reproduced.

This defines \(F^\alpha\) recursively.  The definition uses only \(M_1\)-codes
for the original bookkeeping entries, the canonical localized presentations,
the selected names, and the relevant ordinals.  Hence \(F^\alpha\in M_1\), and
the \(\alpha\)-allowable construction guided by \(F^\alpha\) produces the same
underlying hybrid presentation.
\end{proof}

\begin{lemma}[Allowability of tails]\label{lem:tail-allowability-local}
Let \((\mathcal R_\delta,\dot I_\delta)\) be \(\rho\)-allowable over \(W\),
witnessed by a bookkeeping function \(F\in M_1\), and let
\[
   (\mathcal R_\gamma,\dot I_\gamma),\qquad \gamma\leq\delta,
\]
be the canonical initial hybrid presentations and auxiliary names produced by
this construction.  If \(\beta\leq\gamma\leq\delta\), then
\[
   (\mathcal R_\gamma,\dot I_\gamma)
   \triangleright_\rho
   (\mathcal R_\beta,\dot I_\beta)
\]
as a relative \(\rho\)-allowable extension in \(W\).  Consequently, if
\(G_\beta\subseteq\mathbb P_\beta\) is generic over \(W\), then the quotient
\[
   \mathbb P_\gamma/G_\beta
\]
is a \(\rho\)-allowable tail over \(W[G_\beta]\).  If \(\xi<\rho\), then the
same quotient also has a \(W\)-presentation as a relative \(\xi\)-allowable
tail.  Finally, relative allowable extensions are transitive.
\end{lemma}

\begin{proof}
Restrict the bookkeeping construction witnessing the \(\rho\)-allowability of
\((\mathcal R_\delta,\dot I_\delta)\) to the interval \([\beta,\gamma)\).  This
restricted construction starts with the already built pair
\((\mathcal R_\beta,\dot I_\beta)\) and then appends exactly the same hybrid steps
which occur in the original construction between stages \(\beta\) and
\(\gamma\).  At successor stages, the decision is made from the current full
hybrid forcing, the current auxiliary name, the decoded bookkeeping entry, the
localized presentations of the names involved, and the forcing relation over
that current hybrid forcing.  These are exactly the data available in the
relative construction starting from \((\mathcal R_\beta,\dot I_\beta)\).  If the
stage is a relative Random or amoeba stage, its admissibility witness
\[
   (A_{\mathrm{rb}},\dot W_{\mathrm{rb}},B^0_{\mathrm{rb}},\dot B_{\mathrm{rb}})
\]
is carried with the stage: the tail uses the forcing computed in
\(\dot W_{\mathrm{rb}}\), and the lifted footprint name records the branch
coordinates generated by the local support.  At limit stages both constructions
take the same mixed-support hybrid limit and the same canonical union of the
preceding auxiliary names.

The quotient assertion follows by interpreting the relative presentation in
\(W\) by a \(\mathbb P_\beta\)-generic filter.  The assertion for lower
ranks follows from Lemma~\ref{lem:shrinking-allowability-local}, applied to the
relative tail.  Transitivity is obtained by concatenating the two relative tail
presentations, renaming reservoir coordinates to fresh coordinates if
necessary.
\end{proof}

\begin{lemma}[Products of allowable local hybrid presentations]
\label{lem:rho-allowable-products-local}
Let \(\rho\) be an ordinal.  For \(i<2\), suppose that
\[
   (\mathcal R^i,\dot I^i)
\]
is \(\rho\)-allowable over \(W\), witnessed by a bookkeeping function
\(F_i\in M_1\).  Let \(H\subseteq Br\) be the fixed branch generic used to form
\(W=M_1[H]\).  For \(i<2\), let \(P^i\) be the terminal post-branch forcing over
\(W\) represented by \(\mathcal R^i\).  Then the ordinary product
\(P^0\times P^1\), computed over \(W\), has a canonical \(\rho\)-allowable local
hybrid presentation.

More precisely, there are a local hybrid presentation \(\mathcal R^\otimes\) over
\(W\), an auxiliary name \(\dot I^\otimes\) over its terminal post-branch forcing,
and a bookkeeping function \(F^\otimes\in M_1\), definable from \(F_0\) and
\(F_1\), such that
\[
   (\mathcal R^\otimes,\dot I^\otimes)
\]
is \(\rho\)-allowable over \(W\), witnessed by \(F^\otimes\), and its associated
terminal post-branch forcing is canonically forcing equivalent to
\(P^0\times P^1\).
\end{lemma}

\begin{proof}
We argue by induction on \(\rho\).  Write
\[
   \mathcal R^i=\langle \mathcal R^i_\beta\mid \beta\leq\delta_i\rangle
   \qquad (i<2),
\]
and let \(P^i_\beta\) denote the post-branch forcing over \(W\) represented by
\(\mathcal R^i_\beta\); thus \(P^i=P^i_{\delta_i}\).  The product presentation
\(\mathcal R^\otimes\) is obtained by first running \(\mathcal R^0\), and then
running a shifted copy of \(\mathcal R^1\).  Every name \(\dot\tau\) occurring in the second block is
replaced by its canonical lift \(\dot\tau^\uparrow\) to the corresponding
product initial segment, and all reservoir coordinates in the second block are
renamed to fresh coordinates.  Let \(P^\otimes_\gamma\) denote the post-branch forcing associated with the
initial product presentation \(\mathcal R^\otimes_\gamma\).  Thus, for every
\(\beta\leq\delta_1\),
\[
   P^\otimes_{\delta_0+\beta}
   \cong
   P^0\times P^1_\beta .
\]
The auxiliary names are defined by
\[
   \dot I^\otimes_{\delta_0+\beta}
   =
   \dot I^0\cup(\dot I^1_\beta)^\uparrow
\]
with the evident interpretation at stages \(\leq\delta_0\).  Limit stages use
the same mixed-support limits as the two original presentations.

For \(\rho=0\), this concatenation witnesses \(0\)-allowability.  The ordinary
Cohen, direct coding, well-order coding, and trivial stages are lifted
second-block stages.  If a shifted second-block stage is a relative regularity
stage with tag
\[
   (\mathrm{Rand};\dot a_0,\ldots,\dot a_{k-1}),
   \quad
   (\mathrm{Am}_{\mathcal N};\dot a_0,\ldots,\dot a_{k-1}),
   \quad\text{or}\quad
   (\mathrm{Am}_{\mathcal M};\dot a_0,\ldots,\dot a_{k-1}),
\]
then the product presentation uses the shifted tag obtained by lifting the
finitely many real names to the product initial segment.  If the original stage
has canonical admissibility witness
\[
   (A_{\mathrm{rb}},\dot W_{\mathrm{rb}},B^0_{\mathrm{rb}},\dot B_{\mathrm{rb}}),
\]
where
\[
   \dot W_{\mathrm{rb}}=L[T_2,\dot b_0,\ldots,\dot b_{k-1}]
\]
for the least tuple of names reading the displayed real parameters, then the
shifted stage has the lifted witness obtained from
\[
   A_{\mathrm{rb}}^\otimes,
   \qquad
   \langle \dot b_i^\uparrow\mid i<k\rangle,
   \qquad
   B^0_{\mathrm{rb}},
   \qquad
   \dot B_{\mathrm{rb}}^\uparrow.
\]
Here \(A_{\mathrm{rb}}^\otimes\) is the shifted copy of
\(A_{\mathrm{rb}}\) in the second block, each \(\dot b_i^\uparrow\) is the
canonical lift of \(\dot b_i\) to the product initial segment, and
\(\dot B_{\mathrm{rb}}^\uparrow\) is the canonical lift of the branch-footprint
name.  The lifted regularity base is therefore
\[
   L[T_2,\dot b_0^\uparrow,\ldots,\dot b_{k-1}^\uparrow]
\]
(with the evident meaning \(L[T_2]\) if \(k=0\)).  The Random or amoeba forcing
is computed in this lifted base, not in the larger product extension.  Hence the
displayed product equivalence is preserved at the shifted regularity stage.

Now assume \(\rho=\alpha+1\).  By the successor uniformization clause we mean
the two alternatives (a) and (b) in Definition~\ref{def:rho-allowable-local},
evaluated at a uniformization tag
\[
   (\mathrm{UF},\dot x,\dot z,\dot m).
\]
Thus clause~(a) searches through ranks \(\zeta<\rho\) for a localized name
\(\dot y\) which is forced into the relevant \(A_{\dot m}\)-section and cannot
be forced out by any relative \(\zeta\)-allowable tail; clause~(b) is used when
no such witness exists.

Consider a shifted uniformization stage \(\delta_0+\beta\), and let
\[
   (\mathrm{UF},\dot x,\dot z,\dot m)
\]
be the corresponding second-block bookkeeping value.  For \(\zeta<\rho\) and a
localized \(P^1_\beta\)-name \(\dot y\), write
\[
   \mathsf{Kick}^1_\beta(\zeta,\dot y)
\]
for the assertion that there are a relative \(\zeta\)-allowable extension
\[
   (\mathcal S,\dot J)\triangleright_\zeta
   (\mathcal R^1_\beta,\dot I^1_\beta)
\]
and, writing \(P_{\mathcal S}\) for the terminal post-branch forcing associated
with \(\mathcal S\), a quotient condition \(p\in P_{\mathcal S}/\dot G^1_\beta\)
such that
\[
   p\Vdash
   (\dot x^{\uparrow\mathcal S},
    \dot y^{\uparrow\mathcal S})
      \notin A_{\dot m^{\uparrow\mathcal S}}.
\]
Define \(\mathsf{Kick}^\otimes_\beta(\zeta,\dot y^\uparrow)\) analogously over
the product initial segment
\[
   (\mathcal R^\otimes_{\delta_0+\beta},
    \dot I^\otimes_{\delta_0+\beta}).
\]
Then
\[
   \mathsf{Kick}^1_\beta(\zeta,\dot y)
   \quad\Longleftrightarrow\quad
   \mathsf{Kick}^\otimes_\beta(\zeta,\dot y^\uparrow).
\]
The implication from left to right follows by taking the product of the fixed
first block with the displayed relative \(\zeta\)-allowable tail and applying
the induction hypothesis.  The implication from right to left follows by
restricting the witnessing product tail to the complete subpresentation
generated by the lifted second-block support of
\((\dot x,\dot y,\dot m)\), using Lemmas~\ref{lem:tail-allowability-local}
and~\ref{lem:shrinking-allowability-local}.

Similarly,
\[
   P^1_\beta\Vdash (\dot x,\dot y)\in A_{\dot m}
   \quad\Longleftrightarrow\quad
   P^0\times P^1_\beta
      \Vdash
      (\dot x^\uparrow,\dot y^\uparrow)\in A_{\dot m^\uparrow}.
\]
Hence the set of clause~(a) witnesses in the product construction is exactly
the lift of the set of clause~(a) witnesses in the second construction.  Thus
clause~(a) holds in the shifted product stage iff it holds in the original
second-block stage.  When it holds, the least rank and the \(<_{M_1}\)-least
localized witness are the lifted ones.  When it fails, clause~(b) uses the
lifted bookkeeping guess \(\dot z^\uparrow\).  The auxiliary name is updated by
\[
   \dot I^0\cup(\dot I^1_{\beta+1})^\uparrow.
\]
This proves the successor step.

Finally let \(\lambda\) be a nonzero limit ordinal.  If both inputs are
\(\lambda\)-allowable, then by the induction hypothesis the product
presentation is \(\zeta\)-allowable for every \(\zeta<\lambda\).  Since
\[
   \mathcal A^W_\lambda
   =
   \bigcap_{\zeta<\lambda}\mathcal A^W_\zeta,
\]
the product presentation is \(\lambda\)-allowable.  The bookkeeping witness at
a limit level is obtained by packaging the lower-level witnesses together with
the concatenation description.  Thus \(F^\otimes\in M_1\) is definable from
\(F_0\) and \(F_1\).
\end{proof}

\begin{lemma}[Tentative values remain in their sections]
\label{lem:tentative-values-remain-local}
Let \((\mathcal R,\dot I)\) be \(\alpha\)-allowable over \(W\) with respect to an
\(M_1\)-bookkeeping function \(F\), and let \(\mathbb P\) be the terminal
post-branch forcing associated with \(\mathcal R\).
Let \(G\subseteq\mathbb P\) be generic over \(W\), and let
\(I^G=\dot I^G\).  If
\[
   (x,y,m,\xi)\in I^G
\]
for some \(\xi<\alpha\), then
\[
   W[G]\models (x,y)\in A_m.
\]
Moreover, for every relative \(\xi\)-allowable extension in \(W\)
\[
   (\mathcal S,\dot J)\triangleright_\xi(\mathcal R,\dot I),
\]
the quotient over \(G\) forces preservation of this membership: if \(P'\) is the
terminal post-branch forcing associated with \(\mathcal S\), then
\[
   P'/G\Vdash (x,y)\in A_m.
\]
\end{lemma}

\begin{proof}
Suppose that \(\beta\) is the stage at which the tuple
\((x,y,m,\xi)\) is added to the auxiliary name.  At that stage clause~(a) of
Definition~\ref{def:rho-allowable-local} was used.  Let
\[
   \dot x_\beta,\qquad
   \dot y_\beta,\qquad
   \dot m_\beta
\]
be the canonical names chosen there, and let \(\xi_\beta=\xi\).  Clause~(a)
gives
\[
   \mathbb P_\beta\Vdash
   (\dot x_\beta,\dot y_\beta)\in A_{\dot m_\beta},
\]
and it also gives the preservation statement saying that no relative
\(\xi_\beta\)-allowable tail over the stage-\(\beta\) initial segment can force
the lifted pair out of the relevant section.

By Lemma~\ref{lem:tail-allowability-local}, the interval from stage \(\beta\)
to the terminal forcing is a relative \(\xi_\beta\)-allowable tail.  Therefore
no condition in the terminal forcing can force
\[
   (\dot x_\beta,\dot y_\beta)\notin A_{\dot m_\beta}.
\]
The forcing theorem gives the membership in the final extension.  If
\((\mathcal S,\dot J)\triangleright_\xi(\mathcal R,\dot I)\) is any further relative
\(\xi\)-allowable extension, then the concatenation of the interval from
\(\beta\) to the terminal forcing with this further extension is again a
relative \(\xi\)-allowable extension over stage \(\beta\).  The same
preservation statement therefore rules out any quotient condition forcing the
pair out of \(A_m\).
\end{proof}

\paragraph{Stabilization of the allowability hierarchy.}
We next pass from the transfinite derivative hierarchy to its stabilized level.
The following convention is used only to make the derivative classes into
subsets of one fixed set.  It is unrelated to the projective coding predicate
\(\Phi\).

Two pairs \((\mathcal R,
\dot I)\) and \((\mathcal S,\dot J)\), consisting of a local
hybrid presentation over \(W\) and its auxiliary name, are identified if their
associated forcings are canonically isomorphic over \(W\), the isomorphism
preserves the order of stages and sends the auxiliary name \(\dot I\) to
\(\dot J\).  This identification only removes harmless choices such as the names
of fresh reservoir coordinates.  For each equivalence class we take its
\(<_{M_1}\)-least representative and denote it by
\[
   \operatorname{rep}_W(\mathcal R,\dot I).
\]
Let \(\mathcal C^W\) be the set of all such representatives for pairs of the
form used in the definitions of \(\rho\)-allowability.  For every ordinal
\(\rho\) put
\[
   \mathcal A^W_\rho
   =
   \{\operatorname{rep}_W(\mathcal R,\dot I)
      \mid (\mathcal R,\dot I)\text{ is }\rho\text{-allowable over }W\}
   \subseteq \mathcal C^W .
\]
By Lemma~\ref{lem:shrinking-allowability-local}, the sequence
\(\langle\mathcal A^W_\rho\mid \rho\in\operatorname{Ord}\rangle\) is decreasing,
and at limit stages it is defined by intersection.  Since all terms are subsets
of the single set \(\mathcal C^W\), the sequence stabilizes.

\begin{definition}[The stabilized level and \(\infty\)-allowability]
\label{def:infty-allowable-local}
Let \(\alpha_0\) be the least ordinal \(\alpha\) such that
\[
   \mathcal A^W_\beta=\mathcal A^W_\alpha
   \qquad\text{for every }\beta\geq\alpha .
\]
Set
\[
   \alpha^\ast=\alpha_0+1 .
\]
A pair \((\mathcal R,\dot I)\), consisting of a local hybrid presentation over
\(W\) and its auxiliary name, is called \emph{\(\infty\)-allowable over \(W\)} if
\[
   \operatorname{rep}_W(\mathcal R,\dot I)
      \in \mathcal A^W_{\alpha_0}.
\]
Equivalently,
\[
   \operatorname{rep}_W(\mathcal R,\dot I)
      \in \mathcal A^W_{\alpha^\ast},
\]
since the hierarchy has stabilized at \(\alpha_0\).
\end{definition}

The passage from \(\alpha_0\) to \(\alpha^\ast=\alpha_0+1\) is only a technical
convenience.  The class has already stabilized at \(\alpha_0\), but
\(\alpha^\ast\) is a successor ordinal, so the successor clause of
Definition~\ref{def:rho-allowable-local} can be used literally.  Thus, whenever
an \(\infty\)-allowable construction reaches a uniformization tag
\[
   (\mathrm{UF},\dot x,\dot z,\dot m)
\]
after an initial segment \((\mathcal R_\beta,\dot I_\beta)\), it evaluates
clause~(a) of
Definition~\ref{def:rho-allowable-local} with \(\rho=\alpha^\ast\).  More
explicitly, writing
\[
   A=\operatorname{supp}_{\mathrm{loc}}((\dot x,\dot m),\mathbb P_\beta),
\]
a candidate is a \(\mathbb P_A\)-name \(\dot y\) for a real.  In the displayed
test below, \(P_{\mathcal S}\) denotes the terminal post-branch forcing associated
with the relative extension \(\mathcal S\).  The candidate is accepted at rank
\(\zeta<\alpha^\ast\) iff
\[
   \mathbb P_\beta\Vdash
   (\dot x,\dot y^{\uparrow\beta})\in A_{\dot m}
\]
and
\[
\begin{split}
   \forall (\mathcal S,\dot J)
   \bigl(& (\mathcal S,\dot J)\triangleright_\zeta
            (\mathcal R_\beta,\dot I_\beta) \\[-1mm]
   &\Rightarrow
   \mathbb P_\beta\Vdash
   \text{``there is no }p\in P_{\mathcal S}/\dot G_\beta
   \text{ such that } \\[-1mm]
   &\hspace{33mm}
   p\Vdash
   (\dot x^{\uparrow\mathcal S},\dot y^{\uparrow\mathcal S})
      \notin A_{\dot m^{\uparrow\mathcal S}}\text{''}\bigr).
\end{split}
\]
If such candidates exist, the construction chooses the least possible
\(\zeta\), and then the \(<_{M_1}\)-least localized presentation of such a name
\(\dot y\) in the sense of Definition~\ref{def:localized-presentation}.  If no
candidate satisfies the displayed preservation test, clause~(b) of
Definition~\ref{def:rho-allowable-local} is used.

\paragraph{Countable local killing tails.}
We shall use the following consequence of the local-support convention.  Suppose
that \((\mathcal R,\dot I)\) is a \(\rho\)-allowable pair over \(W\).  Let
\(\mathbb P\) be the terminal post-branch forcing associated with \(\mathcal R\).
Suppose that \(\dot x,\dot y\) are \(\mathbb P\)-names for reals, and that
\(\dot m\) is a \(\mathbb P\)-name for a natural number.  If there is a relative
\(\rho\)-allowable extension
\[
   (\mathcal S,\dot J)\triangleright_\rho(\mathcal R,\dot I)
\]
and a quotient condition \(p\) in the associated quotient which forces
\[
   (\dot x^{\uparrow\mathcal S},\dot y^{\uparrow\mathcal S})
      \notin A_{\dot m^{\uparrow\mathcal S}},
\]
then such a witness may be chosen so that \(\mathcal S\) is obtained from
\(\mathcal R\) by appending a countable local hybrid tail.  Indeed, write the
complement of \(A_m\) as \(\exists r\;\theta(r,x,y,m)\), with
\(\theta\) a \(\Pi^1_3\) formula.  Strengthen the quotient condition, if
necessary, and use the maximum principle to choose a quotient name \(\dot r\)
such that the strengthened condition forces
\(\theta(\dot r,\dot x^{\uparrow\mathcal S},
\dot y^{\uparrow\mathcal S},\dot m^{\uparrow\mathcal S})\).  The condition and
the names \(\dot x,\dot y,\dot m,\dot r\) are read on a countable admissible
support in the appended tail, after closing under the local supports appearing
in regularity-base witnesses and under the names coded at coding stages.  The
regularity stages in this restricted tail keep their lifted footprint names and
still use the forcings computed in their witnessed bases.  The induced complete
subpresentation is a countable local hybrid tail; by
Lemmas~\ref{lem:tail-allowability-local} and
\ref{lem:shrinking-allowability-local}, it is still relative
\(\rho\)-allowable.  Completeness preserves the displayed forcing statement.

For later reference we also isolate the intentional uses of the coding machinery.
Let \((\mathcal R,\dot I)\) be a \(\rho\)-allowable pair over \(W\), witnessed by a
bookkeeping function \(F\).  Let \(\mathbb P\) be the terminal post-branch forcing
associated with \(\mathcal R\), and let \(G\subseteq\mathbb P\) be generic.
We write
\[
   \operatorname{Int}_\Phi(\mathcal R,F,G)
\]
for the set of reals intentionally submitted to the coding machinery by the
presentation.  Thus, if a direct coding stage uses a name \(\dot w\), the
corresponding element is \(\dot w^{G_\beta}\).  If a well-order stage uses a tag
\((\dot z_0,\dot z_1)\), let \((A_i,\sigma_i)\) be the
\(<_{M_1}\)-least localized presentation of \(\dot z_i\), for \(i<2\).  If these
presentations are distinct, the stage contributes the well-order tag with the
smaller presentation placed first.  If a uniformization stage is used, the stage
contributes the real
\[
   \langle \mathrm{UF},x,u,m\rangle
\]
selected by clause~(a) or clause~(b) of
Definition~\ref{def:rho-allowable-local}.  With this notation,
\[
   w\in\operatorname{Int}_\Phi(\mathcal R,F,G)
   \quad\Longrightarrow\quad
   \Phi(w).
\]
Conversely, Lemma~\ref{lem:no-unwanted-codes} applies to the underlying locally
allowable coding presentation: in an allowable extension, every new instance of
\(\Phi(w)\) is produced at an intentional coding stage, except for instances
already present in the previous initial segment.


\section{The final iteration}
\label{sec:final-iteration}

We now define the forcing which will give the final model.  The definition is
the local $M_1$-analogue of the last iteration in the
$\Pi^1_3$-uniformization and $\Delta^1_3$-well-order construction.  The
iteration uses the stabilized class of $\infty$-allowable local hybrid
presentations from Definition~\ref{def:infty-allowable-local}.  The bookkeeping
places ordinary Cohen stages cofinally often and, for every finite tuple of real
names, places the corresponding relative Random, measure-amoeba, and
category-amoeba stages cofinally often over the associated base.  The bases used
in the $M_1$-case are precisely the finite-real-parameter models
\[
   L[T_2,a_0,\ldots,a_{k-1}]
\]
whose real parameters are read on countable admissible supports.  In particular,
the one-parameter bases $L[T_2,a]$ are included.  The stages over these bases
yield the covering statement used with the Martin--Solovay tree $T_2$ in
Hjorth's proof of Lebesgue measurability and the Baire property.

Throughout this section, \(\omega_2\) denotes the ordinal \(\omega_2^{M_1}\), equivalently the same \(\omega_2\) in the preliminary branch extension \(W\).  We fix, in $M_1$, a bookkeeping function
\[
   F:\omega_2\longrightarrow H_{\omega_2}^{M_1}
\]
with the following properties.  The values of $F$ are decoded relative to the
current associated hybrid forcing.  There is a cofinal subset of $\omega_2$ on
which $F$ is the ordinary Cohen tag.  Moreover, whenever finitely many real
names over a countable admissible support determine a regularity stage with
canonical admissibility witness
\[
   (A_{\mathrm{rb}},\dot W_{\mathrm{rb}},B^0_{\mathrm{rb}},\dot B_{\mathrm{rb}}),
\]
where
\[
   \dot W_{\mathrm{rb}}=L[T_2,\dot a_0,\ldots,
      \dot a_{k-1}],
\]
each of the corresponding finite-real-parameter tags
\[
   (\mathrm{Rand};\dot a_0^{\uparrow},\ldots,
      \dot a_{k-1}^{\uparrow}),
   \qquad
   (\mathrm{Am}_{\mathcal N};\dot a_0^{\uparrow},\ldots,
      \dot a_{k-1}^{\uparrow}),
   \qquad
   (\mathrm{Am}_{\mathcal M};\dot a_0^{\uparrow},\ldots,
      \dot a_{k-1}^{\uparrow})
\]
appears cofinally often after the stage at which the tuple becomes meaningful.
Thus the regularity bookkeeping ranges over the allowed bases of the form
\(L[T_2,\vec a]\), not over unspecified named models.  The forcing at such a
stage is the one computed in the lifted base generated by the displayed real
names.  The witness records the support on which the real parameters \(\dot a_i\)
are read, the countable branch-product support \(B^0_{\mathrm{rb}}\) of the
\(W\)-parameters occurring in those names, and the branch footprint
\(\dot B_{\mathrm{rb}}\).  Finally, whenever an $M_1$-description becomes
meaningful over some initial segment as one of the following objects,
\[
   (\mathrm{WO},\dot z_0,
   \dot z_1),
   \qquad
   (\mathrm{UF},\dot x,
   \dot y,
\dot m),
\]
where the displayed symbols are names for reals, and $\dot m$ is a name for a natural number, that code appears cofinally often after the stage at which it becomes meaningful.  We use disjoint recursive real codings for the two kinds of coding actions and write
\[
   \langle \mathrm{WO},z_0,z_1\rangle,
   \qquad
   \langle \mathrm{UF},x,y,m\rangle
\]
for the corresponding tagged reals.  Thus the well-order coding and the uniformization coding are syntactically disjoint.

\begin{definition}[The main iteration]\label{def:main-iteration}
We define, in $W$, an increasing sequence
\[
   \langle (\mathcal R_\beta,\dot I_\beta)
      \mid \beta\leq\omega_2\rangle
\]
of local hybrid presentations in $W$ and auxiliary names.  The subscript $\beta$ denotes the outer bookkeeping stage.  In the second uniformization case below, the extension appended at the stage is required to be a countable local hybrid tail over the current presentation.  Thus each successor step adds either one hybrid coordinate or a countable block of hybrid coordinates.  Since $\omega_2$ is regular, every $\mathcal R_\beta$ is again a single local hybrid presentation of length at most $\omega_2$.

Set $\mathcal R_0$ to be the trivial post-branch presentation, so that
\[
   \mathbb P_0=\mathbf 1,
\]
and the corresponding full forcing over $M_1$ is simply $Br$.
Let $\dot I_0$ be the canonical name for the empty set.  At a limit stage $\lambda\leq\omega_2$, let $\mathcal R_\lambda$ be the mixed-support direct limit of the earlier presentations, with countable support on reservoir coordinates and finite support on the c.c.c. real-adding coordinates, and let $\dot I_\lambda$ be the canonical name for the union of the earlier interpreted auxiliary names.

Assume that $\beta<\omega_2$ and that $(\mathcal R_\beta,\dot I_\beta)$ has been constructed.  Let
\(\mathbb P_\beta\) be the post-branch forcing associated with \(\mathcal R_\beta\).
The next step is chosen by the following cases.

\medskip
\noindent\textbf{Relative regularity and Cohen stages.}
In these cases the hybrid successor step first adds a fresh reservoir
coordinate.  Thus choose a fresh copy \(\mathbb C_\beta\) of
\(\mathbb C_{M_1}\), let \(\dot g_\beta\) be its generic, and put
\[
   \widehat{\mathbb P}_{\beta+1}=\mathbb P_\beta\times\mathbb C_\beta.
\]
If \(F(\beta)\) is one of the relative regularity tags
\[
   (\mathrm{Rand};\dot a^\beta_0,\ldots,
      \dot a^\beta_{k_\beta-1}),
   \qquad
   (\mathrm{Am}_{\mathcal N};\dot a^\beta_0,\ldots,
      \dot a^\beta_{k_\beta-1}),
   \qquad
   (\mathrm{Am}_{\mathcal M};\dot a^\beta_0,\ldots,
      \dot a^\beta_{k_\beta-1}),
\]
and this stage has a canonical admissibility witness in the sense of
Definition~\ref{def:admissible-random-base}, let
\[
   (A^{\mathrm{reg}}_\beta,
    \dot W^{\mathrm{reg}}_\beta,
    B^{0,\mathrm{reg}}_\beta,
    \dot B^{\mathrm{reg}}_\beta)
\]
be that witness.  Let \(\dot W^{\mathrm{reg},+}_\beta\) be the corresponding
\(\widehat{\mathbb P}_{\beta+1}\)-name for the lifted base.  Thus
\[
   \dot W^{\mathrm{reg},+}_\beta
   =L[T_2,(\dot a^{\beta,\mathrm{reg}}_0)^{\uparrow},\ldots,
      (\dot a^{\beta,\mathrm{reg}}_{k_\beta-1})^{\uparrow}]
\]
for the least tuple of names reading the displayed real parameters.  Let
\(\dot{\mathbb Q}_\beta\) be the member of the following list corresponding to
the tag:
\[
   \dot{\mathbb B}_{\mathrm{rand}}^{\dot W^{\mathrm{reg},+}_\beta},
   \qquad
   \dot{\mathbb A}_{\mathcal N}^{\dot W^{\mathrm{reg},+}_\beta},
   \qquad
   \dot{\mathbb A}_{\mathcal M}^{\dot W^{\mathrm{reg},+}_\beta}.
\]
Then
\[
   \mathbb P_{\beta+1}=\widehat{\mathbb P}_{\beta+1}*\dot{\mathbb Q}_\beta.
\]
If the displayed regularity tag has no admissibility witness, the second-step
forcing is trivial and we set
\[
   \mathbb P_{\beta+1}=\widehat{\mathbb P}_{\beta+1}.
\]
If \(F(\beta)\) is the ordinary Cohen
tag, the second-step forcing is ordinary Cohen forcing on \(\omega\):
\[
   \mathbb P_{\beta+1}=\widehat{\mathbb P}_{\beta+1}*\dot{\mathbb C}_{\omega}.
\]
In all these cases \(\dot I_\beta\) is lifted canonically.  Ordinary Cohen
forcing is distinct from the \(M_1\)-Cohen reservoir coordinate attached at the
beginning of the hybrid successor step.

\medskip
\noindent\textbf{Well-order stages.}
Suppose that $F(\beta)$ decodes to
\[
   (\mathrm{WO},\dot z_0,
   \dot z_1),
\]
where $\dot z_0$ and $\dot z_1$ are $\mathbb P_\beta$-names for reals.  Let $(a_i,\sigma_i)$ be the canonical localized presentation of $\dot z_i$, for $i<2$.  As in every hybrid successor step, this case uses the fresh reservoir coordinate attached at the beginning of the stage.  If the two canonical localized presentations are equal, the second-step forcing over that reservoir coordinate is trivial and the auxiliary name is lifted.  Otherwise compare them by $<_{M_1}$.  If
\[
   (a_0,\sigma_0)<_{M_1}(a_1,\sigma_1),
\]
use the second-step coding forcing for
\[
   \operatorname{Code}(\langle\mathrm{WO},\dot z_0,
      \dot z_1\rangle).
\]
If
\[
   (a_1,\sigma_1)<_{M_1}(a_0,\sigma_0),
\]
use the second-step coding forcing for
\[
   \operatorname{Code}(\langle\mathrm{WO},\dot z_1,
      \dot z_0\rangle).
\]
In either nontrivial case this means, as in Definition~\ref{def:zero-local-allowable}, that the nontrivial second step over the reservoir coordinate already attached at this stage is the Jensen--Solovay almost-disjoint coding of the corresponding reshaped David code into a real.  The auxiliary name is lifted unchanged.

\medskip
\noindent\textbf{Uniformization stages.}
Suppose that $F(\beta)$ decodes to
\[
   (\mathrm{UF},\dot x,
   \dot y,
   \dot m),
\]
where $\dot x$ and $\dot y$ are $\mathbb P_\beta$-names for reals and $\dot m$ is a $\mathbb P_\beta$-name for a natural number.  Let
\[
   L_\beta=
   \operatorname{supp}_{\mathrm{loc}}((\dot x,
      \dot m),\mathbb P_\beta).
\]
We apply the stabilized $\alpha^\ast$-test from Definition~\ref{def:infty-allowable-local}.  Thus we ask whether there are an ordinal $\zeta<\alpha^\ast$ and a $\mathbb P_{L_\beta}$-name $\dot y_0$ for a real such that, after canonically lifting $\dot y_0$ to a $\mathbb P_\beta$-name,
\[
   \mathbb P_\beta\Vdash
   (\dot x,\dot y_0^{\uparrow\beta})\in A_{\dot m},
\]
and no relative $\zeta$-allowable extension in $W$ can force this pair out of the relevant section.  Explicitly, for every
\[
   (\mathcal S,
   \dot J)\triangleright_\zeta
   (\mathcal R_\beta,
   \dot I_\beta)
\]
in $W$, write $P_{\mathcal S}$ for the terminal post-branch forcing associated
with $\mathcal S$.  If the displayed names are lifted to $P_{\mathcal S}$, then
\[
   \mathbb P_\beta\Vdash
   \text{``there is no condition in }
   P_{\mathcal S}/\dot G_\beta
   \text{ forcing }
   (\dot x^{\uparrow\mathcal S},
    \dot y_0^{\uparrow\mathcal S})
      \notin A_{\dot m^{\uparrow\mathcal S}}\text{.''}
\]
If such witnesses exist, choose the least possible $\zeta$ and then the $<_{M_1}$-least localized presentation of a witnessing $\dot y_0$.

In this first case, $\dot y_0$ is put into the tentative uniformizing set with rank $\zeta$.  We do not code $\dot y_0$ itself.  Instead choose the $<_{M_1}$-least usable real name $\dot u$ such that
\[
   \mathbb P_\beta\Vdash
   \dot u\neq \dot y_0^{\uparrow\beta}
\]
and such that the tagged name
\[
   \langle\mathrm{UF},\dot x,
      \dot u,
      \dot m\rangle
\]
has not already been intentionally coded at an earlier uniformization stage.  If the bookkeeping guess $\dot y$ is already forced to be different from $\dot y_0^{\uparrow\beta}$ and has not yet been so coded, we take $\dot u=\dot y$.  Now perform
\[
   \operatorname{Code}(\langle\mathrm{UF},\dot x,
      \dot u,
      \dot m\rangle).
\]
The new auxiliary name is the canonical $\mathbb P_{\beta+1}$-name forced to be
\[
   \dot I_\beta^{\dot G_\beta}
   \cup
   \{(\dot x^{\dot G_\beta},
       (\dot y_0^{\uparrow\beta})^{\dot G_\beta},
       \dot m^{\dot G_\beta},
       \check\zeta)
     \}.
\]

If the stabilized test fails, we use the other side of the $\infty$-allowable dichotomy.  Choose the $<_{M_1}$-least description of a countable relative $\alpha^\ast$-allowable local hybrid tail
\[
   (\mathcal S,
   \dot J)\triangleright_{\alpha^\ast}
   (\mathcal R_\beta,
   \dot I_\beta)
\]
together with a quotient condition witnessing that the current bookkeeping candidate can be forced out of the section.  Writing $P_{\mathcal S}$ for the terminal post-branch forcing associated with $\mathcal S$, after restricting the quotient below that condition we have
\[
   P_{\mathcal S}/\dot G_\beta
   \Vdash
   (\dot x^{\uparrow\mathcal S},
    \dot y^{\uparrow\mathcal S})
      \notin A_{\dot m^{\uparrow\mathcal S}}.
\]
The preceding countable-tail observation ensures that such a countable witness exists whenever any relative allowable tail can kill the candidate.  We replace $\mathcal S$ by this canonical restricted countable tail and set
\[
   (\mathcal R_{\beta+1},\dot I_{\beta+1})=(\mathcal S,
   \dot J).
\]
Thus the stage does not merely remember that some condition could kill the candidate; the appended countable quotient itself forces the negation.

If the value of $F(\beta)$ is not meaningful over the current initial segment, the hybrid successor step still attaches its fresh reservoir coordinate, but the second-step forcing is trivial and $\dot I_\beta$ is lifted canonically.
\end{definition}

Let
\[
   \mathbb P^\ast=\mathbb P_{\omega_2},
   \qquad
   \dot I^\ast=\dot I_{\omega_2}.
\]
If $G^\ast\subseteq\mathbb P^\ast$ is $W$-generic, we write
\[
   W^\ast=W[G^\ast],
   \qquad
   I^\ast=(\dot I^\ast)^{G^\ast}.
\]


\begin{lemma}[The final presentation is \(\infty\)-allowable]
\label{lem:main-iteration-infty-allowable}
\label{fact:main-iteration-infty-allowable}
The pair \((\mathcal R_{\omega_2},\dot I_{\omega_2})\) belongs to \(W\) and is
\(\infty\)-allowable.  More precisely, there is an \(M_1\)-bookkeeping function
\(F^\ast\) such that every initial segment of the final presentation is
\(\alpha^\ast\)-allowable relative to \(F^\ast\).  In addition, ordinary Cohen
forcing occurs cofinally often.  If a finite-real-parameter regularity stage becomes
meaningful, with canonical support witness
\((A_{\mathrm{rb}},\dot W_{\mathrm{rb}},B^0_{\mathrm{rb}},\dot B_{\mathrm{rb}})\),
then the relative Random forcing, measure-amoeba forcing, and category-amoeba
forcing computed in \(\dot W_{\mathrm{rb}}\) occur cofinally often in the final
presentation.
\end{lemma}

\begin{proof}
We argue by induction on the construction stages.  At stage \(0\), the
presentation is the trivial post-branch presentation, which is allowable by
Lemma~\ref{lem:basic-zero-allowable}.  At limit stages, the construction takes
the countable-support direct limit of the earlier presentations.  This is
allowed by the limit clause in the derivative hierarchy defining
\(\alpha\)-allowability.

Consider a successor stage.  The relative regularity cases, the ordinary Cohen
case, and the well-order case are instances of the successor clauses in the
local allowable recursion.  In each of these cases the construction first adds
the fresh \(M_1\)-Cohen reservoir coordinate and then performs the specified
second step, or the trivial second step if the relevant bookkeeping value is not
meaningful.  Therefore the resulting presentation remains
\(\alpha^\ast\)-allowable over the preceding initial segment.

The first uniformization case is the stabilized successor clause from
Definition~\ref{def:infty-allowable-local}.  The chosen tentative value is added
to the auxiliary name with its rank, and the construction intentionally codes
only a competing reserved tag.  Thus the new presentation is again
\(\alpha^\ast\)-allowable over the preceding initial segment.

It remains to check the second uniformization case.  In this case the
construction starts with a relative \(\alpha^\ast\)-allowable continuation
\[
   (\mathcal S,\dot J)\triangleright_{\alpha^\ast}
   (\mathcal R_\beta,\dot I_\beta),
\]
a quotient condition \(q\in P_{\mathcal S}/\dot G_\beta\), where
\(P_{\mathcal S}\) is the terminal post-branch forcing associated with
\(\mathcal S\), and a name \(\dot r\) such that
\[
   q\Vdash
   \theta(\dot r,
      \dot x^{\uparrow\mathcal S},
      \dot y^{\uparrow\mathcal S},
      \dot m^{\uparrow\mathcal S}).
\]
By the local-support convention for hybrid presentations, the condition \(q\),
the name \(\dot r\), the lifted names
\[
   \dot x^{\uparrow\mathcal S},\qquad
   \dot y^{\uparrow\mathcal S},\qquad
   \dot m^{\uparrow\mathcal S},
\]
and the data needed to interpret these objects are contained in a countable
localized subpresentation of \((\mathcal S,\dot J)\) over
\((\mathcal R_\beta,\dot I_\beta)\).  The construction chooses the
\(<_{M_1}\)-least such support and restricts the corresponding presentation
below the witnessing condition \(q\).  The resulting presentation
\[
   (\mathcal S^B,\dot J^B)\triangleright_{\alpha^\ast}
   (\mathcal R_\beta,\dot I_\beta)
\]
is still relative \(\alpha^\ast\)-allowable, because the derivative clauses are
closed under passing to localized subpresentations and under restricting below a
condition.  Moreover its quotient forces
\[
   (\dot x^{\uparrow\mathcal S^B},
    \dot y^{\uparrow\mathcal S^B})
      \notin A_{\dot m^{\uparrow\mathcal S^B}}.
\]
Thus the second uniformization case also preserves
\(\alpha^\ast\)-allowability.

It remains to record the bookkeeping.  Whenever the construction appends a
countable presentation in the second uniformization case, we use the
bookkeeping function witnessing that chosen presentation on the newly added
coordinates.  On the coordinates coming from the original stages we use the
original bookkeeping function \(F\).  Since \(\omega_2\) is regular, the
concatenation of \(\omega_2\)-many countable presentations has length at most
\(\omega_2\), and hence can be indexed by \(\omega_2\) in the order in which it
is built.  This gives a single \(M_1\)-bookkeeping function \(F^\ast\) for the
final presentation.

The cofinality requirements are inherited from \(F\).  The original bookkeeping
function places ordinary Cohen tags cofinally often.  Once a finite tuple of real names becomes meaningful as an admissibly based
regularity stage, with witness
\((A_{\mathrm{rb}},\dot W_{\mathrm{rb}},B^0_{\mathrm{rb}},\dot B_{\mathrm{rb}})\),
\(F\) places the three corresponding finite-real-parameter tags cofinally often,
using the Random, measure-amoeba, and category-amoeba forcings computed in the
lifted base generated by that witness.  Adding
countably many coordinates at individual stages does not destroy cofinality in
the final indexing of length \(\omega_2\).  Therefore the final presentation has
all the required cofinal regularity and Cohen stages, and every initial segment
is \(\alpha^\ast\)-allowable relative to \(F^\ast\).  Hence
\((\mathcal R_{\omega_2},\dot I_{\omega_2})\) is \(\infty\)-allowable.
\end{proof}



The next step will be to prove the final dichotomy.  For each $m$ and $x$, either $I^\ast$ contains a tentative value $(x,y,m,\xi)$ of least rank, in which case Lemma~\ref{lem:tentative-values-remain-local} keeps $(x,y)$ in $A_m$ through all further $\infty$-allowable tails, or else the second uniformization case is applied cofinally often to every candidate name for a value in the $x$-section and the final model forces that section of $A_m$ to be empty.

Before proving the dichotomy we isolate the exactness property of the reserved tags used below.


\begin{lemma}[Exactness for reserved tags]\label{lem:reserved-tag-coding-exactness}
Let
\[
   w=\langle\mathrm{UF},x,y,m\rangle
   \quad\text{or}\quad
   w=\langle\mathrm{WO},z_0,z_1\rangle
\]
be a real of one of the reserved tagged forms in \(W^\ast\).  If
\[
   W^\ast\models\Phi(w),
\]
then this instance of \(\Phi(w)\) was introduced by an intentional coding step of
 the construction.  More precisely, a code for a real of the form
\(\langle\mathrm{UF},x,y,m\rangle\) is introduced only by a uniformization
coding step, and a code for a real of the form
\(\langle\mathrm{WO},z_0,z_1\rangle\) is introduced only by a well-order coding
step.
\end{lemma}

\begin{proof}
Fix a real \(r\in W^\ast\) witnessing \(\Phi(w)\).  Thus \(\Psi(r,w)\) holds in
\(W^\ast\), where \(\Psi\) is the predicate used in the definition of \(\Phi\).
By the local-support property of the hybrid presentations, the real \(w\), the
witness \(r\), and the data needed to verify \(\Psi(r,w)\) are read in a
countable localized subpresentation of the final construction over \(W\).  Closing
this support under the local data used by the relevant names gives a
\(0\)-allowable local hybrid presentation \(\mathcal R_A\) over \(W\) and a
generic filter \(G_A\) such that
\[
   r,w\in W[G_A]
   \quad\text{and}\quad
   W[G_A]\models\Psi(r,w).
\]
Hence
\[
   W[G_A]\models\Phi(w).
\]

Apply Lemma~\ref{lem:no-unwanted-codes} to this localized presentation.  There
is an explicit coding coordinate \(\eta\) of \(\mathcal R_A\) and a name
\(\dot u_\eta\) such that
\[
   w=(\dot u_\eta)^{G_A}.
\]
Thus every occurrence of \(\Phi(w)\) for a reserved tag is produced by an
explicit coding coordinate of the actual local presentation which reads the
witness.

It remains only to identify which coding coordinates may write reserved tags.
The construction uses the tags \(\mathrm{UF}\) and \(\mathrm{WO}\) only for the
reserved uniformization and well-order decisions.  Direct coding clauses do not
use reals whose outer tag is \(\mathrm{UF}\) or \(\mathrm{WO}\).  The same
convention is part of every countable localized presentation appended in the
second uniformization case.  Therefore an explicit coding coordinate which writes
a real of the form \(\langle\mathrm{WO},z_0,z_1\rangle\) must be a well-order
coding step, and an explicit coding coordinate which writes a real of the form
\(\langle\mathrm{UF},x,y,m\rangle\) must be a uniformization coding step.

This proves the claimed exactness of the reserved tags.
\end{proof}

\begin{lemma}[Final dichotomy for the uniformization stages]\label{lem:final-dichotomy-local}
In \(W^\ast\), for every real \(x\in 2^\omega\) and every \(m\in\omega\), exactly one of the following alternatives holds.
\begin{enumerate}
\item There are a real \(y\) and an ordinal \(\xi<\alpha^\ast\) such that
\[
   (x,y,m,\xi)\in I^\ast.
\]
In this case there is a unique real \(y_0\) such that
\[
   W^\ast\models
   (x,y_0)\in A_m
   \quad\text{and}\quad
   \neg\Phi(\langle\mathrm{UF},x,y_0,m\rangle).
\]
Moreover, if \(y\ne y_0\) and \((x,y)\in A_m\), then
\[
   W^\ast\models \Phi(\langle\mathrm{UF},x,y,m\rangle).
\]

\item For every real \(y\) and every ordinal \(\xi<\alpha^\ast\),
\[
   (x,y,m,\xi)\notin I^\ast.
\]
In this case
\[
   W^\ast\models A_{m,x}=\emptyset.
\]
\end{enumerate}
Here \(\langle\mathrm{UF},x,y,m\rangle\) denotes the fixed recursive real coding of the tagged quadruple used at the uniformization coding stages.
\end{lemma}

\begin{proof}
The two alternatives are mutually exclusive and exhaustive by their first displayed clauses.  We prove the corresponding conclusions.

Assume first that
\[
   X_{x,m}=\{(y,\xi)\mid \xi<\alpha^\ast\text{ and }(x,y,m,\xi)\in I^\ast\}
\]
is nonempty.  Let \(\xi_0\) be the least ordinal such that \((y,\xi_0)\in X_{x,m}\) for some \(y\).  Among all names whose interpretation gives such a tuple of rank \(\xi_0\), choose the one with \(<_{M_1}\)-least localized presentation.  Let its value be \(y_0\), and let \(\beta_0<\omega_2\) be the first stage at which the corresponding tuple
\[
   (x,y_0,m,\xi_0)
\]
is inserted into the auxiliary name.  By Lemma~\ref{lem:tentative-values-remain-local}, applied to the tail of the final \(\infty\)-allowable presentation after \(\beta_0\),
\[
   W^\ast\models (x,y_0)\in A_m.        \tag{1}
\]

We show that \(\Phi(\langle\mathrm{UF},x,y_0,m\rangle)\) fails in \(W^\ast\).  Suppose otherwise.  By Lemma~\ref{lem:reserved-tag-coding-exactness}, there is an intentional uniformization coding stage \(\eta<\omega_2\) which codes
\[
   w_0=\langle\mathrm{UF},x,y_0,m\rangle .
\]
Choose \(\eta\) least with this property.

First suppose that \(\eta\geq\beta_0\).  At every stage \(\eta'\geq\beta_0\), the lift of the name for \(y_0\) is still a rank \(\xi_0\) tentative value, by Lemma~\ref{lem:tentative-values-remain-local}.  If the construction at \(\eta\) selected some value different from \(y_0\), then that value would have either smaller rank than \(\xi_0\), or rank \(\xi_0\) and a strictly earlier localized presentation.  Its tuple would belong to \(I^\ast\), contradicting the choice of \((\xi_0,y_0)\).  Hence \(y_0\) is the selected value at \(\eta\).  The first uniformization clause never codes the selected value itself, so \(w_0\) is not coded at \(\eta\), a contradiction.

Now suppose that \(\eta<\beta_0\).  If \(w_0\) is coded at \(\eta\) by the second uniformization case, then there are a relative \(\xi_0\)-allowable continuation \((\mathcal S,\dot J)\) over \((\mathcal R_\eta,\dot I_\eta)\), a quotient condition \(q\), and a name \(\dot r\), all contained in the countable localized subpresentation chosen at stage \(\eta\), such that, writing \(P_{\mathcal S}\) for the terminal post-branch forcing associated with \(\mathcal S\),
\[
   q\Vdash_{P_{\mathcal S}/\dot G_\eta}
      \theta(\dot r,
             \dot x^{\uparrow\mathcal S},
             \dot y_0^{\uparrow\mathcal S},
             \check m),
\]
where \((u,v)\notin A_m\) is written as \(\exists r\,\theta(r,u,v,m)\) with \(\theta\in\Pi^1_3\).  By Lemma~\ref{lem:rho-allowable-products-local} and Lemma~\ref{lem:tail-allowability-local}, composing this continuation with the actual interval from \(\eta\) to \(\beta_0\) gives a relative \(\xi_0\)-allowable continuation over \((\mathcal R_{\beta_0},\dot I_{\beta_0})\) which forces the lift of \((\dot x,\dot y_0)\) out of \(A_m\).  This contradicts the fact that \(y_0\) is inserted at stage \(\beta_0\) as a rank \(\xi_0\) tentative value.

Thus stage \(\eta\) used the first uniformization case.  It selected some value \(y_\eta\) and coded \(y_0\) as a different candidate.  If a relative \(\xi_0\)-allowable continuation over \((\mathcal R_\eta,\dot I_\eta)\) forced the lift of \((\dot x,\dot y_0)\) out of \(A_m\), then the same composition argument would contradict the insertion of \((x,y_0,m,\xi_0)\) at \(\beta_0\).  Hence \(y_0\) was already eligible at \(\eta\) as a rank \(\xi_0\) tentative value.  Since the construction chooses first the least rank and then the \(<_{M_1}\)-least localized presentation, the selected value \(y_\eta\) has rank below \(\xi_0\), or has rank \(\xi_0\) with an earlier localized presentation.  Then the tuple for \(y_\eta\) belongs to \(I^\ast\), again contradicting the choice of \((\xi_0,y_0)\).  Therefore
\[
   W^\ast\models \neg\Phi(\langle\mathrm{UF},x,y_0,m\rangle).       \tag{2}
\]

Let now \(y\ne y_0\) and suppose that \(W^\ast\models (x,y)\in A_m\).  Choose localized names \(\dot x,\dot y,\dot y_0\) and, by the usual mixing below a condition in \(G^\ast\), arrange that
\[
   1\Vdash \dot y\ne\dot y_0,
   \qquad
   \dot x^{G^\ast}=x,
   \quad
   \dot y^{G^\ast}=y,
   \quad
   \dot y_0^{G^\ast}=y_0.
\]
By bookkeeping, there is a stage \(\gamma>\beta_0\) with
\[
   F(\gamma)=(\mathrm{UF},\dot x,\dot y,\check m)
\]
after these names are available.  At stage \(\gamma\), the lift of \(y_0\) is the selected tentative value by the preceding minimality argument.  If \(\langle\mathrm{UF},x,y,m\rangle\) has not already been coded, the first uniformization case codes
\[
   \langle\mathrm{UF},\dot x,\dot y,\check m\rangle,
\]
since \(1\Vdash\dot y\ne\dot y_0\).  Hence, in all cases,
\[
   W^\ast\models \Phi(\langle\mathrm{UF},x,y,m\rangle).        \tag{3}
\]
Equations \((1)\)--\((3)\) prove the first alternative.

Assume now that
\[
   X_{x,m}=\emptyset .                                      \tag{4}
\]
Suppose toward a contradiction that \(W^\ast\models (x,y)\in A_m\).  Choose localized names \(\dot x,\dot y\) for \(x,y\), and choose a later stage \(\gamma<\omega_2\) such that
\[
   F(\gamma)=(\mathrm{UF},\dot x,\dot y,\check m).
\]
The first uniformization case cannot occur at \(\gamma\), since it would add some tuple \((x,y',m,\xi)\) with \(\xi<\alpha^\ast\) to \(I^\ast\), contradicting \((4)\).  Therefore the second uniformization case is used.  Thus the construction appends a countable localized relative \(\alpha^\ast\)-allowable presentation, restricted below a witnessing quotient condition \(q\), and a name \(\dot r\) such that
\[
   q\Vdash
      \theta(\dot r,
             \dot x^{\uparrow},
             \dot y^{\uparrow},
             \check m),                                      \tag{5}
\]
where again \((u,v)\notin A_m\) is written as \(\exists r\,\theta(r,u,v,m)\) with \(\theta\in\Pi^1_3\).

The presentation added at \(\gamma\) is an initial segment of the final presentation, and the remaining quotient is an allowable local hybrid tail.  By the \(T_2\)-absoluteness used for allowable tails, the \(\Pi^1_3\)-statement in \((5)\) remains true in the final extension for the same witness \(r=\dot r^{G^\ast}\).  Hence
\[
   W^\ast\models (x,y)\notin A_m,
\]
contradicting the choice of \(y\).  Therefore \(A_{m,x}=\emptyset\) in \(W^\ast\), as required.
\end{proof}

\begin{corollary}[The final \(\Pi^1_4\)-uniformizing relation]\label{cor:final-pi14-uniformization-local}
In \(W^\ast\), the \(\Pi^1_4\)-uniformization property holds for the relations in the fixed universal list.  More precisely, for each \(m<\omega\) define
\[
   U_m(x,y)
   \quad\Longleftrightarrow\quad
   (x,y)\in A_m
   \ \land\
   \neg\Phi(\langle\mathrm{UF},x,y,m\rangle).
\]
Then \(U_m\) is a \(\Pi^1_4\) relation.  Moreover, for every real \(x\), if the \(x\)-section of \(A_m\) is nonempty, then there is exactly one real \(y\) such that \(U_m(x,y)\).
Consequently \(W^\ast\) satisfies the lightface \(\Pi^1_4\)-uniformization property.  By the usual relativization, it also satisfies the boldface \(\Pi^1_4\)-uniformization property.
\end{corollary}

\begin{proof}
Fix \(m<\omega\) and a real \(x\in W^\ast\).  If the \(x\)-section of \(A_m\) is empty, then no real \(y\) satisfies \(U_m(x,y)\), by definition.  Suppose therefore that \(A_{m,x}\neq\emptyset\).  The second alternative of Lemma~\ref{lem:final-dichotomy-local} is then impossible.  Hence the first alternative holds, and the lemma gives a real \(y_0\) such that
\[
   W^\ast\models (x,y_0)\in A_m
   \quad\text{and}\quad
   W^\ast\models \neg\Phi(\langle\mathrm{UF},x,y_0,m\rangle).
\]
Thus \(U_m(x,y_0)\) holds.

The same lemma also gives uniqueness.  Namely, if \(y\neq y_0\) and \((x,y)\in A_m\), then
\[
   W^\ast\models \Phi(\langle\mathrm{UF},x,y,m\rangle),
\]
so \(U_m(x,y)\) fails.  If \((x,y)\notin A_m\), then \(U_m(x,y)\) fails already by its first conjunct.  Therefore \(y_0\) is the unique value selected by \(U_m\) on the \(x\)-section of \(A_m\).

It remains to record the projective complexity.  By the choice of the universal list, \((x,y)\in A_m\) is \(\Pi^1_4\).  The local coding predicate \(\Phi(w)\) was defined in the coding section as a \(\Sigma^1_4\) predicate.  Therefore
\[
   \neg\Phi(\langle\mathrm{UF},x,y,m\rangle)
\]
is \(\Pi^1_4\).  Since the class \(\Pi^1_4\) is closed under finite conjunctions, \(U_m(x,y)\) is \(\Pi^1_4\).

For the boldface conclusion, let \(A(x,y)\) be \(\Pi^1_4(a)\).  Using the fixed recursive pairing of reals, regard this as a lightface \(\Pi^1_4\) relation in the pair \((a,x)\) and the value \(y\), say as a section of some member of the universal list.  Applying the preceding lightface uniformization to the relation on \(((a,x),y)\) and then fixing the parameter \(a\) gives a \(\Pi^1_4(a)\) uniformizing relation for \(A\).
\end{proof}

\begin{lemma}[The final \(\Delta^1_4\)-well-order]\label{lem:final-delta14-wellorder}
In \(W^\ast\) the reals admit a \(\Delta^1_4\)-definable well-order.
\end{lemma}

\begin{proof}
We first describe the underlying order independently of the projective definition.  If \(z\in W^\ast\) is a real, let \(\pi(z)\) be the \(<_{M_1}\)-least localized presentation, in the final hybrid presentation, of a name whose interpretation is \(z\).  Such a presentation exists because every real in the final mixed-support hybrid iteration has a bounded local support.  If \(z_0\ne z_1\), then \(\pi(z_0)\ne\pi(z_1)\), since the same localized name has only one interpretation in the fixed generic extension.  Hence
\[
   z_0\triangleleft z_1
   \quad\Longleftrightarrow\quad
   z_0\ne z_1\text{ and }\pi(z_0)<_{M_1}\pi(z_1)
\]
defines a well-order of the reals in \(W^\ast\).  It is the order induced by the fixed $<_{M_1}$-well-order of the canonical localized presentations.

We claim that, for distinct reals \(z_0,z_1\in W^\ast\),
\[
   z_0\triangleleft z_1
   \quad\Longleftrightarrow\quad
   W^\ast\models \Phi(\langle\mathrm{WO},z_0,z_1\rangle).
\]
Assume first that \(z_0\triangleleft z_1\).  Choose localized names \(\dot z_0,\dot z_1\) whose final canonical localized presentations are \(\pi(z_0)\) and \(\pi(z_1)\).  The supports of these presentations are bounded in \(\omega_2\), so there is an initial segment \(\mathbb P_\beta\) over which these names are already available.  By the bookkeeping property of \(F\), there is a later stage \(\gamma>\beta\) at which
\[
   F(\gamma)=(\mathrm{WO},\dot z_0,\dot z_1).
\]
At stage \(\gamma\) the well-order clause compares the same canonical localized presentations.  The localization convention and the product/tail lemmas ensure that passing to the larger initial segment, or later appending a countable local tail, does not change the \(<_{M_1}\)-comparison of these presentations.  Since \(\pi(z_0)<_{M_1}\pi(z_1)\), the stage explicitly codes
\[
   \langle\mathrm{WO},\dot z_0,\dot z_1\rangle.
\]
After interpreting by \(G^\ast\), Lemma~\ref{lem:basic-zero-allowable} yields
\[
   W^\ast\models \Phi(\langle\mathrm{WO},z_0,z_1\rangle).
\]

Conversely suppose that \(z_0\ne z_1\) and
\[
   W^\ast\models \Phi(\langle\mathrm{WO},z_0,z_1\rangle).
\]
By Lemma~\ref{lem:reserved-tag-coding-exactness}, this code was introduced at an intentional well-order coding stage of the final presentation.  Thus, at some stage, the construction decoded a well-order tag, compared the canonical localized presentations of two names whose interpretations are \(z_0\) and \(z_1\), and coded the ordered tag with the smaller presentation placed first.  Again the localization and product/tail conventions imply that this comparison agrees with the comparison of the final canonical presentations \(\pi(z_0)\) and \(\pi(z_1)\).  Therefore \(\pi(z_0)<_{M_1}\pi(z_1)\), and hence \(z_0\triangleleft z_1\).  This proves the claim.

Now define, in \(W^\ast\),
\[
   z_0<_{\Delta}z_1
   \quad\Longleftrightarrow\quad
   z_0\ne z_1\ \land\
   \Phi(\langle\mathrm{WO},z_0,z_1\rangle).
\]
By the claim, \(<_{\Delta}\) is exactly \(\triangleleft\), and hence is a well-order of the reals.

It remains to record the projective complexity.  Since \(\Phi\) is the local \(\Sigma^1_4\)-coding predicate, the displayed definition of \(<_{\Delta}\) is \(\Sigma^1_4\).  The claim also gives, for distinct reals \(z_0,z_1\), that exactly one of
\[
   \Phi(\langle\mathrm{WO},z_0,z_1\rangle),
   \qquad
   \Phi(\langle\mathrm{WO},z_1,z_0\rangle)
\]
holds.  Hence, in \(W^\ast\), the same relation is equivalently defined by
\[
   z_0<_{\Delta}z_1
   \quad\Longleftrightarrow\quad
   z_0\ne z_1\ \land\
   \neg\Phi(\langle\mathrm{WO},z_1,z_0\rangle).
\]
This second definition is \(\Pi^1_4\).  Thus \(<_{\Delta}\) is both \(\Sigma^1_4\) and \(\Pi^1_4\), and consequently is a \(\Delta^1_4\)-definable well-order of the reals.
\end{proof}




\begin{lemma}[Localized covering over $L\lbrack T_2,a\rbrack$]\label{lem:random-cohen-covering-final}
Let $a\in W^\ast$ be a real and put
\[
   N_a=L[T_2,a].
\]
Then there is a Borel null set $Z_a^{\mathcal N}\in W^\ast$ such that
\[
   \bigcup\{B\mid B\text{ is a Borel null set coded in }N_a\}
   \subseteq Z_a^{\mathcal N}.
\]
There is also a Borel meager set $Z_a^{\mathcal M}\in W^\ast$ such that
\[
   \bigcup\{B\mid B\text{ is a Borel meager set coded in }N_a\}
   \subseteq Z_a^{\mathcal M}.
\]
Consequently
\[
   R_a=\{z\in\omega^\omega:z\text{ is not Random-generic over }N_a\}
\]
is null, and
\[
   C_a=\{z\in\omega^\omega:z\text{ is not Cohen-generic over }N_a\}
\]
is meager.
\end{lemma}

\begin{proof}
Fix $a\in W^\ast$.  By
Lemma~\ref{lem:countable-localization-real-names}, choose a countable
admissible support $A\subseteq\omega_2$ and a $\mathbb P_A$-name $\sigma$ such
that
\[
   \sigma^{G_A}=a,
\]
where $G_A=G^\ast\cap\mathbb P_A$.  Since $\omega_2$ is regular, choose
$\beta<\omega_2$ with $A\subseteq\beta$.  Let
\[
   \dot a=(\sigma)^{\uparrow\beta}
\]
and let $\dot N$ be the $\mathbb P_\beta$-name for $L[T_2,\dot a]$.
Then the corresponding regularity tags for $\dot N$ are admissibly based.  Indeed, this is the
one-parameter instance of Definition~\ref{def:admissible-random-base}: take
\[
   A_{\mathrm{rb}}=A,
   \qquad k=1,
   \qquad \dot a_0=\sigma,
   \qquad \dot W_{\mathrm{rb}}=L[T_2,\sigma].
\]
No further presentation data are needed.  The set \(B^0_{\mathrm{rb}}\) is the
countable union of the branch-product supports of the \(W\)-parameters occurring
in the name \(\sigma\), and \(\dot B_{\mathrm{rb}}\) is the corresponding branch
footprint generated by \(B^0_{\mathrm{rb}}\) and by the explicit coding
coordinates in \(A\).  Thus the bookkeeping for admissible bases applies to
\(\dot N\), and it must eventually schedule the relative measure- and
category-amoeba stages over this base.

By the final bookkeeping, at some stage $\lambda_{\mathcal N}>\beta$ the tag
\[
   (\mathrm{Am}_{\mathcal N},\dot N^{\uparrow\lambda_{\mathcal N}})
\]
is used.  At this stage the forcing is
\[
   \mathbb A_{\mathcal N}^{N_a},
\]
computed in the interpreted base $N_a=L[T_2,a]$.  Let
$Z_a^{\mathcal N}$ be the Borel null set added by this amoeba forcing.  By the
definition of amoeba forcing for the null ideal,
\[
   B\subseteq Z_a^{\mathcal N}
\]
for every Borel null set $B$ coded in $N_a$.  Later forcing preserves the Borel
code for $Z_a^{\mathcal N}$ and the statement that it is null.  Hence the first
displayed inclusion holds in $W^\ast$.

The meager case is identical.  The bookkeeping later uses the tag
\[
   (\mathrm{Am}_{\mathcal M},\dot N^{\uparrow\lambda_{\mathcal M}})
\]
for some $\lambda_{\mathcal M}>\beta$.  The forcing
$\mathbb A_{\mathcal M}^{N_a}$ adds a Borel meager set
$Z_a^{\mathcal M}$ covering every Borel meager set coded in $N_a$, and this
covering relation remains true in the final extension.

Finally, a real is not Random-generic over $N_a$ iff it belongs to some Borel
null set coded in $N_a$.  Hence
\[
   R_a\subseteq Z_a^{\mathcal N},
\]
so $R_a$ is null.  Similarly, a real is not Cohen-generic over $N_a$ iff it
belongs to some Borel meager set coded in $N_a$.  Hence
\[
   C_a\subseteq Z_a^{\mathcal M},
\]
so $C_a$ is meager.
\end{proof}

\begin{lemma}[Hjorth--Solovay regularity in the final model]\label{lem:sigma13-regularity-final}
In $W^\ast$, every boldface $\boldsymbol{\Sigma}^1_3$ set of reals is Lebesgue measurable and has the Baire property.
\end{lemma}

\begin{proof}
We give the argument in the form in which it is used in Hjorth's proof of Corollary~2.4 of \cite{Hjorth}, with the use of Martin's Axiom replaced by Lemma~\ref{lem:random-cohen-covering-final}.  Let $a\in W^\ast$ be a real parameter and put
\[
   N_a=L[T_2,a].
\]
Here $T_2$ is the Martin--Solovay tree fixed in the preliminaries.  By Lemma~\ref{lem:random-cohen-covering-final}, almost every real in $W^\ast$ is Random over $N_a$, and comeagerly many reals in $W^\ast$ are Cohen over $N_a$.

Let $A\subseteq\omega^\omega$ be $\Sigma^1_3(a)$, say
\[
   A=\{z: W^\ast\models\varphi(z,a)\},
\]
where $\varphi(v,a)$ is a $\Sigma^1_3(a)$ formula.  We prove first that $A$ is Lebesgue measurable.  Work in $N_a$ with the Random algebra $\mathbb B_{\mathrm{rand}}^{N_a}$ and let $\dot r$ be its canonical name for the Random real.  The Boolean value
\[
   \|\varphi(\dot r,a)\|_{\mathbb B_{\mathrm{rand}}^{N_a}}
\]
is represented by a Borel set $B\subseteq\omega^\omega$, coded in $N_a$, modulo null sets.  If $z$ is Random over $N_a$, then the forcing theorem for $\mathbb B_{\mathrm{rand}}^{N_a}$ gives
\[
   z\in B
   \quad\Longleftrightarrow\quad
   N_a[z]\models\varphi(z,a).
\]
The Martin--Solovay tree $T_2$, together with its absolute complement from Corollary~\ref{cor:t2-sigma13-absoluteness}, gives the required $\Sigma^1_3$ generic absoluteness through all forcing extensions used here.  Hence, for every $z$ Random over $N_a$,
\[
   N_a[z]\models\varphi(z,a)
   \quad\Longleftrightarrow\quad
   W^\ast\models\varphi(z,a).
\]
Thus $A$ and $B$ can differ only on $R_a$, the set of reals which are not Random over $N_a$.  By Lemma~\ref{lem:random-cohen-covering-final}, $R_a$ is null.  Therefore $A\triangle B$ is null, and $A$ is Lebesgue measurable.

The proof of the Baire property is the category analogue.  Work in $N_a$ with ordinary Cohen forcing $\mathbb C^{N_a}$ and let $\dot c$ be its canonical Cohen real.  The Boolean value
\[
   \|\varphi(\dot c,a)\|_{\mathbb C^{N_a}}
\]
is represented, modulo meager sets, by a Borel set $C\subseteq\omega^\omega$ coded in $N_a$.  If $z$ is Cohen over $N_a$, then the forcing theorem gives
\[
   z\in C
   \quad\Longleftrightarrow\quad
   N_a[z]\models\varphi(z,a).
\]
The same $T_2$-absoluteness yields
\[
   N_a[z]\models\varphi(z,a)
   \quad\Longleftrightarrow\quad
   W^\ast\models\varphi(z,a)
\]
for every such Cohen real $z$.  Hence $A\triangle C\subseteq C_a$, where $C_a$ is the meager set of reals which are not Cohen over $N_a$.  Therefore $A$ differs from the Borel set $C$ by a meager set, and $A$ has the Baire property.

Since the parameter $a$ was arbitrary, the conclusion holds for all boldface $\boldsymbol{\Sigma}^1_3$ sets of reals in $W^\ast$.
\end{proof}

\subsection{\(\Sigma^1_3\)-uniformization in small \(M_1\)-generic extensions}
\label{subsec:sigma13-uniformization-m1-small}

It remains to record the lower-level \(\Sigma\)-uniformization conclusion.
Unlike the \(\Pi^1_4\)-uniformization theorem, this part does not use the
coding predicate \(\Phi\).  It follows from the relativized Steel capture
analysis of \(M_1(s)\).

What is used is the standard Steel analysis of \(M_1\) and of its
relativizations.  We only need this analysis for the small generic extensions
which occur in the construction.  Thus, throughout this subsection, a
\emph{small generic extension of \(M_1\)} means an extension \(M_1[G]\) by a
set forcing whose size in \(M_1\) is below the Woodin cardinal of \(M_1\) and
to which the canonical iteration strategy of \(M_1\) lifts.  All intermediate
models used in the construction, and in particular the final model \(W^\ast\),
are of this form.

For a real \(s\) in such an extension, let \(M_1(s)\) denote the canonical
proper class mouse over \(s\) with one Woodin cardinal.  The relevant mice
exist and are iterable by the lifted \(M_1\)-strategy.  The argument below uses
a capture theorem for this canonical mouse over \(s\); it does not identify
\(M_1(s)\) with the forcing extension generated by \(s\).

\begin{definition}
A \emph{simple \(s\)-mouse} is a sound, \(\Pi^1_2\)-iterable
\(s\)-premouse which projects to \(\omega\) and is an initial segment of the
canonical construction of \(M_1(s)\).
\end{definition}

The assertion that a real \(c\) codes a sound simple \(s\)-mouse is a
\(\Pi^1_2(s,c)\) condition.  Any two such mice compare by initial segment.
Hence the reals of \(M_1(s)\) carry the usual good \(\Sigma^1_3(s)\)
well-order \(<^1_s\): first compare the least simple \(s\)-mouse containing
the real, and then use the canonical well-order of that mouse.

We shall use the following standard form of Steel's relativized capture theorem
at the first mouse level; see the projective definability and comparison
analysis of \(M_1\) in \cite{Steel2,Steel3}.

\begin{fact}\label{fact:m1-capture-and-good-wellorder}
Let \(N=M_1[G]\) be a small generic extension of \(M_1\), and let
\(s\in\mathbb R^N\).
\begin{enumerate}
    \item If \(N\) satisfies a \(\Sigma^1_3(s)\) statement of the form
    \[
       \exists u\,\exists v\,\theta(s,u,v),
    \]
    where \(\theta\) is \(\Pi^1_2\), then there are such witnesses
    \(u,v\in M_1(s)\).

    \item For reals from \(M_1(s)\), \(\Pi^1_2\) truth is computed correctly by
    \(M_1(s)\) and is preserved to \(N\).  Equivalently, if
    \(u,v\in M_1(s)\) and \(\theta\) is \(\Pi^1_2\), then
    \[
       M_1(s)\models\theta(s,u,v)
       \quad\Longleftrightarrow\quad
       N\models\theta(s,u,v).
    \]

    \item The well-order \(<^1_s\) is good for \(\Sigma^1_3\) definitions.
    More explicitly, for every \(\Pi^1_2\) formula \(\theta(s,u,v)\), the
    relation
    \[
       \operatorname{Least}_{\theta}(s,u)
    \]
    saying that \(u\in M_1(s)\) and \(u\) is the \(<^1_s\)-least real for which
    there is a \(v\in M_1(s)\) with
    \(M_1(s)\models\theta(s,u,v)\) is uniformly \(\Sigma^1_3(s)\).
\end{enumerate}
\end{fact}

For completeness, let us spell out why this is the right theorem to apply.  The
mouse \(M_1(s)\) is the canonical mouse over the parameter \(s\).  The
comparison theorem for \(\Pi^1_2\)-iterable \(s\)-mice identifies its countable
initial segments with the segments captured by the \(\Pi^1_2\) mouse condition.
Therefore a small generic extension of \(M_1\) cannot create a new
\(\Sigma^1_3(s)\) witness without some countable initial segment of \(M_1(s)\)
already capturing the corresponding branch.  The same comparison analysis gives
the \(\Pi^1_2\) correctness used in item (2), and the usual definition of the
well-order by least simple \(s\)-mouse gives item (3).

\begin{theorem}\label{thm:small-m1-sigma13-uniformization}
Every small generic extension of \(M_1\) satisfies boldface
\(\Sigma^1_3\)-uniformization.  In particular, \(W^\ast\) satisfies
\(\Sigma^1_3\)-uniformization.
\end{theorem}

\begin{proof}
Let \(N=M_1[G]\) be a small generic extension of \(M_1\).  Work in \(N\), and
let \(A\subseteq\mathbb R^2\) be a boldface \(\Sigma^1_3\) relation.  The same
argument also works for lightface \(\Sigma^1_3\)-relations.  Fix a real
parameter \(a\) and a \(\Pi^1_2\) formula \(\psi\) such that
\[
   A(x,y)\quad\Longleftrightarrow\quad \exists z\,\psi(a,x,y,z).
\]
For each real \(x\), put \(s=a\oplus x\).  Define \(A^*(x,y)\) to hold iff
\[
   \operatorname{Least}_{\theta}(s,y),
\]
where \(\theta(s,y,z)\) is the \(\Pi^1_2\) formula obtained from
\(\psi(a,x,y,z)\) after decoding \(s=a\oplus x\).  By
Fact~\ref{fact:m1-capture-and-good-wellorder}(3), the relation \(A^*\) is
\(\Sigma^1_3(a)\), hence boldface \(\Sigma^1_3\).

We verify that \(A^*\) uniformizes \(A\).  If \(A^*(x,y)\) holds, then by
definition there is some \(z\in M_1(a\oplus x)\) such that
\[
   M_1(a\oplus x)\models \psi(a,x,y,z).
\]
By Fact~\ref{fact:m1-capture-and-good-wellorder}(2),
\(N\models\psi(a,x,y,z)\).  Hence \(A(x,y)\) holds.

The relation \(A^*\) is single-valued because \(<^1_{a\oplus x}\) is a
well-order and \(A^*(x,y)\) asserts that \(y\) is the
\(<^1_{a\oplus x}\)-least real in \(M_1(a\oplus x)\) for which a suitable
\(z\) exists.

Finally suppose that the \(x\)-section of \(A\) is nonempty in \(N\).  Then
\[
   N\models \exists y\exists z\,\psi(a,x,y,z).
\]
By Fact~\ref{fact:m1-capture-and-good-wellorder}(1), there are witnesses
\(y,z\in M_1(a\oplus x)\).  Therefore the set of \(<^1_{a\oplus x}\)-candidates
is nonempty, so it has a least element.  For this least element \(y_0\),
Fact~\ref{fact:m1-capture-and-good-wellorder}(3) gives \(A^*(x,y_0)\).

Thus \(A^*\) is a boldface \(\Sigma^1_3\) uniformization of \(A\) in \(N\).  The
final model \(W^\ast\) is a small generic extension of \(M_1\) by the forcings
fixed in the construction, so the final assertion follows.
\end{proof}


\begin{theorem}[The main theorem in the $M_1$ case]\label{thm:main-m1-final}
Assume that $M_1$ exists.  There is a forcing extension $W^\ast$ such that
\[
   W^\ast\models 2^{\aleph_0}=\aleph_2,
\]
every boldface $\boldsymbol\Sigma^1_3$ set of reals is Lebesgue measurable and has the Baire property, the $\Sigma^1_3$- and $\Pi^1_4$-uniformization properties hold, and the reals admit a $\Delta^1_4$-definable well-order.
\end{theorem}

\begin{proof}
Let $W^\ast=W[G^\ast]$ be the final model obtained from Definition~\ref{def:main-iteration}.  By Lemma~\ref{lem:basic-zero-allowable} and Fact~\ref{fact:main-iteration-infty-allowable}, the final forcing is a proper, $\aleph_2$-c.c. hybrid presentation of length $\omega_2=\omega_2^{M_1}$ with iterands of size at most $\aleph_1$.  Since $W$ satisfies $\CH$, the final forcing has size at most $\omega_2$.  Hence $W^\ast\models 2^{\aleph_0}\leq\aleph_2$.  On the other hand, ordinary Cohen forcing occurs cofinally often in the final presentation, and each such nontrivial coordinate adds a new real.  Therefore $W^\ast$ has at least $\omega_2$ many reals.  Thus $W^\ast\models 2^{\aleph_0}=\aleph_2$.

Theorem~\ref{thm:small-m1-sigma13-uniformization} gives the boldface $\Sigma^1_3$-uniformization property.  Corollary~\ref{cor:final-pi14-uniformization-local} gives the boldface $\Pi^1_4$-uniformization property.  Lemma~\ref{lem:final-delta14-wellorder} gives a $\Delta^1_4$-definable well-order of the reals.  Finally, Lemma~\ref{lem:sigma13-regularity-final} gives Lebesgue measurability and the Baire property for all boldface $\boldsymbol\Sigma^1_3$ sets of reals.  These are exactly the asserted conclusions.
\end{proof}



\section{The uniform $M_n$-version}\label{sec:uniform-mn-version}

The preceding sections were written in full detail for the first nontrivial case, namely for the construction over $M_1$.  We now record the uniform form of the argument.  The point of this section is not to introduce a new forcing construction, but to make explicit the replacements which turn the $M_1$-argument into the general $M_n$-argument and to isolate the few places where the shift in projective complexity is used.

Fix for the rest of this section a natural number $n\geq 1$, and assume that $M_n$ exists.  Let
\[
   \kappa_n=(\omega_2)^{M_n}.
\]
We work over the canonical $M_n$-ground in exactly the same way as above.  Thus we let
\[
   \vec S^n=\langle S^n_\xi\mid \xi<\omega_1^{M_n}\rangle
\]
be the $M_n$-least independent sequence of Suslin trees obtained from the canonical $M_n$-diamond sequence of Lemma~\ref{lem:mn-diamond}, and we first pass to the finite-support branch extension
\[
   W_n=M_n[\langle b_\xi\mid \xi<\omega_1^{M_n}\rangle].
\]
All subsequent forcing is a hybrid coding iteration over $W_n$.  It uses the same reservoir convention as the detailed $M_1$-construction.  Thus every successor stage first attaches a fresh $M_n$-Cohen reservoir coordinate, added by the $M_n$-computed countably closed forcing.  The second-step forcing is then a finite-support real-adding coordinate: a relative Random algebra, a relative amoeba forcing for measure or category, ordinary Cohen forcing, a Jensen--Solovay almost-disjoint coding forcing, or the trivial forcing.  Random and amoeba coordinates are understood in the relative sense: the stage is tagged by an admissible base of the form $L[T_{n+1},\vec a]$, with finitely many real parameters read on a countable admissible support, and uses the forcing computed in that base, not the corresponding forcing recomputed in the ambient universe.  At an explicit coding coordinate, the fresh reservoir generic chooses the coding area in $\vec S^n$; the construction writes the relevant real into the associated $\omega$-blocks of the Suslin sequence, reshapes in the sense of David, and finally almost-disjointly codes the reshaped set by a real.

The following dictionary will be used throughout this section:
\[
   M_1\rightsquigarrow M_n,
   \qquad
   T_2\rightsquigarrow T_{n+1},
\]
\[
   \boldsymbol\Sigma^1_3\rightsquigarrow \boldsymbol\Sigma^1_{n+2},
   \qquad
   \Pi^1_4\rightsquigarrow \Pi^1_{n+3},
   \qquad
   \Phi\rightsquigarrow \Phi_n.
\]
Here $T_{n+1}$ is the canonical weakly homogeneous tree fixed in Definition~\ref{def:canonical-tn}; by Theorem~\ref{thm:canonical-tn-existence} and Lemma~\ref{lem:tn-small-generic-absoluteness}, it represents the universal $\boldsymbol\Sigma^1_{n+2}$ set and has the required small-generic absoluteness properties.  In the $M_n$-version, Definition~\ref{def:admissible-random-base} is read with $M_1$ replaced by $M_n$ and with $T_2$ replaced by $T_{n+1}$.  Thus a relative regularity stage is generated by finitely many real names over a countable admissible support, and its base has the form
\[
   L[T_{n+1},\dot a_0,\ldots,\dot a_{k-1}].
\]

\begin{definition}[The $M_n$-localized coding predicate]\label{def:mn-localized-coding-predicate}
Let $\Phi_n(r)$ be the predicate obtained from the construction of $\Phi$ in Section~\ref{sec:M1-large-continuum-coding} by replacing $M_1$ by $M_n$, the $M_1$-least Suslin sequence by $\vec S^n$, and the $M_1$-local presentation order by the $M_n$-local presentation order.  Thus $\Phi_n(r)$ asserts that $r$ is decoded by a correct $M_n$-localized coding witness: a lower-part $M_n$-mouse code, a bounded localized presentation of the relevant hybrid forcing, a fresh $M_n$-Cohen coding area, a David-reshaped set of ordinals, and a final Jensen--Solovay almost-disjoint code.

For the two tags used in the final construction we reserve the forms
\[
   \langle\mathrm{WO},n,z_0,z_1\rangle
   \qquad\text{and}\qquad
   \langle\mathrm{UF},n,x,y,m\rangle .
\]
These tags are not used as direct coding tags except at the corresponding well-order and uniformization stages.
\end{definition}

\begin{lemma}[Complexity and exactness of $\Phi_n$]\label{lem:mn-coding-complexity-exactness}
In the final $M_n$-construction, $\Phi_n$ is a $\Sigma^1_{n+3}$ predicate.  Moreover, for the reserved tags the following exactness statement holds.  If
\[
   W_n^\ast\models \Phi_n(\langle\mathrm{WO},n,z_0,z_1\rangle),
\]
then this code was introduced at an intentional well-order coding stage comparing the canonical localized presentations of $z_0$ and $z_1$.  If
\[
   W_n^\ast\models \Phi_n(\langle\mathrm{UF},n,x,y,m\rangle),
\]
then this code was introduced at an intentional uniformization stage for the triple $(x,y,m)$.  Conversely, every intentional coding of either of these reserved tags gives the corresponding instance of $\Phi_n$ in the final model.
\end{lemma}

\begin{proof}
The complexity calculation is the uniform version of the calculation for $\Phi$ in the $M_1$ case.  A witness to $\Phi_n(r)$ is a real coding a countable localized presentation together with the associated decoding data.  The correctness of the lower-part mouse and of the localized $M_n$-initial segment is expressed using the $\Pi_n$-iterability analysis from Definition~\ref{def:pin-iterability}, Fact~\ref{fact:steel-pin-iterability}, and Lemma~\ref{lem:definable-mn-initial-segments}.  Steel's definability theorem places the relevant verification one projective level above the definition of the $M_n$-well-order, hence at level $\Sigma^1_{n+3}$.  The remaining clauses--that the displayed Suslin-tree pattern is present, that the reshaping is correct, and that the almost-disjoint code decodes the reshaped set--are arithmetic or projective of lower complexity relative to that mouse witness.  Thus the whole predicate is $\Sigma^1_{n+3}$.

The exactness assertion is the $M_n$-version of Lemma~\ref{lem:reserved-tag-coding-exactness}.  The proof of no unwanted codes uses independence of the Suslin sequence, freshness of the $M_n$-Cohen coding areas, and branch omission for selected coordinates.  The branch-footprint part is unchanged: the finitely many real names generating a regularity base are read on countable admissible supports, and hence generate only countably many coding areas; the footprint records the selected coordinates coming from those areas.  Hence the proof of Lemma~\ref{lem:no-unwanted-codes} applies after replacing $M_1$ by $M_n$.  Since the two reserved tag forms are excluded from direct coding and are used only by the well-order and uniformization clauses, any final occurrence of $\Phi_n$ on such a tag is intentional, and every intentional code is read by $\Phi_n$.
\end{proof}

Let
\[
   \langle A^n_m\mid m\in\omega\rangle
\]
be a fixed universal enumeration of lightface $\Pi^1_{n+3}$ subsets of $(\omega^\omega)^2$.  Boldface parameters are handled, as usual, by allowing the bookkeeping to range over names for real parameters and by coding the parameter into the first coordinate.  Thus it is enough to define uniformizing relations for the displayed universal family.

\begin{definition}[$M_n$-$\rho$-allowability and $M_n$-$\infty$-allowability]\label{def:mn-rho-infty-allowability}
The classes of $M_n$-$\rho$-allowable and $M_n$-$\infty$-allowable hybrid presentations are obtained from Definitions~\ref{def:zero-local-allowable}, \ref{def:limit-lambda-allowable-local}, \ref{def:relative-rho-allowable-local}, \ref{def:rho-allowable-local}, and~\ref{def:infty-allowable-local} by making the following replacements:
\[
   M_1\text{-localized presentation}
\rightsquigarrow M_n\text{-localized presentation},
\]
\[
   T_2\rightsquigarrow T_{n+1},
   \qquad
   L[T_2,\dot a_0,\ldots,\dot a_{k-1}]
   \rightsquigarrow
   L[T_{n+1},\dot a_0,\ldots,\dot a_{k-1}],
\]
\[
   \Phi\rightsquigarrow\Phi_n,
   \qquad
   \Pi^1_4\rightsquigarrow\Pi^1_{n+3},
   \qquad
   A_m\rightsquigarrow A^n_m.
\]
In particular, the regularity bases in the $M_n$-version are precisely the finite-real-parameter models
\[
   L[T_{n+1},a_0,\ldots,a_{k-1}]
\]
whose real parameters are read on countable admissible supports.
At a successor uniformization stage for a triple $(\dot x,\dot y,\dot m)$ the construction asks the same stabilized question as before: is there, locally and relative to the current $M_n$-allowable presentation, a tentative value in the $\dot x$-section of $A^n_{\dot m}$ which cannot be removed by the relevant future $M_n$-allowable tails?  If yes, the least such value, ordered by rank and by the $M_n$-canonical order of localized presentations, is protected and all competing candidates are coded with their reserved $\mathrm{UF}$-tags.  If no, a countable relative $M_n$-$\rho$-allowable local hybrid tail is chosen which forces the candidate out of the section.
\end{definition}

\begin{lemma}[Structural lemmas for $M_n$-allowability]\label{lem:mn-allowability-structural-lemmas}
The shrinking, tail, product, and persistence lemmas for allowability hold for the $M_n$-allowable hierarchy.  More explicitly:
\begin{enumerate}
   \item if a presentation is $M_n$-$\beta$-allowable and $\alpha<\beta$, then it is $M_n$-$\alpha$-allowable;
   \item tails of $M_n$-$\rho$-allowable presentations are relative $M_n$-$\rho$-allowable; moreover, any tail witnessing that a candidate can be forced out of an $A^n_m$-section can be replaced by a countable local hybrid tail with the same forcing property;
   \item products of $M_n$-$\rho$-allowable presentations are again represented by $M_n$-$\rho$-allowable presentations after the usual concatenation and renaming of the hybrid supports;
   \item tentative values remain in their $\Pi^1_{n+3}$ sections through all later $M_n$-$\infty$-allowable tails.
\end{enumerate}
\end{lemma}

\begin{proof}
The first three assertions are formal consequences of the recursive definition, exactly as in Lemmas~\ref{lem:shrinking-allowability-local}, \ref{lem:tail-allowability-local}, and~\ref{lem:rho-allowable-products-local}.  The proof uses only the fact that the class at a later derivative stage has more restrictions, that the hybrid presentation can be factored into an initial part and a relative tail, that relative regularity tags carry finitely many real parameters and branch footprints under products, and that the fresh coding areas for a product can be renamed to fresh coordinates.  The countable-tail refinement in (2) is the same local-support argument used in the $M_1$ construction: choose a name for the witness to the relevant $\Sigma^1_{n+3}$ statement and restrict to the countable complete subpresentation generated by that name, the condition, and the parameters.

For (4), suppose a value has been declared tentative at a successor stage.  By definition, every relevant future $M_n$-$\rho$-allowable tail preserves its membership in the relevant $\Pi^1_{n+3}$ section.  The only possible loss of membership would be witnessed by a $\Sigma^1_{n+3}$ statement in a later extension.  This statement is represented, relative to the real parameters already present, by the canonical tree $T_{n+1}$ and its Martin--Solovay absolute complement.  Lemma~\ref{lem:tn-small-generic-absoluteness} therefore gives the same small-generic absoluteness used in the $T_2$ argument.  Hence the proof of Lemma~\ref{lem:tentative-values-remain-local} lifts with $T_2$ replaced by $T_{n+1}$.
\end{proof}

\begin{definition}[The final $M_n$-iteration]\label{def:main-iteration-mn}
Let
\[
   F_n:\kappa_n\longrightarrow H_{\kappa_n}^{M_n}
\]
be the \(M_n\)-least bookkeeping function with the following properties.  It
places ordinary Cohen tags cofinally often.  Moreover, whenever finitely many
real names over a countable admissible support determine a regularity stage with
canonical admissibility witness
\[
   (A_{\mathrm{rb}},\dot W_{\mathrm{rb}},B^0_{\mathrm{rb}},\dot B_{\mathrm{rb}}),
\]
where
\[
   \dot W_{\mathrm{rb}}=L[T_{n+1},\dot a_0,\ldots,
      \dot a_{k-1}],
\]
the corresponding finite-real-parameter tags
\[
   (\mathrm{Rand};\dot a_0^{\uparrow},\ldots,
      \dot a_{k-1}^{\uparrow}),
   \qquad
   (\mathrm{Am}_{\mathcal N};\dot a_0^{\uparrow},\ldots,
      \dot a_{k-1}^{\uparrow}),
   \qquad
   (\mathrm{Am}_{\mathcal M};\dot a_0^{\uparrow},\ldots,
      \dot a_{k-1}^{\uparrow})
\]
appear cofinally often after that point.  The forcing is computed in the lifted
base generated by the displayed real names, and the witness records the support
on which the real parameters are read, the countable branch-product support
\(B^0_{\mathrm{rb}}\) of their \(W_n\)-parameters, and the branch footprint
\(\dot B_{\mathrm{rb}}\).  The bookkeeping also lists, cofinally often and in all
possible localized presentations, well-order tags
\[
   (\mathrm{WO},\dot z_0,
   \dot z_1)
\]
and uniformization tags
\[
   (\mathrm{UF},\dot x,
   \dot y,
   \dot m).
\]
We define by recursion a hybrid presentation
\[
   \langle (\mathcal R^n_\beta,\dot I^n_\beta)
      \mid \beta\leq\kappa_n\rangle .
\]
At every single-coordinate successor stage of this presentation we first attach the fresh $M_n$-Cohen reservoir coordinate.  At relative regularity stages the second-step forcing is the Random or amoeba forcing computed in the lifted base \(L[T_{n+1},\vec a]\) generated by the displayed finite tuple of real names, and at Cohen stages the second-step forcing is the ordinary Cohen iterand.  At well-order stages we compare the final canonical localized presentations of the two interpreted reals and, if they are distinct, use the second-step almost-disjoint coding forcing to intentionally code exactly one of
\[
   \langle\mathrm{WO},n,z_0,z_1\rangle,
   \qquad
   \langle\mathrm{WO},n,z_1,z_0\rangle,
\]
namely the tag with the smaller canonical presentation first.  At uniformization stages we apply the $M_n$-$\infty$-allowable successor rule from Definition~\ref{def:mn-rho-infty-allowability}.  Thus, if a tentative value is found, the least tentative value is protected and the non-selected candidates are coded with their tags
\[
   \langle\mathrm{UF},n,x,y,m\rangle .
\]
If no tentative value is found, we append the chosen countable relative $M_n$-$\infty$-allowable local hybrid tail, restricted below the quotient condition which forces the current candidate out of its $A^n_m$-section.  At limits we take the mixed-support hybrid limit, with countable support on reservoir coordinates and finite support on the c.c.c. real-adding coordinates.  The final forcing is denoted by
\[
   \mathbb P^n_\infty=\mathbb P^n_{\kappa_n},
\]
and if $G^n_\infty$ is generic, we write
\[
   W_n^\ast=W_n[G^n_\infty],
   \qquad
   I_n^\ast=(\dot I^n_{\kappa_n})^{G^n_\infty}.
\]
\end{definition}

\begin{fact}\label{fact:main-iteration-mn-infty-allowable}
The final presentation $(\mathcal R^n_{\kappa_n},\dot I^n_{\kappa_n})$ is $M_n$-$\infty$-allowable.  Moreover ordinary Cohen forcing occurs cofinally often.  If a finite-real-parameter regularity stage becomes meaningful, with canonical support witness
\[
   (A_{\mathrm{rb}},\dot W_{\mathrm{rb}},B^0_{\mathrm{rb}},\dot B_{\mathrm{rb}}),
\]
then the relative Random forcing, measure-amoeba forcing, and category-amoeba forcing computed in the lifted base generated by that witness occur cofinally often in the final presentation.
\end{fact}

\begin{proof}
This is the same countable-tail argument as in Fact~\ref{fact:main-iteration-infty-allowable}.  The bookkeeping function $F_n$ supplies cofinally many ordinary Cohen coordinates, cofinally many relative Random and amoeba coordinates for every finite-real-parameter regularity stage once it becomes meaningful, and cofinally many names for the well-order and uniformization tasks.  At a relative regularity stage the forcing is the one computed in the witnessed base \(L[T_{n+1},\vec a]\).  The witness carries the support on which the finitely many real parameters \(\vec a\) are read, the countable branch-product support of the \(W_n\)-parameters occurring in those names, and the branch footprint through later tails and products.  A second-case uniformization stage appends a countable relative $M_n$-$\infty$-allowable local hybrid tail rather than a single coordinate.  Since $\kappa_n$ is regular, inserting countable blocks at $\kappa_n$ many stages still gives a presentation of length at most $\kappa_n$, and the witnessing bookkeeping functions for the countable tails can be concatenated with the original bookkeeping into one global bookkeeping function.  The $M_n$-allowability clauses are preserved by this countable-tail construction.
\end{proof}

\begin{lemma}[The $M_n$ final dichotomy]\label{lem:final-dichotomy-mn}
Let $x\in\omega^\omega$, $m\in\omega$, and work in $W_n^\ast$.  Exactly one of the following alternatives holds.
\begin{enumerate}
   \item The final auxiliary object contains a tentative value for the pair $(x,m)$.  In this case there is a rank-and-presentation least such value $y_0$, and
   \[
      (x,y_0)\in A^n_m.
   \]
   Moreover, for every $y$,
   \[
      (x,y)\in A^n_m\ \text{ and }\ y\neq y_0
      \quad\Longrightarrow\quad
      \Phi_n(\langle\mathrm{UF},n,x,y,m\rangle).
   \]
   Also
   \[
      \neg\Phi_n(\langle\mathrm{UF},n,x,y_0,m\rangle).
   \]

   \item The final auxiliary object contains no tentative value for $(x,m)$.  In this case
   \[
      A^n_m(x,\cdot)=\emptyset .
   \]
\end{enumerate}
\end{lemma}

\begin{proof}
The proof is the proof of Lemma~\ref{lem:final-dichotomy-local}, with the complexity shifted by $n-1$.  We indicate the points where something has to be checked.

Suppose first that a tentative value occurs.  By Lemma~\ref{lem:mn-allowability-structural-lemmas}(4), tentative values remain in their $\Pi^1_{n+3}$ sections through later $M_n$-$\infty$-allowable tails.  Let $y_0$ be the least tentative value according to the fixed rank and $M_n$-localized-presentation order.  The bookkeeping revisits every later candidate $y$ in the same section.  Whenever $y\neq y_0$ is still a member of the $x$-section, the uniformization clause intentionally codes the tag $\langle\mathrm{UF},n,x,y,m\rangle$.  By Lemma~\ref{lem:mn-coding-complexity-exactness}, these intentional codes are exactly the final instances of $\Phi_n$ on reserved $\mathrm{UF}$-tags.  The selected value $y_0$ is never coded with its own $\mathrm{UF}$-tag, again by the exactness lemma.  This gives the first alternative.

Now suppose that no tentative value occurs.  Let $y$ be any real in the final model and choose a bounded localized name for $y$.  By the cofinality of the bookkeeping, some later uniformization stage considers the corresponding triple $(x,y,m)$.  Since no tentative value is available, the construction uses the second case and appends a countable relative $M_n$-$\infty$-allowable local hybrid tail, restricted below a quotient condition forcing
\[
   (x,y)\notin A^n_m.
\]
The complement of $A^n_m$ is $\Sigma^1_{n+3}$.  The witness to this $\Sigma^1_{n+3}$ statement is preserved through the remaining tail by the $T_{n+1}$ absoluteness supplied by Lemma~\ref{lem:tn-small-generic-absoluteness}.  Hence $(x,y)\notin A^n_m$ in $W_n^\ast$.  Since $y$ was arbitrary, the section is empty.
\end{proof}

\begin{corollary}[The $M_n$ uniformizing relations]\label{cor:final-pin3-uniformization-mn}
In $W_n^\ast$, every $\Pi^1_{n+3}$ set of pairs of reals has a $\Pi^1_{n+3}$ uniformization.  More precisely, for each $m\in\omega$, define
\[
   U^n_m(x,y)
   \quad\Longleftrightarrow\quad
   (x,y)\in A^n_m\ \land\
   \neg\Phi_n(\langle\mathrm{UF},n,x,y,m\rangle).
\]
Then $U^n_m$ uniformizes $A^n_m$.
\end{corollary}

\begin{proof}
Since $A^n_m$ is $\Pi^1_{n+3}$ and $\Phi_n$ is $\Sigma^1_{n+3}$, the displayed relation is $\Pi^1_{n+3}$.  If the $x$-section of $A^n_m$ is empty, then the $x$-section of $U^n_m$ is empty.  If the $x$-section of $A^n_m$ is nonempty, Lemma~\ref{lem:final-dichotomy-mn} gives a unique value $y_0$ in the section whose $\mathrm{UF}$-tag is not $\Phi_n$-coded.  Hence $U^n_m(x,y)$ holds exactly for $y=y_0$.  This is uniformization.  The boldface version follows by the usual coding of real parameters into the universal family and by the fact that the bookkeeping ranges over names for those parameters.
\end{proof}

\begin{lemma}[The $\Delta^1_{n+3}$ well-order in the $M_n$ extension]\label{lem:final-deltan3-wellorder-mn}
In $W_n^\ast$, the reals admit a $\Delta^1_{n+3}$-definable well-order.
\end{lemma}

\begin{proof}
For a real $z\in W_n^\ast$, let $\pi_n(z)$ be the least canonical $M_n$-localized presentation of a name whose interpretation is $z$, ordered by the fixed $M_n$-well-order of localized presentations.  Every real has such a presentation because the final hybrid forcing has countable support and every real name has bounded local support.  Distinct reals have distinct least presentations.

Define the underlying order by
\[
   z_0\triangleleft_n z_1
   \quad\Longleftrightarrow\quad
   z_0\neq z_1\ \text{ and }\ \pi_n(z_0)<_{M_n}\pi_n(z_1).
\]
This is a well-order of the reals.  The well-order stages of Definition~\ref{def:main-iteration-mn} intentionally code
\[
   \langle\mathrm{WO},n,z_0,z_1\rangle
\]
exactly when $z_0\triangleleft_n z_1$.  By Lemma~\ref{lem:mn-coding-complexity-exactness}, for distinct reals $z_0,z_1$,
\[
   z_0\triangleleft_n z_1
   \quad\Longleftrightarrow\quad
   \Phi_n(\langle\mathrm{WO},n,z_0,z_1\rangle).
\]
Thus
\[
   z_0<_{\Delta,n}z_1
   \quad\Longleftrightarrow\quad
   z_0\neq z_1\ \land\
   \Phi_n(\langle\mathrm{WO},n,z_0,z_1\rangle)
\]
is a $\Sigma^1_{n+3}$ definition of the well-order.  Since exactly one of the two opposite well-order tags is intentionally coded, the same relation is equivalently given by
\[
   z_0<_{\Delta,n}z_1
   \quad\Longleftrightarrow\quad
   z_0\neq z_1\ \land\
   \neg\Phi_n(\langle\mathrm{WO},n,z_1,z_0\rangle).
\]
This is a $\Pi^1_{n+3}$ definition.  Hence the well-order is $\Delta^1_{n+3}$.
\end{proof}

\begin{lemma}[Localized covering over \(L\lbrack T_{n+1},a\rbrack\)]\label{lem:random-cohen-covering-mn}
Let \(a\in W_n^\ast\) be a real and put
\[
   N^n_a=L[T_{n+1},a].
\]
Then there is a Borel null set \(Z^{\mathcal N}_{n,a}\in W_n^\ast\) such that
\[
   \bigcup\{B\mid B\text{ is a Borel null set coded in }N^n_a\}
   \subseteq Z^{\mathcal N}_{n,a}.
\]
There is also a Borel meager set \(Z^{\mathcal M}_{n,a}\in W_n^\ast\) such that
\[
   \bigcup\{B\mid B\text{ is a Borel meager set coded in }N^n_a\}
   \subseteq Z^{\mathcal M}_{n,a}.
\]
Consequently the set of reals which are not Random-generic over \(N^n_a\) is
null, and the set of reals which are not Cohen-generic over \(N^n_a\) is
meager.
\end{lemma}

\begin{proof}
Fix \(a\in W_n^\ast\).  By the \(M_n\)-version of
Lemma~\ref{lem:countable-localization-real-names}, choose a countable
admissible support \(A\subseteq\kappa_n\) and a \(\mathbb P_A\)-name
\(\sigma\) such that
\[
   \sigma^{G_A}=a,
\]
where \(G_A=G^n_\infty\cap\mathbb P_A\).  Choose \(\beta<\kappa_n\) with
\(A\subseteq\beta\).  Let \(\dot a=\sigma^{\uparrow\beta}\), and let
\(\dot N\) be the \(\mathbb P_\beta\)-name for \(L[T_{n+1},\dot a]\).
Then the corresponding regularity tags for \(\dot N\) are admissibly based.  This is the
higher-level version of Definition~\ref{def:admissible-random-base}, with
\[
   A_{\mathrm{rb}}=A,
   \qquad k=1,
   \qquad \dot a_0=\sigma,
   \qquad \dot W_{\mathrm{rb}}=L[T_{n+1},\sigma].
\]
The set \(B^0_{\mathrm{rb}}\) is the countable union of the branch-product
supports of the \(W_n\)-parameters occurring in \(\sigma\), and
\(\dot B_{\mathrm{rb}}\) is the corresponding branch footprint generated by
\(B^0_{\mathrm{rb}}\) and by the explicit coding coordinates in \(A\).

By the final bookkeeping, at some stage \(\lambda_{\mathcal N}>\beta\) the tag
\[
   (\mathrm{Am}_{\mathcal N},\dot N^{\uparrow\lambda_{\mathcal N}})
\]
is used.  At this stage the forcing is
\[
   \mathbb A_{\mathcal N}^{N^n_a},
\]
computed in the interpreted base \(N^n_a\).  Let \(Z^{\mathcal N}_{n,a}\) be
the Borel null set added by this amoeba forcing.  By the definition of amoeba
forcing for the null ideal,
\[
   B\subseteq Z^{\mathcal N}_{n,a}
\]
for every Borel null set \(B\) coded in \(N^n_a\).  Later forcing preserves
the Borel code for \(Z^{\mathcal N}_{n,a}\) and the statement that it is null.
Hence the first displayed inclusion holds in \(W_n^\ast\).

The meager case is the same.  The bookkeeping later uses the tag
\[
   (\mathrm{Am}_{\mathcal M},\dot N^{\uparrow\lambda_{\mathcal M}})
\]
for some \(\lambda_{\mathcal M}>\beta\).  The forcing
\(\mathbb A_{\mathcal M}^{N^n_a}\) adds a Borel meager set
\(Z^{\mathcal M}_{n,a}\) covering every Borel meager set coded in \(N^n_a\),
and this covering relation remains true in the final extension.

Finally, a real is not Random-generic over \(N^n_a\) iff it belongs to a Borel
null set coded in \(N^n_a\).  Hence the non-Random reals over \(N^n_a\) are
contained in \(Z^{\mathcal N}_{n,a}\).  Similarly, the non-Cohen reals over
\(N^n_a\) are contained in \(Z^{\mathcal M}_{n,a}\).  This proves the lemma.
\end{proof}

\begin{lemma}[Hjorth--Solovay regularity in the $M_n$ extension]\label{lem:sigma-n2-regularity-mn}
In $W_n^\ast$, every boldface $\boldsymbol\Sigma^1_{n+2}$ set of reals is Lebesgue measurable and has the Baire property.
\end{lemma}

\begin{proof}
Let $A\subseteq\omega^\omega$ be $\Sigma^1_{n+2}(a)$, where $a\in W_n^\ast$ is a real parameter, and put $N^n_a=L[T_{n+1},a]$.  By Lemma~\ref{lem:random-cohen-covering-mn}, almost every real is Random over $N^n_a$ and comeagerly many reals are Cohen over $N^n_a$.

For Lebesgue measurability, work in $N^n_a$ with the Random algebra $\mathbb B_{\mathrm{rand}}^{N^n_a}$ and let $\dot r$ be the canonical Random real.  The Boolean value of the statement ``$\dot r\in A$'' is represented, modulo null sets, by a Borel set $B$ coded in $N^n_a$.  If $z$ is Random over $N^n_a$, the forcing theorem for $\mathbb B_{\mathrm{rand}}^{N^n_a}$ gives
\[
   z\in B
   \quad\Longleftrightarrow\quad
   N^n_a[z]\models z\in A.
\]
Since $A$ is $\Sigma^1_{n+2}(a)$ and $T_{n+1}$ is the canonical weakly homogeneous tree representing the universal $\Sigma^1_{n+2}$ set, Lemma~\ref{lem:tn-small-generic-absoluteness} identifies this truth with the truth of $z\in A$ in the final model.  Therefore $A\triangle B$ is contained in the null set of reals which are not Random over $N^n_a$.  Hence $A$ is Lebesgue measurable.

For the Baire property, the same argument is carried out in $N^n_a$ with Cohen forcing.  The Boolean value of ``$\dot c\in A$'' for the canonical Cohen real is represented by a Borel set $C$ coded in $N^n_a$, and for every Cohen real $z$ over $N^n_a$ we have
\[
   z\in C
   \quad\Longleftrightarrow\quad
   z\in A
\]
by the same $T_{n+1}$ absoluteness.  Thus $A\triangle C$ is contained in the meager set of reals which are not Cohen-generic over $N^n_a$.  Hence $A$ has the Baire property.
\end{proof}

\subsection{\(\Sigma^1_{n+2}\)-uniformization in small \(M_n\)-generic extensions}
\label{subsec:sigman2-uniformization-mn-small}

We shall argue for \(\Sigma^1_{n+2}\)-uniformization now.  Again, as in the
case \(n=1\), this follows from a more general fact.  Throughout this
subsection, a small generic extension of \(M_n\) means an extension \(M_n[G]\)
by a set forcing in \(M_n\) of \(M_n\)-cardinality below the least Woodin
cardinal, in the class of forcing extensions considered here, so that the
canonical \(M_n\)-strategy and the relevant comparison arguments are preserved.

\begin{proposition}\label{prop:mn-relativized-capture}
Let \(N\) be such a small generic extension of \(M_n\), and let
\(s\in\mathbb R^N\).  Let \(M_n(s)\) be the canonical proper class
\(s\)-mouse with \(n\) Woodin cardinals.  Then the following hold in \(N\).
\begin{enumerate}
    \item If
    \[
       N\models \exists u\,\exists v\,\theta(s,u,v),
    \]
    where \(\theta\) is \(\Pi^1_{n+1}\), then there are witnesses
    \(u,v\in\mathbb R\cap M_n(s)\) such that
    \[
       M_n(s)\models \theta(s,u,v).
    \]

    \item For reals \(u,v\in M_n(s)\), \(\Pi^1_{n+1}\) truth is computed
    correctly by \(M_n(s)\) and is preserved to \(N\).  Thus, for every
    \(\Pi^1_{n+1}\) formula \(\theta\),
    \[
       M_n(s)\models\theta(s,u,v)
       \quad\Longleftrightarrow\quad
       N\models\theta(s,u,v).
    \]

    \item The canonical mouse order \(<_{n,s}\) on
    \(\mathbb R\cap M_n(s)\) is a good \(\Sigma^1_{n+2}(s)\) well-order.  More
    explicitly, if \(\theta(s,u,v)\) is \(\Pi^1_{n+1}\), then the relation
    \[
       \operatorname{Least}^{n}_{\theta}(s,u)
    \]
    saying that \(u\in M_n(s)\) and \(u\) is the \(<_{n,s}\)-least real for
    which there is a \(v\in M_n(s)\) with
    \(M_n(s)\models\theta(s,u,v)\) is uniformly \(\Sigma^1_{n+2}(s)\).
\end{enumerate}
\end{proposition}

\begin{proof}
This is the relativized Steel analysis of the canonical mice \(M_n(s)\), in
the same form as Fact~\ref{fact:m1-capture-and-good-wellorder} for \(n=1\).
We recall the ingredients in order to make clear that no new coding argument
is being used here.

The projective description of the sound initial segments of \(M_n(s)\) is the
relativized version of the \(\Pi_n\)-iterability analysis recalled in
Section~\ref{sec:canonical-inner-models}.  In the codes, the relevant
\(s\)-premice are described by a \(\Pi^1_{n+1}(s)\) condition: they are sound,
project to \(\omega\), are \(n\)-small over the parameter \(s\), and have the
required \(\Pi_n\)-iterability.  Steel comparison linearly orders these premice
by initial segment and identifies the correctly iterable ones with the initial
segments of the canonical mouse \(M_n(s)\).

The usual comparison-and-capture argument then gives the first two clauses.  If
a small generic extension satisfies a \(\Sigma^1_{n+2}(s)\) assertion, write it
in the form \(\exists u\exists v\,\theta(s,u,v)\) with \(\theta\)
\(\Pi^1_{n+1}\).  The tree or mouse witnessing this assertion is captured by a
countable initial segment of \(M_n(s)\); otherwise comparison with the canonical
construction would produce the same contradiction as in Steel's proof of the
projective correctness of \(M_n(s)\).  Conversely, once the witnesses belong to
\(M_n(s)\), \(\Pi^1_{n+1}\) correctness follows from the same comparison theorem
and the preservation convention for the generic extensions considered here.

Finally, define \(<_{n,s}\) by first taking the least sound initial segment of
\(M_n(s)\) containing the real in question, and then using the canonical
well-order of that premouse.  Since membership in the relevant initial segment
class is \(\Pi^1_{n+1}(s)\) and comparison gives initial-segment linearity,
initial segments of \(<_{n,s}\) are uniformly \(\Sigma^1_{n+2}(s)\).  Therefore
least-witness assertions for \(\Pi^1_{n+1}\) matrices are again
\(\Sigma^1_{n+2}(s)\).
\end{proof}

\begin{theorem}\label{thm:small-mn-sigman2-uniformization}
Every small generic extension of \(M_n\) satisfies boldface
\(\Sigma^1_{n+2}\)-uniformization.  In particular, \(W_n^\ast\) satisfies
\(\Sigma^1_{n+2}\)-uniformization.
\end{theorem}

\begin{proof}
Let \(N\) be a small generic extension of \(M_n\).  Work in \(N\), and let
\(A\subseteq\mathbb R^2\) be a boldface \(\Sigma^1_{n+2}\) relation.  Choose a
real parameter \(a\) and a \(\Pi^1_{n+1}\) formula \(\psi\) such that
\[
   A(x,y)
   \quad\Longleftrightarrow\quad
   \exists z\,\psi(a,x,y,z).
\]
For each real \(x\), put \(s=a\oplus x\), and let \(\theta(s,y,z)\) be the
\(\Pi^1_{n+1}\) formula obtained from \(\psi(a,x,y,z)\) after decoding
\(s\) as \(a\oplus x\).  Define \(U(x,y)\) to hold iff
\[
   \operatorname{Least}^{n}_{\theta}(a\oplus x,y).
\]
By Proposition~\ref{prop:mn-relativized-capture}(3), \(U\) is
\(\Sigma^1_{n+2}(a)\), hence boldface \(\Sigma^1_{n+2}\).

If \(U(x,y)\) holds, then for some \(z\in M_n(a\oplus x)\),
\[
   M_n(a\oplus x)\models\psi(a,x,y,z).
\]
By Proposition~\ref{prop:mn-relativized-capture}(2),
\(N\models\psi(a,x,y,z)\), and therefore \(A(x,y)\) holds.  Thus
\(U\subseteq A\).

The relation \(U\) is single-valued because \(<_{n,a\oplus x}\) is a
well-order and \(U(x,y)\) asserts that \(y\) is the least real in
\(M_n(a\oplus x)\) for which a suitable \(z\) exists.

Finally suppose that the \(x\)-section of \(A\) is nonempty in \(N\).  Then
\[
   N\models\exists y\exists z\,\psi(a,x,y,z).
\]
By Proposition~\ref{prop:mn-relativized-capture}(1), there are witnesses
\(y,z\in M_n(a\oplus x)\).  Hence the set of \(<_{n,a\oplus x}\)-candidates is
nonempty, and it has a least element \(y_0\).  By definition, \(U(x,y_0)\)
holds.  Thus \(U\) uniformizes \(A\).
\end{proof}


\begin{theorem}[The main theorem, uniform form]\label{thm:main-mn-final}
Assume that $M_n$ exists, where $1\leq n<\omega$.  There is a forcing extension $W_n^\ast$ such that
\[
   W_n^\ast\models 2^{\aleph_0}=\aleph_2,
\]
every boldface $\boldsymbol\Sigma^1_{n+2}$ set of reals is Lebesgue measurable and has the Baire property, the $\Sigma^1_{n+2}$- and $\Pi^1_{n+3}$-uniformization properties hold, and the reals admit a $\Delta^1_{n+3}$-definable well-order.
\end{theorem}

\begin{proof}
Let $W_n^\ast$ be the final extension from Definition~\ref{def:main-iteration-mn}.  The forcing has length $\kappa_n=(\omega_2)^{M_n}$, preserves $\omega_1$, has the $\aleph_2$-chain condition by the same size and $\Delta$-system argument used in Lemma~\ref{lem:basic-zero-allowable}, and every iterand has size at most $\aleph_1$ in the relevant intermediate model.  Since the preparatory model satisfies $\CH$, the final forcing has size at most $\kappa_n$, so
\[
   W_n^\ast\models 2^{\aleph_0}\leq\aleph_2.
\]
Cofinally many ordinary Cohen coordinates add new reals, and hence $W_n^\ast$ has at least $\kappa_n$ many reals.  Therefore
\[
   W_n^\ast\models 2^{\aleph_0}=\aleph_2.
\]

Theorem~\ref{thm:small-mn-sigman2-uniformization} gives $\Sigma^1_{n+2}$-uniformization.  Corollary~\ref{cor:final-pin3-uniformization-mn} gives $\Pi^1_{n+3}$-uniformization.  Lemma~\ref{lem:final-deltan3-wellorder-mn} gives a $\Delta^1_{n+3}$ well-order of the reals.  Lemma~\ref{lem:sigma-n2-regularity-mn} gives Lebesgue measurability and the Baire property for all boldface $\boldsymbol\Sigma^1_{n+2}$ sets of reals.  This proves the theorem.
\end{proof}


\bibliographystyle{plain}
\bibliography{references}




\end{document}
